\documentclass[11pt]{article}
\usepackage[margin=1in]{geometry}
\usepackage[dvipsnames]{xcolor}
\usepackage{amsmath,amsfonts,amsthm,amssymb,mathtools}
\usepackage{booktabs,enumitem,float,array}
\usepackage{tikz}
\usetikzlibrary{arrows.meta}
\usepackage[colorlinks=true,linkcolor=MidnightBlue,citecolor=ForestGreen,urlcolor=RoyalBlue]{hyperref}
\usepackage[capitalise]{cleveref}
\crefname{assumption}{Assumption}{Assumptions}
\emergencystretch=2em

\newtheorem{theorem}{Theorem}[section]
\newtheorem{corollary}[theorem]{Corollary}
\newtheorem{lemma}[theorem]{Lemma}
\newtheorem{claim}[theorem]{Claim}
\newtheorem{assumption}[theorem]{Assumption}
\newtheorem{definition}[theorem]{Definition}
\theoremstyle{remark}
\newtheorem{remark}[theorem]{Remark}

\usepackage{dsfont}
\renewcommand{\AA}{\mathbf A}
\newcommand{\BB}{\mathbf B}
\newcommand{\BBtil}{{\widetilde{\BB}}}
\newcommand{\cchi}{{\boldsymbol\chi}}
\newcommand{\dd}{\mathbf d}
\newcommand{\DD}{\mathbf D}
\renewcommand{\gg}{{\mathbf g}}
\newcommand{\GG}{\mathbf G}
\newcommand{\II}{\mathbf I}
\newcommand{\LL}{\mathbf L}
\newcommand{\LLtil}{{\widetilde{\LL}}}
\newcommand{\MM}{\mathbf M}
\newcommand{\MMtil}{{\widetilde{\MM}}}
\newcommand{\NN}{\mathbf N}
\newcommand{\PPi}{{\boldsymbol\Pi}}
\newcommand{\QQ}{\mathbf Q}
\newcommand{\SSS}{\mathbf S}
\newcommand{\uu}{\mathbf u}
\newcommand{\ww}{\mathbf w}
\newcommand{\WW}{\mathbf W}
\newcommand{\WWtil}{{\widetilde{\WW}}}
\newcommand{\wwbar}{\overline{\ww}}
\newcommand{\xx}{\mathbf x}
\newcommand{\XX}{\mathbf X}
\newcommand{\xxproxy}{\xx^{\mathrm{prox}}}
\newcommand{\xxrand}{\xx^{\mathrm{rand}}}
\newcommand{\xxres}{\xx^{\mathrm{res}}}
\newcommand{\xxvert}{\xx^{\mathrm{vert}}}
\newcommand{\yy}{\mathbf y}
\newcommand{\zz}{\mathbf z}
\newcommand{\PP}{\mathbf P}
\newcommand{\TT}{\mathbf T}
\newcommand{\UU}{\mathbf U}
\newcommand{\YY}{\mathbf{Y}}
\newcommand{\YYbar}{{\overline{\YY}}}
\newcommand{\YYtil}{{\widetilde{\YY}}}
\newcommand{\ZZ}{\mathbf Z}
\newcommand{\dheavy}{\theta_{\deg}}
\newcommand{\Ccal}{\mathcal{C}}
\newcommand{\Ecal}[1]{\mathcal E_{#1}}
\newcommand{\Ecaltil}{\widetilde{\mathcal{E}}}
\newcommand{\Ecalvert}[1]{\mathcal{E}^{\mathrm{vert}}_{#1}}
\newcommand{\Ecut}{E^{\mathrm{cut}}}
\newcommand{\Gh}{G^{\mathrm{heavy}}}
\newcommand{\Gl}{G^{\mathrm{light}}}
\newcommand{\Gret}{G^{\mathrm{ret}}}
\newcommand{\Gnew}{G^{\mathrm{new}}}
\newcommand{\R}{\mathbb R}
\newcommand{\Rcal}{\mathcal{R}}
\newcommand{\Reff}{\Rcal^{\mathrm{eff}}}
\newcommand{\Reffgap}{\mathcal R^{\mathrm{gap}}}
\newcommand{\Reffref}{\mathrlap{\widehat{\phantom{\Rcal}\mkern4mu}}\Rcal^{\mathrm{eff}}}
\newcommand{\defeq}{\coloneqq}
\newcommand{\diag}{\textsc{Diag}}
\newcommand{\Ex}{\mathbb E}
\newcommand{\norm}[1]{\lVert#1\rVert}
\newcommand{\Odouble}{\ddot{O}}
\renewcommand{\O}{\ddot{O}}
\newcommand{\oone}{\mathds{1}}
\newcommand{\tr}{\textsc{Tr}}
\newcommand{\vol}{\textsc{Vol}}
\newcommand{\cut}{\textsc{Cut}}
\newcommand{\bb}{\mathbf b}
\newcommand{\vv}{\mathbf v}
\newcommand{\Z}{\mathbb Z}
\newcommand{\dist}{\textsc{Dist}}
\newcommand{\poly}{\mathrm{poly}}
\newcommand{\gap}[1]{\item \textbf{Open.} #1}
\newcommand{\todo}[1]{\item \textbf{Fixable.} #1}

\title{$n$-Batch Dynamic Adaptive Spectral Sparsification\\
in $\Odouble(n\log^{4}n+\epsilon^{-2}n\log^{3}n)$ Time}
\date{2 November 2026}
\author{}

\begin{document}
\maketitle

\begin{abstract}
We maintain a weighted graph on $n$ vertices under blocks of at most $n$
edge insertions, deletions and weight changes, each block followed by a
request for a $(1\pm\epsilon)$ spectral sparsifier, against an adaptive
adversary.
Using a query-driven partition of the edges, a heavy/light split of the
edges and a three-step residual estimator of effective resistances -- all
proved here -- and no expander decomposition, a block costs
$\Odouble(n\log^{4}n+\epsilon^{-2}n\log^{3}n)$ time, and the output has
$O(\epsilon^{-2}n\log n)$ edges.
The bound comes from a block lemma: a block whose updates decrease no
quadratic form of the Laplacian by more than the factor $10$ costs
$\Odouble(\epsilon^{-2}\,n\log^{3}n)$, and an arbitrary block is simulated
by $O(\log n)$ such blocks that halve the deleted weights.
Removing the last logarithm reduces to sampling Gaussian vectors with covariance
$\LL(G-F)^{\dag}-\LL(G)^{\dag}$ for the deleted edge set $F$; this step is open.
Every logarithmic factor is counted; $\Odouble(\cdot)$ hides
$(\log\log n)^{O(1)}$ factors only.
\end{abstract}

\begin{center}\small
This manuscript was generated by Richard Peng and Alicia Stepin using GPT6
Astra and the Devin Harness from a manuscript on short-step interior point
methods.
\end{center}

\section{Main result}\label{sec:main}

\paragraph{The model.}
A weighted undirected graph $G=(V,E,\ww)$ on a fixed set $V$ of $n\ge2^{20}$
vertices, with $\ww_e\in[n^{-10},n^{10}]$, starts empty and changes through
\emph{blocks}.
A block is a sequence of at most $n$ updates, each an insertion
$\textsc{Insert}(e,\ww_e)$, a deletion $\textsc{Delete}(e)$ or a weight change
$\textsc{Weight}(e,\ww'_e)$, followed by one query $\textsc{Sparsify}(\epsilon)$,
$n^{-10}\le\epsilon\le\frac12$ (so $1+\epsilon^{-2}\le2\epsilon^{-2}$, and bounds
are written with $\epsilon^{-2}$ alone), which must return a graph $H$ on $V$ with
$\LL(H)\approx_{\epsilon}\LL(G)$ for the current $G$.
The adversary is adaptive: the updates of a block may depend on every graph
returned so far.
The graph is connected at every query.
Logarithms are to the base $2$, $\log n$ and $\log m$ are identified
($m\le n^{2}$), and $\Odouble(\cdot)$ denotes $O(\cdot)$ up to factors of
$(\log\log n)^{O(1)}$; no other factor is hidden anywhere.
The machine holds records linked by pointers, reals in fixed point and
integers of $O(\log n)$ bits, with unit-cost comparison, addition,
subtraction, multiplication, integer division and remainder; only the packed
distance oracle of \Cref{sec:packed} uses multiplication of packed integers,
and it runs on any word RAM with words of $O(\log n)$ bits.

Write $G_{k}$ for the graph at the $k$-th query ($G_0$ is the empty graph) and
split the edges of $G_{k}$ according to the previous query:
\begin{itemize}[leftmargin=*]
\item the \emph{retained part} $\Gret_k$ has the edges of $G_k$ that are also
edges of $G_{k-1}$, each with the weight $\min\{\ww^{(k-1)}_e,\ww^{(k)}_e\}$;
\item the \emph{new part} $\Gnew_k$ has the edges of $G_k$ that are not edges of
$G_{k-1}$, with their weights, and every edge of $\Gret_k$ whose weight
increased, with the weight $\ww^{(k)}_e-\ww^{(k-1)}_e$;
\end{itemize}
so that $\LL(G_k)=\LL(\Gret_k)+\LL(\Gnew_k)$ and $\Gnew_k$ has at most $n$
edges.
The block lemma below needs a block to be \emph{stable}:
\begin{definition}[stable block]\label{def:stable}
Block $k$ is \emph{stable} if $\LL(\Gret_k)\succeq\frac1{10}\LL(G_{k-1})$.
\end{definition}
Stability restricts only the deletions and weight decreases of a
block: insertions and weight increases are unconstrained.
It holds when the deleted and decreased parts of a block have total leverage
$\sum_{e}(\ww^{(k-1)}_e-\ww^{\mathrm{ret}}_e)\Reff(G_{k-1})_e\le\frac9{10}$, and
when every weight changes by at most the factor $10$, the regime of
interior-point methods whose weights drift by a constant factor between
refreshes.
It fails when a block deletes more than nine tenths of the weight of some cut.

\begin{lemma}[stable blocks]\label{lem:stable}
If every block is stable, there is a randomized data structure that answers
every query $\textsc{Sparsify}(\epsilon)$ of the model with a graph $H$ of at
most $10^{12}\epsilon^{-2}n\log n$ edges satisfying $\LL(H)\approx_{\epsilon}\LL(G)$,
and whose total running time over $Q$ blocks with $U$ updates is
\[
\Odouble\Bigl(U\log^{3}n+n\log^{3}n\sum_{k\le Q}\epsilon_k^{-2}\Bigr),
\]
with probability $1-Qn^{-10}$, against an adaptive adversary.
A block of $n$ updates and one query with constant $\epsilon$ costs
$\Odouble(n\log^{3}n)$.
\end{lemma}
The three parts of the bound are the maintenance of the partition structure
of \Cref{lem:partition} under the updates ($O(\log^{3}n)$ amortized per
update), the refresh of the resistance estimate at every query
($\Odouble(n\log^{3}n)$: the $\log^{3}n$ heaviest edges at every vertex,
scanned once, and $10^{3}\log n$ partition queries of $O(n\log^{2}n)$ each),
and the output itself ($O(\epsilon^{-2}n\log^{2}n)$ proposals from the
residual route, $O(\log n)$ each).

\paragraph{Arbitrary updates.}
Stability can be dispensed with at the price of one logarithm: a
block whose deletions and decreases remove the weights $\delta_e$, $e\in F$,
is simulated by $30\log n+1$ internal blocks, the $i$-th leaving the weight
$2^{-i}\delta_e$ in place, which are stable with the constant $2$, and the
last of which removes the remaining $n^{-30}\delta_e$ and is stable because
$\sum_{e\in F}\delta_e\Reff(G_k)_e\le n^{22}$.
The structure of \Cref{lem:stable} run with this schedule gives the main
theorem, proved in \Cref{sec:general}.
\begin{theorem}[main]\label{thm:main}
Without any hypothesis on the updates, there is a randomized data structure
that answers every query $\textsc{Sparsify}(\epsilon)$ with
a graph $H$ of at most $10^{12}\epsilon^{-2}n\log n$ edges satisfying
$\LL(H)\approx_{\epsilon}\LL(G)$, and its total running time over $Q$ blocks with
$U$ updates is
\[
\Odouble\Bigl(U\log^{4}n+Qn\log^{4}n+n\log^{3}n\sum_{k\le Q}\epsilon_k^{-2}\Bigr),
\]
with probability $1-Qn^{-9}$, against an adaptive adversary.
A block of $n$ updates and one query with constant $\epsilon$ costs
$\Odouble(n\log^{4}n)$.
\end{theorem}

\paragraph{Where the logarithms come from.}
Every refresh rebuilds the effective resistance estimate from scratch with
fresh randomness, so the adversary learns nothing it can exploit about the
randomness of later blocks; the only information carried from one block to
the next is a public estimate $\Reffref$ of the effective resistances of the
previous graph, which overestimates them within a constant factor
(\Cref{lem:public}).
Stability makes $10\Reffref$ an upper bound on the resistances of the
retained edges in the new graph (\Cref{lem:transfer}), the new edges are kept
exactly, and sampling with these intensities gives the sparse references of
the new refresh.
The estimate is the following: for the
$\dheavy$ heaviest edges at every vertex it is read off $10^{3}\log n$ random
projections of a sparse reference, for the remaining light edges it is a
proxy at the two endpoints plus a residual whose total energy is
$O(n/\dheavy+n/\log^{3}n)$.
The light edges with a large residual are found by $\ell_2$ queries to the
query-driven partition with the residual vectors, whose certificate costs a
factor $\log^{4}n$ over the energy; this is where $\dheavy=\log^{3}n$ is
forced, and the three terms $n\dheavy$, $n\log^{6}n/\dheavy$ and $n\log^{3}n$ of
the block cost meet.

\paragraph{Comparison with previous results.}
The previous bounds, each converted to the cost of $n$ updates followed by
one query, are the following.
\begin{itemize}[leftmargin=*]
\item Abraham, Durfee, Koutis, Krinninger and Peng~\cite{ADKKP16:conf}
maintain a $(1\pm\epsilon)$-spectral sparsifier against an oblivious
adversary with amortized update time $\poly(\log n,\epsilon^{-1})$; traced
through their theorem with $\log\rho=O(\log n)$ for the density $\rho=m/n$,
the update time is $O(\epsilon^{-2}\log^{9}n)$, so $n$ updates and one query
cost $O(\epsilon^{-2}n\log^{9}n)$ (the exponent $9$ is our reading of the
proof; the source states $\poly(\log n,\epsilon^{-1})$).
The adversary is oblivious.
\item Bernstein, van den Brand, Probst Gutenberg, Nanongkai, Saranurak,
Sidford and Sun \cite{BBGNSSS22:conf} maintain, against an adaptive
adversary, a $\log^{O(1)}n$-approximate spectral sparsifier in
$\log^{O(1)}n$ amortized update time, and an $O(k)$-approximate cut
sparsifier in $\widetilde O(n^{1/k})$ amortized update time; for $n$ updates
this gives $n\log^{O(1)}n$, with the exponent unresolved because the source
states its bounds with $\widetilde O$ and $\log^{O(1)}n$.
The query is more limited: the approximation factor is polylogarithmic, not
$1\pm\epsilon$.
\item Beimel, Kaplan, Mansour, Nissim, Saranurak and Stemmer~\cite{BKMNSS22:conf}
(Lemma~4.11 and the discussion following it) refresh a
$(1+\epsilon)$-spectral sparsifier on top of a dynamic expander
decomposition with conductance $\phi=O(1/\log^{4}m)$, maintained in
$O(\phi^{-2}\log^{7}n)$ amortized time per update (their Theorem~4.10, from
\cite{BBGNSSS22:conf}), in $\widetilde O(n+t)$ time per $t$ edge changes of
the expanders; per $n$ updates this is $n\log^{c}n$ with $c\ge15$ from the
decomposition alone and the rest unresolved ($\widetilde O$ of the source).
The query is more limited: the sparsifier is accurate only if the updates
between two refreshes are fixed when the refresh is called (their
``blinking adversary''), whereas here the adversary chooses every update
after seeing every previous output.
\item The leverage score maintenance of van den Brand, Lee, Nanongkai, Peng,
Saranurak, Sidford, Song and Wang~\cite{BLNPSSSW20:conf} (Theorem~B.1 there)
is adaptive but works on a fixed edge set whose weights drift by bounded
amounts, with a total time stated with $\widetilde O$; it does not support
insertions or deletions, so it gives no bound in this regime.
\item The batch-dynamic literature processes batches of updates for
connectivity, dynamic trees, $k$-cores and the massively parallel
model~\cite{AABD19:conf,TDB19:conf,ILMP19:conf,DDKPSS20:conf,AABDW20:conf,NO21:conf,LSYDS22:conf};
none of these maintains a spectral sparsifier, so they give no bound in
this regime.
\end{itemize}

\paragraph{Organization.}
\Cref{sec:prelim} fixes notation and states the tools:
the Laplacian solvers, the heavy/light split, the query-driven partition and
the packed distance oracle, each proved in place.
\Cref{sec:state} describes the maintained state.
\Cref{sec:refresh} is the refresh of one block: references, random vectors,
public estimate.
\Cref{sec:transfer} carries the estimate across a block and samples the next
references and the output.
\Cref{sec:proof} assembles the block procedure and proves \Cref{lem:stable}.
\Cref{sec:general} removes the stability hypothesis by the halving schedule,
proves \Cref{thm:main}, and records what stands between the result and one
logarithm less.
\Cref{sec:packed} and \Cref{sec:solver} prove the packed distance oracle and
the solver statement; the table of symbols is the last appendix.

\section{Preliminaries}\label{sec:prelim}

\paragraph{Graphs and matrices.}
For $G=(V,E,\ww)$, $\BB\in\R^{E\times V}$ is the signed incidence matrix,
$\cchi_e=\BB^{\top}\oone_e\in\R^{V}$ the signed indicator of $e$,
$\LL(G)=\BB^{\top}\diag(\ww)\BB=\sum_e\ww_e\cchi_e\cchi_e^{\top}$ the
Laplacian, $\dd(G)_v$ the weighted degree, $\DD(G)=\diag(\dd(G))$ and
$\AA(G)=\DD(G)-\LL(G)$ the adjacency matrix.
$\LL(G)^{\dag}$ is the pseudoinverse, $\PPi=\II-\frac1n\oone\oone^{\top}$ the
projection onto $\oone^{\perp}$, and
$\Reff(G)_e\defeq\cchi_e^{\top}\LL(G)^{\dag}\cchi_e$ the effective resistance,
with $\sum_e\ww_e\Reff(G)_e=n-1$ for connected $G$.
For a vector $\xx\in\R^{V}$ and a set of edges $S$,
$\Ecal2(\xx,S)\defeq\sum_{e=uv\in S}\ww_e(\xx_u-\xx_v)^{2}$ is the energy.
For positive semidefinite $\AA,\BB$ and $\epsilon\in(0,1)$,
$\AA\approx_{\epsilon}\BB$ means $(1-\epsilon)\BB\preceq\AA\preceq(1+\epsilon)\BB$,
and $a\approx_{c}b$ for scalars means $b/c\le a\le cb$.
If $\AA\succeq\frac1{c}\BB$ and $\AA\preceq\BB$ then $\ker\AA=\ker\BB$ and
$\AA^{\dag}\preceq c\BB^{\dag}$; in particular, by Rayleigh monotonicity,
adding edges to a connected graph does not increase any effective
resistance.

\paragraph{Solvers.}
\Cref{lem:chain} (\Cref{sec:solver}) gives the routine
$\textsc{LsolverKLPSS}$ of Kyng, Lee, Peng, Sachdeva and
Spielman~\cite{KLPSS16:conf}: on an SDDM matrix $\MM$ with $n$ rows, $m'$
nonzero entries and polynomially bounded entries, and an error
$\epsilon\in(0,\frac12]$, it returns in $O(m'\log n)+O(n\log^{2}n)$ time
implicit access to a symmetric matrix $\ZZ(\MM)$ with
$\ZZ(\MM)\approx_{\epsilon}\MM^{\dag}$ on the range of $\MM$, each
application of which takes $O((m'+n\log n)\log(1/\epsilon))$ time.
The operator $\ZZ(\MM)$ is a fixed matrix once built, so all its applications
are consistent; this is what makes the random vectors below Gaussian with a
fixed covariance.
The source states its failure probability as ``with high probability''
without an exponent; we write it as $n^{-14}$ throughout, and the
exponent is the one quantity of this paper that rests on an unquantified
statement of a source (\Cref{sec:solver}).

\paragraph{Sampling.}
\begin{lemma}[\cite{SS11:journal}; Theorem~1.1 of~\cite{T12:journal}]\label{lem:sparsify}
Let $G=(V,E,\ww)$ be connected, $\epsilon\in(0,\frac12]$, and let
$p_e\in(0,1]$ satisfy $p_e\ge\min\{1,\ 10^{3}\epsilon^{-2}\log n\cdot\ww_e\Reff(G)_e\}$
for every edge.
Retain every edge independently with probability $p_e$, with weight
$\ww_e/p_e$ when retained.
The retained graph $H$ satisfies $\LL(H)\approx_{\epsilon}\LL(G)$ with
probability $1-n^{-12}$, and has at most $2\sum_ep_e$ edges with probability
$1-\exp(-\frac13\sum_ep_e)$ when $\sum_ep_e\ge10^{3}\log n$.
\end{lemma}
The first statement is the matrix Chernoff bound of~\cite{T12:journal}
applied to the centered rank-one matrices
$(\ww_e/p_e)\LL(G)^{\dag/2}\cchi_e\cchi_e^{\top}\LL(G)^{\dag/2}$, as in~\cite{SS11:journal};
the second is the multiplicative Chernoff bound for independent indicators.

\paragraph{Heavy and light edges.}
\Cref{def:heavy} (\Cref{sec:resistance}) splits the edges of $G$ into the
heavy graph $\Gh$ of the $\dheavy$ heaviest edges at every vertex and the
light graph $\Gl$, and sets $\MM\defeq\LL(\Gh)+\DD(\Gl)$; \Cref{lem:topk}
collects the facts used about this split, in particular $\LL(G)\preceq2\MM$,
$\cchi_e^{\top}\MM^{\dag}\cchi_e\le2\Reff(G)_e$ and
$\tr[\MM^{\dag}\AA(\Gl)\MM^{\dag}\AA(\Gl)]\le32n/\dheavy$.

\paragraph{The packed distance oracle.}
\Cref{lem:packed} (\Cref{sec:packed}) returns, for the rows of a matrix
$\XX\in\R^{n\times10^{3}\log n}$ whose nonzero pairwise distances lie in
$[n^{-c},n^{c}]$, the squared distance of two rows within the factor $10^{3}$
in $\Odouble(1)$ worst-case time after $O(n\log^{2}n)$ randomized
preprocessing, for all pairs simultaneously with probability $1-n^{-14}$; its
answers are a fixed function of the preprocessing, so adaptive queries are
allowed.

\section{The query-driven partition}
\label{sec:partition}

This section builds the data structure \textsc{QueryDrivenPartition} whose
interface is used throughout the paper.
It maintains a weighted graph $G$ on the fixed vertex set $V$, $|V| = n$,
under edge insertions and deletions, with at most one edge between any two
vertices at any time and all weights in $[n^{-10}, n^{10}]$.
The edge set is kept partitioned into clusters, and a query with a vector
$\xx \in \R^{V}$ refines the partition until the quadratic form of $\xx$
over the edges of $G$ is certified by distances of the vertices to one
center per cluster.
The structure uses only pointers, comparisons and arithmetic on stored
reals, each at unit cost in the model of \Cref{sec:main}.
The two tools are the binary logarithmic method of Bentley and
Saxe~\cite{BS80:journal}, which keeps the number of clusters at a vertex
logarithmic, and the sweep-cut form of Cheeger's
inequality~\cite{C97:book}, which turns a large quadratic form into a
low-conductance split.

\begin{definition}
\label{def:cluster}
Let $G$ be a weighted graph on $V$ with weights in $[n^{-10}, n^{10}]$.
A \emph{cluster} $C$ is a nonempty set of edges of $G$ whose weights lie in
one interval $[\ww_C, 2\ww_C)$ with $\ww_C$ a power of two; $\ww_C$ is the
\emph{weight} of $C$.
Its vertex set $V(C)$ consists of the endpoints of its edges, the
\emph{degree} $\deg_C(v)$ is the number of edges of $C$ at $v$, and the
\emph{volume} $\vol(C) = \sum_{v \in V(C)} \deg_C(v) = 2|C|$; for
$S \subseteq V(C)$, $\vol_C(S) = \sum_{v \in S} \deg_C(v)$ and
$\cut_C(S)$ is the number of edges of $C$ with exactly one endpoint in $S$.
A \emph{partition} $\Ccal$ of $G$ is a set of pairwise edge-disjoint
clusters whose union is the edge set of $G$; its \emph{copies} are the pairs
$(C, v)$ with $C \in \Ccal$ and $v \in V(C)$.
For a vector $\xx \in \R^{V}$ and a cluster $C$ the \emph{weighted median}
$\mathrm{median}_{\xx}(C)$ is the smallest value $x_u$, $u \in V(C)$, with
$\vol_C(\{v \in V(C) : x_v \le x_u\}) \ge \vol(C)/2$; the vertices with
$x_v < \mathrm{median}_{\xx}(C)$ and those with
$x_v > \mathrm{median}_{\xx}(C)$ then each have volume at most
$\vol(C)/2$.
Each cluster carries a \emph{center} $\mathrm{center}_{\xx}(C) \in \R$ for
every vector $\xx$ queried since the last update, defined by the split
recursion: when the query of $\xx$ returns,
$\mathrm{center}_{\xx}(C) = \mathrm{median}_{\xx}(C)$ for the clusters
present at that moment, and whenever a cluster $C$ is later split into two
clusters both children receive the value $\mathrm{center}_{\xx}(C)$.
The \emph{vertex energy} of a copy with respect to $\xx$ is
\[
\Ecalvert{2}(\xx, C, v)
= \ww_C \left( x_v - \mathrm{center}_{\xx}(C) \right)^{2},
\qquad
\Ecalvert{2}(\xx, \Ccal)
= \sum_{C \in \Ccal} \sum_{v \in V(C)} \deg_C(v) \, \Ecalvert{2}(\xx, C, v),
\]
and the \emph{cluster energy} of the partition is
\[
\Ecal{2}(\xx, \Ccal)
= \sum_{C \in \Ccal} \ww_C \sum_{uv \in C} \left( x_u - x_v \right)^{2}
\le \Ecal{2}(\xx, G).
\]
\end{definition}

\begin{figure}[H]
\centering
\fbox{\begin{minipage}{0.96\textwidth}
\textbf{State of \textsc{QueryDrivenPartition}} (vertex set $V$, weights in
$[n^{-10}, n^{10}]$).
\begin{itemize}
\item The live edges of $G$ with their weights, one record per edge, with
  pointers to its two endpoints, to its cluster, to its block, and to its
  nodes in the lists and trees of the two copies at its endpoints.
\item For every integer $i$ with $-10 \log n \le i \le 10 \log n$ the
  \emph{weight class} $i$, the edges of weight in $[2^{i}, 2^{i+1})$, kept
  as at most one \emph{block} of each \emph{rank} $r \ge 0$; a block of
  rank $r$ was formed from exactly $2^{r}$ insertions and holds the live
  edges among them, partitioned into clusters.
\item For every cluster $C$: its weight $\ww_C$, its block, a doubly linked
  list of its edges, a doubly linked list of its copies $(C, v)$, and the
  list of its centers $\mathrm{center}_{\xx}(C)$ for the vectors $\xx$
  queried since the last update, each stored with the number of updates
  performed so far (a center stored before the last update is treated as
  absent).
\item For every copy $(C, v)$: the degree $\deg_C(v)$, a pointer to $v$ and
  to $C$, its node in the global doubly linked list of all copies, a doubly
  linked list of the edges of $C$ at $v$, and a balanced binary search tree
  over these edges whose nodes store subtree sizes.
\item For every vertex $v$: the list of its copies, and the values $x_v$ of
  the vectors $\xx$ queried since the last update, each stored with the
  number of updates performed so far.
\item During a query of $\xx$ only: for every cluster $C$ a balanced binary
  search tree $T_C$ over $V(C)$ keyed by $x_v$, whose nodes store the sums
  of $\deg_C(v)$, of $\deg_C(v)\, x_v$ and of $\deg_C(v)\, x_v^{2}$ over
  their subtrees; the value $\mathrm{median}_{\xx}(C)$; two heaps $H_{+}$
  and $H_{-}$ over the clusters; and the running sums
  $\Ecalvert{2}(\xx, \Ccal)$ and, for each sign $s \in \{+1, -1\}$,
  $\sum_{C} \ww_C \sum_{v : s (x_v - \mathrm{median}_{\xx}(C)) > 0}
  \deg_C(v) (x_v - \mathrm{median}_{\xx}(C))^{2}$.
\end{itemize}
\textbf{Accessors} (each in $O(1)$ per item returned): the list of all
copies; for a copy $(C, v)$ the values $\ww_C$, $\deg_C(v)$,
the pairs $(\mathrm{center}_{\xx}(C), \Ecalvert{2}(\xx, C, v))$ over the
vectors $\xx$ queried since the last update, in query order; the list of
edges of $C$ at $v$; and, given
an integer $k \in [1, \deg_C(v)]$, the $k$-th edge of $C$ at $v$ in the order
of the tree of $(C, v)$, in $O(\log n)$.
\end{minipage}}
\caption{The state of the query-driven partition.
Its guarantees are in Lemma~\ref{lem:partition}.}
\label{fig:partstate}
\end{figure}

\begin{figure}[H]
\centering
\fbox{\begin{minipage}{0.96\textwidth}
\textbf{Insert}$(e = uv, \ww_e)$.
\textbf{Input:} a vertex pair $uv$ not joined by a live edge, a weight
$\ww_e \in [n^{-10}, n^{10}]$.
\textbf{Output:} none; $G$ gains the edge.
\begin{enumerate}
\item Increase the update counter.
\item Let $i$ be the integer with $2^{i} \le \ww_e < 2^{i+1}$.
  Form a block of rank $0$ in class $i$ holding the single cluster
  $C = \{e\}$ with $\ww_C = 2^{i}$, $V(C) = \{u, v\}$, degrees $1$, and
  the trees of $(C, u)$ and $(C, v)$ holding $e$.
\item While class $i$ holds two blocks of the same rank $r$: merge them
  into one block of rank $r + 1$ whose live edges form one cluster $C$
  with $\ww_C = 2^{i}$; build the lists of $C$, the copies $(C, v)$ with
  their degrees, and the tree of each copy over the edges of $C$ at $v$;
  discard the clusters of the two old blocks.
\end{enumerate}
\medskip
\textbf{Delete}$(e)$.
\textbf{Input:} a live edge $e = uv$.
\textbf{Output:} none; $G$ loses the edge.
\begin{enumerate}
\item Increase the update counter.
\item Let $C$ be the cluster of $e$.
  Remove $e$ from the list of $C$, from the trees of $(C, u)$ and $(C, v)$,
  and from the live edges; decrease $\deg_C(u)$ and $\deg_C(v)$ by one.
\item Remove a copy whose degree became $0$ from $C$ and from the global
  list of copies; remove $C$ from its block if it has no edge left.
\end{enumerate}
\end{minipage}}
\caption{The updates of the query-driven partition.
Its guarantees are in Lemma~\ref{lem:partition}.}
\label{fig:partops}
\end{figure}

\begin{figure}[H]
\centering
\fbox{\begin{minipage}{0.96\textwidth}
\textbf{PartitionTwo}$(\xx; \Ecaltil)$.
\textbf{Input:} $\xx \in \R^{V}$, and $\Ecaltil$ with
$\Ecaltil \approx_{10} \Ecal{2}(\xx, G)$.
\textbf{Output:} the set $\Ecut$ of edges removed from $G$ by this query.
\begin{enumerate}
\item Store $x_v$ at every vertex $v$ with the current update counter.
  Sort $V$ by $x_v$; for every cluster $C$ build $T_C$ from the sorted
  order, compute $\mathrm{median}_{\xx}(C)$ by descending $T_C$ on the
  degree sums, and compute the energies of $C$ from the sums at the root of
  $T_C$ and at the median.
  For every cluster $C$ and each sign $s \in \{+1, -1\}$ insert into the
  heap $H_{s}$ the vertex of $C$ nearest to $\mathrm{median}_{\xx}(C)$ among
  those with $s(x_v - \mathrm{median}_{\xx}(C)) > 0$, if any, keyed by
  $\ww_C (x_v - \mathrm{median}_{\xx}(C))^{2}$.
  Set $\Ecut = \emptyset$.
\item While $\Ecalvert{2}(\xx, \Ccal) > 10^{8} \log^{4} n \cdot \Ecaltil$
  (with $\mathrm{center}_{\xx}(C) = \mathrm{median}_{\xx}(C)$):
\begin{enumerate}
\item Let $s \in \{+1, -1\}$ be a sign for which the vertices with
  $s(x_v - \mathrm{median}_{\xx}(C)) > 0$ carry at least half of
  $\Ecalvert{2}(\xx, \Ccal)$.
  Set $\mathrm{Vol} = 0$ and $\mathrm{Bd} = 0$.
\item Repeat: extract the maximum $(C, v)$ of $H_{s}$, mark $(C, v)$ as
  selected, add $\deg_C(v)$ to $\mathrm{Vol}$, and for every edge $vu$ of
  $C$ at $v$ add $1$ to $\mathrm{Bd}$ if $(C, u)$ is not selected and
  subtract $1$ otherwise; insert into $H_{s}$ the next vertex of $C$ on the
  side $s$ of the median in the order of $T_C$, if any.
  Stop as soon as $\mathrm{Bd} \le 10^{-3} \mathrm{Vol} / \log^{2} n$.
\item For every cluster $C$ with a selected vertex, let $S_C$ be its
  selected vertices.
  Add the edges of $C$ with exactly one endpoint in $S_C$ to $\Ecut$ and
  remove them from $G$ as in \textbf{Delete} (without changing the update
  counter).
  Move the copies $(C, v)$, $v \in S_C$, with their trees to a new cluster
  $C'$ of weight $\ww_C$ in the block of $C$, remove the vertices of $S_C$
  from $T_C$ and insert them into a new tree $T_{C'}$, give $C'$ the centers
  of $C$ for the earlier vectors, recompute the medians, the energies and
  the entries in $H_{+}$ and $H_{-}$ of $C$ and $C'$, and discard a cluster
  left without edges.
  Clear the selected marks.
\end{enumerate}
\item Set $\mathrm{center}_{\xx}(C) = \mathrm{median}_{\xx}(C)$ for every
  cluster $C$, stored with the current update counter, discard the trees
  $T_C$, and return $\Ecut$.
\end{enumerate}
\end{minipage}}
\caption{The query of the query-driven partition.
Its guarantees are in Lemma~\ref{lem:partition}.}
\label{fig:partquery}
\end{figure}

\begin{lemma}
\label{lem:partition}
There is a deterministic data structure \textsc{QueryDrivenPartition}
(\Cref{fig:partstate,fig:partops,fig:partquery}) that maintains a graph
$G$ on $V$ with at most one edge between any two vertices and weights in
$[n^{-10}, n^{10}]$, starting from the empty graph, together with a
partition $\Ccal$ of $G$ into clusters in the sense of \Cref{def:cluster},
under the operations
\begin{itemize}
\item $\textsc{Insert}(e, \ww_e)$ and $\textsc{Delete}(e)$, the
  \emph{updates};
\item $\textsc{PartitionTwo}(\xx; \Ecaltil)$ for $\xx \in \R^{V}$ and
  $\Ecaltil \approx_{10} \Ecal{2}(\xx, G)$, the \emph{queries}, each
  returning a set $\Ecut$ of edges that it removes from $G$.
\end{itemize}
Over any sequence of at most $n^{10}$ operations:
\begin{itemize}
\item there are at most $10^{3} n \log^{2} n$ copies at all times;
\item the sets $\Ecut$ returned over the sequence contain in total at most
  $10^{-1}$ times the number of updates in the sequence edges;
\item between two consecutive updates the clusters only split, and for
  every vector $\xx$ queried since the last update every cluster $C$ has a
  center $\mathrm{center}_{\xx}(C)$ as in \Cref{def:cluster} (set by the
  query of $\xx$, inherited by both halves at every later split) with
  \[
  \sum_{C \in \Ccal} \sum_{v \in V(C)}
  \deg_C(v)\, \Ecalvert{2}(\xx, C, v)
  \le 10^{8} \log^{4} n \cdot \Ecaltil ,
  \]
  where $\Ecaltil$ is the estimate passed with $\xx$;
\item the structure returns in $O(1)$ per item the list of all copies and,
  for a copy $(C, v)$, the values $\ww_C$ and $\deg_C(v)$, the list of the
  edges of $C$ at $v$, and the pairs
  $(\mathrm{center}_{\xx}(C), \Ecalvert{2}(\xx, C, v))$ over the vectors
  $\xx$ queried since the last update in query order; given an integer
  $k \in [1, \deg_C(v)]$ it returns the $k$-th edge of $C$ at $v$ in a fixed
  order in $O(\log n)$, so one uniformly random integer $k$ yields a
  uniformly random live edge of $C$ at $v$ in $O(\log n)$;
\item a sequence of $U$ updates and $Q$ queries takes
  $O(U \log^{3} n + Q n \log^{2} n)$ time in total, that is, a query costs
  amortized $O(n \log^{2} n)$ and an update amortized $O(\log^{3} n)$.
\end{itemize}
\end{lemma}

For an edge $e = uv$ in a cluster $C$ of a query vector $\xx$,
$\ww_e < 2\ww_C$ and
$(x_u - x_v)^{2}
\le 2(x_u - \mathrm{center}_{\xx}(C))^{2}
+ 2(x_v - \mathrm{center}_{\xx}(C))^{2}$
give the \emph{envelope}
\begin{equation}
\label{eq:envelope}
\ww_e (x_u - x_v)^{2}
\le 4 \left( \Ecalvert{2}(\xx, C, u) + \Ecalvert{2}(\xx, C, v) \right).
\end{equation}

The proof of \Cref{lem:partition} occupies the rest of this section.
The first tool is the sweep-cut form of Cheeger's inequality; we state and
prove the version we need, for a multigraph with unit edge weights and a
nonnegative vector, following the standard argument~\cite{C97:book}.

\begin{lemma}
\label{lem:cheeger}
Let $H$ be a multigraph with unit edge weights on a vertex set $W$, with
degrees $d(v)$, and let $\yy \in \R^{W}$ be nonnegative and not identically
zero.
There is a threshold $\theta > 0$ such that
$S = \{ v \in W : y_v > \theta \}$ is nonempty and the number of edges of
$H$ with exactly one endpoint in $S$ is at most
\[
\sqrt{ \frac{ 2 \sum_{uv \in H} (y_u - y_v)^{2} }
{ \sum_{v \in W} d(v) y_v^{2} } }
\cdot \sum_{v \in S} d(v) .
\]
\end{lemma}

\begin{proof}
Write $\rho$ for the square root in the statement, and for $t \ge 0$ let
$S_t = \{ v \in W : y_v^{2} > t \}$, with $c(t)$ edges of $H$ having
exactly one endpoint in $S_t$ and $\mathrm{vol}(S_t) = \sum_{v \in S_t} d(v)$.
An edge $uv$ with $y_u^{2} > y_v^{2}$ has exactly one endpoint in $S_t$
precisely for $t \in [y_v^{2}, y_u^{2})$, and a vertex $v$ lies in $S_t$
precisely for $t \in [0, y_v^{2})$, so
\[
\int_{0}^{\infty} c(t) \, dt
= \sum_{uv \in H} \left| y_u^{2} - y_v^{2} \right| ,
\qquad
\int_{0}^{\infty} \mathrm{vol}(S_t) \, dt
= \sum_{v \in W} d(v) y_v^{2} .
\]
By the Cauchy--Schwarz inequality and $(y_u + y_v)^{2} \le 2y_u^{2} + 2y_v^{2}$,
\[
\sum_{uv \in H} \left| y_u^{2} - y_v^{2} \right|
= \sum_{uv \in H} |y_u - y_v| (y_u + y_v)
\le \sqrt{ \sum_{uv \in H} (y_u - y_v)^{2} }
\sqrt{ 2 \sum_{v \in W} d(v) y_v^{2} }
= \rho \sum_{v \in W} d(v) y_v^{2} .
\]
Let $M = \max_v y_v^{2} > 0$.
Suppose that $c(t) > \rho \, \mathrm{vol}(S_t)$ for every $t \in [0, M)$.
Both $c$ and $\mathrm{vol}(S_t)$ are step functions of $t$ with finitely
many steps that vanish for $t \ge M$, so integrating over $[0, M)$ gives
$\int_{0}^{\infty} c(t)\, dt > \rho \int_{0}^{\infty} \mathrm{vol}(S_t)\, dt$,
contradicting the two displays.
Hence some $t \in [0, M)$ has $c(t) \le \rho \, \mathrm{vol}(S_t)$;
the set $S_t$ is nonempty because $t < M$.
If $t > 0$ take $\theta = \sqrt{t}$; if $t = 0$ take $\theta$ to be half the
smallest positive value of $\yy$.
In both cases $S = \{ v : y_v > \theta \} = S_t$, because $\yy \ge 0$.
\end{proof}

The second tool is the guarantee of one sweep, Steps 2(a)--(b) of
\textsc{PartitionTwo}.

\begin{lemma}
\label{lem:sweep}
Let $\Ccal$ be the current partition and $\xx \in \R^{V}$ a vector with
$\mathrm{center}_{\xx}(C) = \mathrm{median}_{\xx}(C)$ for all $C \in \Ccal$
and
$\Ecalvert{2}(\xx, \Ccal) > 10^{7} \log^{4} n \cdot \Ecal{2}(\xx, \Ccal)$.
Then Steps 2(a)--(b) of \textsc{PartitionTwo} terminate, and the sets $S_C$
of selected vertices satisfy
\begin{itemize}
\item $s (x_v - \mathrm{median}_{\xx}(C)) > 0$ for every $v \in S_C$, hence
  $\vol_C(S_C) \le \vol(C)/2$ for every $C$;
\item $\sum_{C \in \Ccal} \cut_C(S_C)
  \le 10^{-3} \log^{-2} n \sum_{C \in \Ccal} \vol_C(S_C)$;
\item $\sum_{C \in \Ccal} \vol_C(S_C) \ge 1$.
\end{itemize}
If $H_s$ holds, for every cluster, its copy nearest to the median on the
side $s$ when the sweep starts, the sweep takes $O(\log n)$ time per
selected copy plus $O(1)$ per edge of $C$ at a selected vertex.
\end{lemma}

\begin{proof}
Let $H$ be the multigraph with unit edge weights whose vertices are the
copies $(C, v)$ and whose edges are the edges $uv$ of every cluster $C$,
joining $(C, u)$ and $(C, v)$; the degree of $(C, v)$ in $H$ is $\deg_C(v)$,
and $H$ is the disjoint union of the clusters.
Fix the sign $s$ of Step 2(a) and define $\yy \in \R^{V(H)}$ by
$y_{(C, v)} = \sqrt{\ww_C} \, s (x_v - \mathrm{median}_{\xx}(C))$ if this
quantity is positive and $y_{(C, v)} = 0$ otherwise.
Then $\sum_{(C, v)} \deg_C(v) y_{(C, v)}^{2} \ge \Ecalvert{2}(\xx, \Ccal)/2$
by the choice of $s$, so $\yy \ne 0$, and since
$|\max(a, 0) - \max(b, 0)| \le |a - b|$,
\[
\sum_{uv \in C,\, C \in \Ccal} \left( y_{(C, u)} - y_{(C, v)} \right)^{2}
\le \sum_{C \in \Ccal} \ww_C \sum_{uv \in C} (x_u - x_v)^{2}
= \Ecal{2}(\xx, \Ccal) .
\]
\Cref{lem:cheeger} gives $\theta > 0$ such that
$S = \{ (C, v) : y_{(C, v)} > \theta \}$ is nonempty and has at most
\[
\sqrt{ \frac{ 4 \, \Ecal{2}(\xx, \Ccal) }{ \Ecalvert{2}(\xx, \Ccal) } }
\sum_{(C, v) \in S} \deg_C(v)
< \sqrt{ \frac{4}{10^{7} \log^{4} n} } \sum_{(C, v) \in S} \deg_C(v)
< \frac{10^{-3}}{\log^{2} n} \sum_{(C, v) \in S} \deg_C(v)
\]
edges of $H$ with exactly one endpoint in $S$.

The sweep extracts copies in decreasing order of
$\ww_C (x_v - \mathrm{median}_{\xx}(C))^{2} = y_{(C, v)}^{2}$ among the copies
with $s(x_v - \mathrm{median}_{\xx}(C)) > 0$: within a cluster these copies
are visited in the order of $T_C$ away from the median, along which the key
decreases, and the heap $H_s$ holds the next copy of every cluster, so the
global order is decreasing as well.
Ties are broken arbitrarily.
The quantity $\mathrm{Vol}$ is the degree sum of the selected copies, and
$\mathrm{Bd}$ is the number of edges of $H$ with exactly one selected
endpoint, as an edge is counted $+1$ when its first endpoint is selected and
$-1$ when its second one is.
Consider the moment when the last copy with $y_{(C, v)}^{2} > \theta^{2}$
has been extracted (it exists since $S \ne \emptyset$, and the heap is not
exhausted before it).
The selected set is then exactly $S$, so $\mathrm{Bd} \le 10^{-3} \log^{-2} n
\cdot \mathrm{Vol}$ and the sweep stops there if it has not stopped earlier.
At the stopping moment the selected copies form the sets $S_C$; every
selected copy satisfies $s(x_v - \mathrm{median}_{\xx}(C)) > 0$, the test
gives the second item, and at least one copy was extracted, which gives the
third item because degrees of copies are at least $1$.
The vertices $v \in V(C)$ with $x_v$ strictly on one side of
$\mathrm{median}_{\xx}(C)$ have $\vol_C$ at most $\vol(C)/2$ by
\Cref{def:cluster}, which gives the first item.
Each extraction and each heap insertion costs $O(\log n)$, locating the next
copy of a cluster in $T_C$ costs $O(\log n)$, and the update of
$\mathrm{Bd}$ visits the edges of $C$ at the selected vertex once.
\end{proof}

\begin{proof}[Proof of \Cref{lem:partition}]
We verify the invariants of the state, then the five items in turn.
Throughout, $\log n \ge 20$ and $N$, $D$, $Q$ denote the numbers of
insertions, deletions and queries of the sequence, with
$N + D + Q \le n^{10}$ and $U = N + D$.

\paragraph{Invariants.}
A block of rank $r$ is formed from exactly $2^{r}$ insertions: a rank-$0$
block from one, and a merge of two blocks of rank $r - 1$ from the
disjoint union of their insertions.
Hence $2^{r} \le N \le n^{10}$, every rank is at most $10 \log n$, each
weight class holds at most $10 \log n + 1$ blocks, and the number of weight
classes is $20 \log n + 1$.
The clusters of one block have pairwise disjoint vertex sets: this holds
after a merge, when the block has one cluster, and a split moves the copies
$(C, v)$, $v \in S_C$, to the new cluster $C'$ after removing every edge
between $S_C$ and $V(C) \setminus S_C$, so $V(C') = S_C$ and the remaining
vertex set of $C$ is $V(C) \setminus S_C$ (a vertex left without edges is
removed from both).
Every cluster has weights in $[\ww_C, 2\ww_C)$ with $\ww_C = 2^{i}$ the lower
end of its weight class, and the clusters of all blocks partition the live
edges, so $\Ccal$ is a partition in the sense of \Cref{def:cluster}.
Between two updates, \textsc{Delete} is not called, and the only changes
to $\Ccal$ are the splits of Step 2(c), each replacing a cluster $C$ by the
two clusters with vertex sets $S_C$ and $V(C) \setminus S_C$ (one of them
possibly empty) after the removal of the edges between them; so the
clusters only split.

\paragraph{Number of copies.}
By the invariants a vertex has at most one copy per block, so the number of
copies is at most
\[
n \left( 10 \log n + 1 \right) \left( 20 \log n + 1 \right)
\le n \cdot 10.05 \log n \cdot 20.05 \log n
\le 202 \, n \log^{2} n
\le 10^{3} n \log^{2} n .
\]

\paragraph{Potential.}
Define
\[
\Phi = \sum_{C \in \Ccal} \vol(C) \log_2 \vol(C) \ge 0 ,
\]
which is $0$ for the empty graph.
Since $t \mapsto t \log_2 t$ is nondecreasing on $\{0\} \cup [1, \infty)$
and volumes are even integers, $\Phi$ does not increase when edges are
removed from clusters (deletions, cut edges) or clusters are discarded.
At a merge forming a cluster with $|R|$ live edges, $\Phi$ increases by at
most $2|R| \log_2 (2|R|) \le 2|R| \cdot 2 \log n = 4|R| \log n$, because the
live graph has at most one edge per vertex pair, so $2|R| \le n(n-1) \le n^{2}$.
At a split of $C$ with $V = \vol(C)$ and $\sigma = \vol_C(S_C) \le V/2$
(\Cref{lem:sweep}), the two resulting clusters have volumes at most
$\sigma$ and $V - \sigma$, so $\Phi$ decreases by at least
\[
V \log_2 V - (V - \sigma) \log_2 (V - \sigma) - \sigma \log_2 \sigma
\ge V \log_2 V - (V - \sigma) \log_2 V - \sigma \log_2 \sigma
= \sigma \log_2 (V/\sigma)
\ge \sigma .
\]
Every sweep decreases $\Phi$ by at least $\sum_C \vol_C(S_C) \ge 1$.
The blocks of rank $r$ formed during the sequence have disjoint sets of
$2^{r}$ insertions each, so at most $N / 2^{r}$ of them are formed, each
holding at most $2^{r}$ live edges; summing over the ranks $0, \dots,
10 \log n$, the total number of live edges over all blocks formed is at
most $N (10 \log n + 1)$, and the total increase of $\Phi$ over the
sequence is at most
\[
4 \log n \cdot N (10 \log n + 1) \le 41 \, N \log^{2} n .
\]
Hence the total volume selected by all sweeps of the sequence is at most
$41 N \log^{2} n$, and by \Cref{lem:sweep} the total number of cut edges
returned is at most
\[
10^{-3} \log^{-2} n \cdot 41 N \log^{2} n \le 0.041 \, N \le 10^{-1} U .
\]

\paragraph{Termination and the certificate.}
Consider a query of $\xx$ with estimate $\Ecaltil$.
At each test of the loop the centers are the current medians, and the
cluster energy satisfies
$\Ecal{2}(\xx, \Ccal) \le \Ecal{2}(\xx, G) \le 10 \Ecaltil$, where $G$ is the
graph at the start of the query, because the sweeps only remove edges and
all weights are nonnegative; only the lower bound
$\Ecaltil \ge \Ecal{2}(\xx, G)/10$ of the hypothesis
$\Ecaltil \approx_{10} \Ecal{2}(\xx, G)$ is used.
When the loop body is entered,
$\Ecalvert{2}(\xx, \Ccal) > 10^{8} \log^{4} n \cdot \Ecaltil
\ge 10^{7} \log^{4} n \cdot \Ecal{2}(\xx, \Ccal)$, so \Cref{lem:sweep}
applies, the sweep terminates and selects volume at least $1$, and the
split decreases $\Phi$ by at least $1$; as $\Phi \ge 0$, the loop
terminates.
At exit, $\Ecalvert{2}(\xx, \Ccal) \le 10^{8} \log^{4} n \cdot \Ecaltil$
with $\mathrm{center}_{\xx}(C) = \mathrm{median}_{\xx}(C)$, which Step 3
stores; this is the certificate for $\xx$ at the end of its query.
For the maintained sums: the median of $C$ is found by descending $T_C$ on
the subtree degree sums, and for a center $c$,
$\sum_{v \in V(C)} \deg_C(v) \ww_C (x_v - c)^{2}
= \ww_C ( \sum_v \deg_C(v) x_v^{2} - 2c \sum_v \deg_C(v) x_v
+ c^{2} \sum_v \deg_C(v) )$ is read from the sums at the root of $T_C$,
and the same expression restricted to the keys strictly above (or below)
the median is a range sum in $T_C$, in $O(\log n)$.

\paragraph{Later splits preserve the certificate.}
Let $\xx$ be a vector queried since the last update, and consider a later
split of a cluster $C$, before the next update, into $C'$ with
$V(C') = S_C$ and the rest of $C$, both inheriting
$c = \mathrm{center}_{\xx}(C)$.
A vertex $v \in V(C)$ belongs to at most one of the two clusters, with
degree at most $\deg_C(v)$ there, so
\[
\sum_{v \in V(C')} \deg_{C'}(v) \ww_C (x_v - c)^{2}
+ \sum_{v \in V(C) \setminus S_C} \deg_{C}(v) \ww_C (x_v - c)^{2}
\le \sum_{v \in V(C)} \deg_C(v) \ww_C (x_v - c)^{2}
\]
where the degrees on the left are those after the split.
Thus $\sum_C \sum_{v \in V(C)} \deg_C(v) \Ecalvert{2}(\xx, C, v)$ never
increases between updates, and the certificate established at the end of
the query of $\xx$ holds until the next update.
The stored values $x_v$ and centers carry the update counter at the time
of their query, so the structure recognizes which stored vectors are since
the last update, and a split copies only those centers.

\paragraph{Accessors.}
The list of all copies is a global doubly linked list; $\ww_C$ and
$\deg_C(v)$ are fields; the edges of $C$ at $v$ form a doubly linked list;
the centers of $C$ and the values $x_v$ at $v$ for the vectors queried since
the last update are two lists in query order, and the pair
$(\mathrm{center}_{\xx}(C), \Ecalvert{2}(\xx, C, v))$ is computed from
their entries and $\ww_C$ in $O(1)$ while traversing them in parallel.
The tree of $(C, v)$ is a balanced binary search tree whose nodes store
subtree sizes, over the edges of $C$ at $v$ in an arbitrary fixed order;
it has at most $n$ nodes, so the $k$-th edge is found by descending on the
sizes in $O(\log n)$, and the edge found for a uniformly random
$k \in [1, \deg_C(v)]$ is a uniformly random live edge of $C$ at $v$.
The tree is built in linear time at a merge, updated in $O(\log n)$ per
edge insertion or removal, and moves with its copy at a split.

\paragraph{Running time.}
\textsc{Insert} forms a rank-$0$ block in $O(1)$.
A merge into a block of rank $r$ visits the at most $2^{r}$ live edges of
the two blocks (through the edge lists of their clusters), groups them by
endpoint using a scratch pointer at each vertex to its copy in the cluster
being built, and builds the lists and trees of the copies in linear time,
so it costs $O(2^{r})$; since at most $N / 2^{r}$ blocks of rank $r$ are
formed, all merges cost $O(N (10 \log n + 1)) = O(N \log n)$.
\textsc{Delete} costs $O(\log n)$ for the two tree removals.
In a query, Step 1 stores $n$ values, sorts $V$ in $O(n \log n)$, appends
every vertex in sorted order to the sequence of each of its copies in
$O(1)$ per copy, builds each $T_C$ from its sorted sequence with its sums
in time linear in $|V(C)|$, computes the median and the energies of each
cluster by one traversal of $T_C$ with running sums in time linear in
$|V(C)|$, and builds the two heaps $H_{+}$, $H_{-}$ from the first copy of
every cluster on each side of its median in time linear in the number of
clusters; the total is $O(n \log n)$ plus $O(1)$ per copy, hence
$O(n \log^{2} n)$ by the copy bound.
The heaps persist over the iterations of the loop: the entries of the
clusters untouched by a sweep remain valid, and for a cluster touched by
the sweep (one with a selected vertex) its stale entries are deleted and the
entries of the two resulting clusters inserted, in $O(\log n)$ per touched
cluster, of which there are at most as many as selected copies.
By \Cref{lem:sweep} a sweep costs $O(\log n)$ per selected copy and $O(1)$
per edge at a selected vertex, that is, $O(\log n)$ per unit of selected
volume.
The split of Step 2(c) removes each cut edge from two trees and two lists
in $O(\log n)$, moves each selected copy in $O(1)$, moves the edges inside
$S_C$ to the list of $C'$ by traversing the lists of the selected copies in
$O(1)$ per edge, removes the $|S_C| \le \vol_C(S_C)$ selected vertices from
$T_C$ and inserts them into $T_{C'}$ in $O(\log n)$ each, and recomputes
the medians, energies and heap entries of $C$ and $C'$ in $O(\log n)$;
the number of cut edges is at most the selected volume, so the split also
costs $O(\log n)$ per unit of selected volume.
Summed over the sequence, the sweeps and splits cost
$O(\log n \cdot 41 N \log^{2} n) = O(N \log^{3} n)$ by the potential
argument.
It remains to account for copying the centers of earlier vectors at a
split, $O(1)$ per center.
Let $\Psi$ be the number of copies minus twice the number of clusters, which
is nonnegative since a cluster has at least one edge, hence at least two
copies.
A split that produces a nonempty $C'$ increases the number of clusters by
one and does not increase the number of copies, so it decreases $\Psi$ by
at least $2$; removing an edge decreases the number of copies by the number
of endpoints left without edges and discards a cluster only when both its
endpoints are left without edges in it, so $\Psi$ does not increase.
Between two consecutive updates $\Psi$ starts at most
$10^{3} n \log^{2} n$, so at most $500 n \log^{2} n$ splits with nonempty
$C'$ occur, and each copies at most as many centers as there are queries
since the last update; if $q$ queries fall between two consecutive updates,
their center copying costs $O(q \cdot n \log^{2} n)$ in total.
Adding up, the sequence costs
$O(N \log n + D \log n + N \log^{3} n + Q n \log^{2} n)
= O(U \log^{3} n + Q n \log^{2} n)$.
\end{proof}

\section{Heavy and light edges, and the residual estimator}
\label{sec:resistance}

Throughout this section $G=(V,E,\ww)$ is a connected weighted graph on
$n\ge2^{20}$ vertices, so that $\log n\ge20$, with weights in
$[n^{-10},n^{10}]$; $\PPi\defeq\II-\frac1n\oone\oone^{\top}$ is the
orthogonal projection onto $\oone^{\perp}$, $\LL(G)^{\dag}$ is the
pseudoinverse of the Laplacian, $\LL(G)^{\dag/2}\defeq(\LL(G)^{\dag})^{1/2}$,
and $\Reff(G)_e=\cchi_e^{\top}\LL(G)^{\dag}\cchi_e$ is the effective
resistance of the edge $e$.
For a symmetric positive semidefinite matrix $\NN$, $\NN^{\dag}$ is its
pseudoinverse, $\norm{\NN}$ its operator norm, $\norm{\NN}_F$ its Frobenius
norm and $\norm{\NN}_{S_6}\defeq\tr[(\NN^{\top}\NN)^{3}]^{1/6}$ its
Schatten $6$-norm, defined for every square matrix.
The section splits the edges of $G$ into a heavy part, which is read
directly, and a light part, whose contribution to the effective
resistances is estimated by a residual; the main statement is
\Cref{lem:refresh}.

\subsection{Heavy and light edges}

\begin{definition}[heavy and light edges]
\label{def:heavy}
For $\dheavy\ge200$, the \emph{heavy graph} $\Gh$ consists of the
$\dheavy$ heaviest edges at every vertex (all edges at a vertex of degree at
most $\dheavy$; ties broken by a fixed order), the \emph{light graph} is
$\Gl\defeq E\setminus\Gh$, and
\[
\MM\ \defeq\ \LL(\Gh)+\DD(\Gl).
\]
\end{definition}
An edge is heavy if at least one of its endpoints counts it among its
$\dheavy$ heaviest edges; a light edge is light at both endpoints.

\begin{lemma}[the heavy and light parts]
\label{lem:topk}
Let $G=(V,E,\ww)$ be connected with the split of \Cref{def:heavy}.
\begin{enumerate}[label=(\alph*),leftmargin=*]
\item $|E(\Gh)|\le n\dheavy$.
\item $\ww_e\le\min\{\dd(\Gh)_u,\dd(\Gh)_v\}/\dheavy$ for every light edge
$e=uv$.
\item $-\DD(\Gl)\preceq\AA(\Gl)\preceq\DD(\Gl)\preceq\MM$ and
$\LL(G)\preceq2\MM$.
\item If $\Gl=\emptyset$ then $\MM=\LL(G)$ and the kernel of $\MM$ is spanned
by $\oone$; if $\Gl\neq\emptyset$ then $\MM$ is positive definite.
In both cases the range of $\MM$ contains $\oone^{\perp}$.
\item $\xx^{\top}\MM^{\dag}\xx\le2\,\xx^{\top}\LL(G)^{\dag}\xx$ for every
$\xx\perp\oone$; in particular
$\cchi_e^{\top}\MM^{\dag}\cchi_e\le2\Reff(G)_e$ for every edge $e$.
\item $\tr[\DD(\Gl)\MM^{\dag}]\le n$.
\item $\tr\bigl[\MM^{\dag}\AA(\Gl)\MM^{\dag}\AA(\Gl)\bigr]
=\norm{\MM^{\dag/2}\AA(\Gl)\MM^{\dag/2}}_F^{2}\le32n/\dheavy$.
\end{enumerate}
\end{lemma}
\begin{proof}
(a) Each vertex contributes at most $\dheavy$ heavy edges.

(b) A light edge $e=uv$ is not among the $\dheavy$ heaviest edges at $u$,
so its weight is at most the weight of each of those $\dheavy$ edges, and
$\dheavy\ww_e$ is at most their total weight, which is at most $\dd(\Gh)_u$;
likewise at $v$.

(c) The Laplacian $\LL(\Gl)=\DD(\Gl)-\AA(\Gl)$ and the signless Laplacian
$\DD(\Gl)+\AA(\Gl)$ are positive semidefinite (the quadratic forms are
$\sum_{uv\in\Gl}\ww_{uv}(\xx_u-\xx_v)^{2}$ and
$\sum_{uv\in\Gl}\ww_{uv}(\xx_u+\xx_v)^{2}$), so
\begin{align*}
-\DD(\Gl)&\preceq\AA(\Gl)\preceq\DD(\Gl)\preceq\LL(\Gh)+\DD(\Gl)=\MM,\\
\LL(G)&=\LL(\Gh)+\LL(\Gl)=\MM-\AA(\Gl)\preceq\MM+\DD(\Gl)\preceq2\MM .
\end{align*}

(d) If $\Gl=\emptyset$ then $\Gh=G$ and $\MM=\LL(G)$, whose kernel is
spanned by $\oone$ since $G$ is connected.
Let $\Gl\neq\emptyset$ and $\MM\uu=0$.
Then $\uu^{\top}\LL(\Gh)\uu=0$ and $\uu^{\top}\DD(\Gl)\uu=0$, so $\uu$ is
constant on every connected component of $\Gh$ and vanishes at every vertex
of positive light degree.
A component of $\Gh$ containing an endpoint of a light edge therefore
carries the value $0$.
A component of $\Gh$ containing no endpoint of a light edge has all its
edges in $\Gh$ and is a component of $G$, hence all of $V$, which forces
$\Gl=\emptyset$.
So $\uu=0$, and $\MM$ is positive definite.

(e) Let $\xx\perp\oone$.
By (d), $\xx$ lies in the range of $\MM$, so
\begin{align*}
\xx^{\top}\MM^{\dag}\xx
&=\max_{\uu\in\R^{V}}\bigl(2\xx^{\top}\uu-\uu^{\top}\MM\uu\bigr)\\
&\le\max_{\uu\in\R^{V}}\bigl(2\xx^{\top}\uu-\tfrac12\uu^{\top}\LL(G)\uu\bigr)
=2\,\xx^{\top}\LL(G)^{\dag}\xx ,
\end{align*}
where the inequality is $\LL(G)\preceq2\MM$ from (c), and the last maximum
is attained at $\uu=2\LL(G)^{\dag}\xx$: the objective is invariant under
adding multiples of $\oone$ to $\uu$ (as $\xx^{\top}\oone=0$) and strictly
concave on $\oone^{\perp}$.
Taking $\xx=\cchi_e$ gives the second claim.

(f) If $\Gl=\emptyset$ then $\DD(\Gl)=0$.
Otherwise $\MM$ is positive definite, and for a vertex $v$ with
$\dd(\Gl)_v>0$, $\MM\succeq\DD(\Gl)\succeq\dd(\Gl)_v\oone_v\oone_v^{\top}$ gives
\[
(\MM^{-1})_{v,v}
=\max_{\uu\in\R^{V}}\bigl(2\uu_v-\uu^{\top}\MM\uu\bigr)
\le\max_{\uu\in\R^{V}}\bigl(2\uu_v-\dd(\Gl)_v\uu_v^{2}\bigr)
=\frac1{\dd(\Gl)_v},
\]
so $\tr[\DD(\Gl)\MM^{-1}]=\sum_{v:\dd(\Gl)_v>0}\dd(\Gl)_v(\MM^{-1})_{v,v}\le n$.

(g) The identity
$\tr[\MM^{\dag}\AA(\Gl)\MM^{\dag}\AA(\Gl)]=\norm{\MM^{\dag/2}\AA(\Gl)\MM^{\dag/2}}_F^{2}$
is the cyclicity of the trace.
If $\Gl=\emptyset$ then $\AA(\Gl)=0$ and there is nothing to prove, so let
$\Gl\neq\emptyset$; then $\MM$ is positive definite by (d).

\emph{A subgraph of the heavy graph.}
Write $v\to u$ if $vu$ is among the $\dheavy$ heaviest edges at $v$ (``$v$
chose $vu$''), and $\max\{\ww_f:f\in\Gl,\ v\in f\}$ for the heaviest light
weight at $v$ (zero if $v$ has no light edge).
Let $H^{\mathrm{heavy}}$ be the graph on $V$ whose edges are the edges
$uv\in\Gh$ with positive weight
\begin{align*}
\ww(H^{\mathrm{heavy}})_{uv}\ \defeq\ \tfrac12\bigl(
&\max\{\ww_f:f\in\Gl,\ u\in f\}\cdot[u\to v]\\
&+\max\{\ww_f:f\in\Gl,\ v\in f\}\cdot[v\to u]\bigr),
\end{align*}
where $[\cdot]$ is the indicator.
Then $\ww(H^{\mathrm{heavy}})_{uv}\le\ww_{uv}$: an edge chosen by $u$ weighs
at least every light edge at $u$, so each of the two terms is at most
$\ww_{uv}$.
Hence $\LL(H^{\mathrm{heavy}})\preceq\LL(\Gh)$ and
\[
\LL(H^{\mathrm{heavy}})+\DD(\Gl)\ \preceq\ \MM .
\]
If $v$ has a light edge then $v$ has more than $\dheavy$ edges, so $v$ chose
exactly $\dheavy$ edges, each of weight at least
$\frac12\max\{\ww_f:f\in\Gl,\ v\in f\}$ in $H^{\mathrm{heavy}}$; therefore
\begin{equation}\label{eq:halfdeg}
\dd(H^{\mathrm{heavy}})_v\ \ge\ \tfrac12\dheavy\max\{\ww_f:f\in\Gl,\ v\in f\}
\qquad\text{for every $v$ with a light edge.}
\end{equation}

\emph{The graph $H$.}
Let $H$ be the graph with edge set $E(H^{\mathrm{heavy}})\cup E(\Gl)$ (the two
sets are disjoint), weights $\ww(H^{\mathrm{heavy}})$ on the first part and
$\ww$ on the second, and vertex set $V_H$ the vertices of positive degree in
it.
A vertex outside $V_H$ has no light edge and no edge of $H^{\mathrm{heavy}}$,
so the corresponding rows and columns of $\AA(\Gl)$, $\DD(\Gl)$ and
$\LL(H^{\mathrm{heavy}})$ vanish; all matrices below are restricted to
$V_H\times V_H$, which changes none of the traces in question.
On $V_H$,
\[
\LL(H^{\mathrm{heavy}})+\DD(\Gl)=\DD(H)-\AA(H^{\mathrm{heavy}}),
\]
and this matrix is positive definite: a vector in its kernel is constant
on every component of $H^{\mathrm{heavy}}[V_H]$ and zero at every vertex
with a light edge, and every component of $H^{\mathrm{heavy}}[V_H]$ contains
a vertex with a light edge, because every edge of $H^{\mathrm{heavy}}$ has
positive weight only if a choosing endpoint has a light edge.

\emph{Normalized Frobenius norms.}
Orient every edge of $H$: an edge $uv$ of $H^{\mathrm{heavy}}$ away from a
choosing endpoint $u$ whose heaviest light weight is largest among the
choosing endpoints, and a light edge $uv$ away from either endpoint.
In both cases the weight $\ww(H)_{uv}$ of the oriented edge $u\to v$ is at
most $\max\{\ww_f:f\in\Gl,\ u\in f\}$: for a light edge this is the
definition of the heaviest light weight, and for an edge of
$H^{\mathrm{heavy}}$ the weight is the half-sum of at most two terms each at
most $\max\{\ww_f:f\in\Gl,\ u\in f\}$.
Moreover $u$ has a light edge (the weight is positive), so
\Cref{eq:halfdeg} gives
\[
\frac{\ww(H)_{uv}}{\dd(H)_u}
\le\frac{\max\{\ww_f:f\in\Gl,\ u\in f\}}{\dd(H^{\mathrm{heavy}})_u}
\le\frac2{\dheavy}
\qquad\text{for every oriented edge $u\to v$ of $H$.}
\]
Hence
\[
\norm{\DD(H)^{-1/2}\AA(H)\DD(H)^{-1/2}}_F^{2}
=2\sum_{uv\in E(H)}\frac{\ww(H)_{uv}^{2}}{\dd(H)_u\dd(H)_v}
\le\frac4{\dheavy}\sum_{u\to v}\frac{\ww(H)_{uv}}{\dd(H)_v}
\le\frac{4n}{\dheavy},
\]
because the total weight of the edges oriented into $v$ is at most
$\dd(H)_v$ and $|V_H|\le n$.
The matrices $\AA(H^{\mathrm{heavy}})$ and $\AA(\Gl)$ have disjoint supports
and sum to $\AA(H)$, so
\begin{equation}\label{eq:twonorms}
\norm{\DD(H)^{-1/2}\AA(H^{\mathrm{heavy}})\DD(H)^{-1/2}}_F^{2}
+\norm{\DD(H)^{-1/2}\AA(\Gl)\DD(H)^{-1/2}}_F^{2}
\le\frac{4n}{\dheavy}.
\end{equation}

\emph{An orthogonal similarity.}
The matrix
$(\LL(H^{\mathrm{heavy}})+\DD(\Gl))^{-1/2}\AA(\Gl)(\LL(H^{\mathrm{heavy}})+\DD(\Gl))^{-1/2}$
lies between $-\II$ and $\II$, because
$\LL(H^{\mathrm{heavy}})+\DD(\Gl)\pm\AA(\Gl)=\LL(H^{\mathrm{heavy}})+(\DD(\Gl)\pm\AA(\Gl))\succeq0$.
The matrix
\[
\UU\defeq\bigl(\DD(H)^{-1/2}(\LL(H^{\mathrm{heavy}})+\DD(\Gl))\DD(H)^{-1/2}\bigr)^{-1/2}
\DD(H)^{-1/2}(\LL(H^{\mathrm{heavy}})+\DD(\Gl))^{1/2}
\]
satisfies $\UU\UU^{\top}=\II$, so it is orthogonal, and
\begin{equation}\label{eq:tracesim}
\begin{aligned}
&\UU\,(\LL(H^{\mathrm{heavy}})+\DD(\Gl))^{-1/2}\AA(\Gl)(\LL(H^{\mathrm{heavy}})+\DD(\Gl))^{-1/2}\,\UU^{\top}\\
&\quad=\bigl(\DD(H)^{-1/2}(\LL(H^{\mathrm{heavy}})+\DD(\Gl))\DD(H)^{-1/2}\bigr)^{-1/2}
\,\DD(H)^{-1/2}\AA(\Gl)\DD(H)^{-1/2}\,\\
&\qquad\cdot\bigl(\DD(H)^{-1/2}(\LL(H^{\mathrm{heavy}})+\DD(\Gl))\DD(H)^{-1/2}\bigr)^{-1/2}.
\end{aligned}
\end{equation}
Both sides have the same Frobenius norm and lie between $-\II$ and $\II$;
we write $\TT$ for the right-hand side.

\emph{Spectral split.}
Let $\PP$ be the spectral projection of
\[
\DD(H)^{-1/2}(\LL(H^{\mathrm{heavy}})+\DD(\Gl))\DD(H)^{-1/2}
=\II-\DD(H)^{-1/2}\AA(H^{\mathrm{heavy}})\DD(H)^{-1/2}
\]
onto its eigenvalues below $\frac12$.
On the range of $\PP$ the eigenvalues of the symmetric matrix
$\DD(H)^{-1/2}\AA(H^{\mathrm{heavy}})\DD(H)^{-1/2}$ exceed $\frac12$, and a
symmetric matrix has at most four times its squared Frobenius norm many
eigenvalues above $\frac12$, so
\[
\tr[\PP]\le4\,\norm{\DD(H)^{-1/2}\AA(H^{\mathrm{heavy}})\DD(H)^{-1/2}}_F^{2}.
\]
On the range of $\II-\PP$ the inverse of
$\DD(H)^{-1/2}(\LL(H^{\mathrm{heavy}})+\DD(\Gl))\DD(H)^{-1/2}$ has eigenvalues
at most $2$, and both $\PP$ and $\II-\PP$ commute with its inverse square
root.
Write the right-hand side of \Cref{eq:tracesim} as the sum of its four
compressions by $\PP$ and $\II-\PP$ on the two sides; these four matrices are
orthogonal in the Frobenius inner product, so the squared Frobenius norm of
the sum is the sum of their squared Frobenius norms.
\begin{itemize}[leftmargin=*]
\item The compression $(\II-\PP)\TT(\II-\PP)$ is the product of the matrix
\[
(\II-\PP)\bigl(\DD(H)^{-1/2}(\LL(H^{\mathrm{heavy}})+\DD(\Gl))\DD(H)^{-1/2}\bigr)^{-1/2},
\]
of operator norm at most $\sqrt2$, then $\DD(H)^{-1/2}\AA(\Gl)\DD(H)^{-1/2}$,
then the transpose of the first factor; so
\[
\norm{(\II-\PP)\TT(\II-\PP)}_F^{2}\le4\,\norm{\DD(H)^{-1/2}\AA(\Gl)\DD(H)^{-1/2}}_F^{2}.
\]
\item The three compressions touching $\PP$ have total squared Frobenius
norm at most $2\tr[\PP]$: $\norm{\PP\TT}_F^{2}=\tr[\PP\TT^{2}\PP]\le\tr[\PP]$
since $\TT^{2}\preceq\II$, this covers the compressions $\PP\TT\PP$ and
$\PP\TT(\II-\PP)$, and
$\norm{(\II-\PP)\TT\PP}_F=\norm{\PP\TT(\II-\PP)}_F\le\norm{\PP\TT}_F$.
\end{itemize}
Therefore, by \Cref{eq:twonorms},
\begin{align*}
\norm{\TT}_F^{2}
&\le4\,\norm{\DD(H)^{-1/2}\AA(\Gl)\DD(H)^{-1/2}}_F^{2}\\
&\quad+8\,\norm{\DD(H)^{-1/2}\AA(H^{\mathrm{heavy}})\DD(H)^{-1/2}}_F^{2}\\
&\le\frac{32n}{\dheavy}.
\end{align*}

\emph{Monotonicity.}
The map $\NN\mapsto\tr[\NN\AA(\Gl)\NN\AA(\Gl)]$ is nondecreasing on the
positive semidefinite cone: its derivative at $\NN\succeq0$ along
$\boldsymbol\Delta\succeq0$ is
$2\tr[\boldsymbol\Delta\,\AA(\Gl)\NN\AA(\Gl)]\ge0$, as
$\AA(\Gl)\NN\AA(\Gl)\succeq0$.
For every $t>0$,
$\LL(H^{\mathrm{heavy}})+\DD(\Gl)+t\II\preceq\MM+t\II$ gives
$(\MM+t\II)^{-1}\preceq(\LL(H^{\mathrm{heavy}})+\DD(\Gl)+t\II)^{-1}$, on
$V_H$ as on $V$; the latter matrix is block diagonal with respect to
$V_H$ and $V\setminus V_H$ and $\AA(\Gl)$ is supported on $V_H\times V_H$,
so the trace only involves the $V_H$ block, which converges to the inverse
of the positive definite matrix $\LL(H^{\mathrm{heavy}})+\DD(\Gl)$ on $V_H$
as $t\downarrow0$.
Hence
\begin{align*}
\tr[\MM^{-1}\AA(\Gl)\MM^{-1}\AA(\Gl)]
&\le\tr\bigl[(\LL(H^{\mathrm{heavy}})+\DD(\Gl))^{-1}\AA(\Gl)\\
&\qquad\cdot(\LL(H^{\mathrm{heavy}})+\DD(\Gl))^{-1}\AA(\Gl)\bigr]\\
&=\norm{\TT}_F^{2}\\
&\le\frac{32n}{\dheavy}.
\qedhere
\end{align*}
\end{proof}

\subsection{Two tools}

The acceptance tests and the all-pairs estimate of \Cref{lem:refresh} are
certified with the following tail bound for Gaussian quadratic forms.
The upper tail is Theorem~1 of Hsu, Kakade and Zhang~\cite{HKZ12:journal},
specialized to a standard Gaussian vector; the lower tail is Lemma~1 of
Laurent and Massart~\cite{LM00:journal}, applied to the eigenvalues of
$\NN$.
\begin{lemma}[Theorem~1 of~\cite{HKZ12:journal}; Lemma~1 of~\cite{LM00:journal}]
\label{lem:quadtail}
Let $\NN\in\R^{d\times d}$ be symmetric positive semidefinite and let
$\gg\in\R^{d}$ have independent standard Gaussian entries.
Then $\tr[\NN^{2}]\le\norm{\NN}\tr[\NN]$, and for every $t>0$
\begin{align*}
\Pr\bigl[\gg^{\top}\NN\gg>\tr[\NN]+2\sqrt{\tr[\NN^{2}]\,t}+2\norm{\NN}t\bigr]
&\le e^{-t},\\
\Pr\bigl[\gg^{\top}\NN\gg<\tr[\NN]-2\sqrt{\tr[\NN^{2}]\,t}\bigr]
&\le e^{-t}.
\end{align*}
\end{lemma}
\begin{proof}
Let $\lambda_1,\dots,\lambda_d\ge0$ be the eigenvalues of $\NN$, so that
$\tr[\NN^{2}]=\sum_i\lambda_i^{2}\le\max_i\lambda_i\sum_i\lambda_i=\norm{\NN}\tr[\NN]$.
Theorem~1 of~\cite{HKZ12:journal} states, for a matrix $\AA'$, the vector
$\xx$ with $\Ex[\exp(\boldsymbol\alpha^{\top}\xx)]\le\exp(\sigma^{2}\norm{\boldsymbol\alpha}^{2}/2)$
for all $\boldsymbol\alpha$ and $\boldsymbol\Sigma\defeq\AA'^{\top}\AA'$, that
$\Pr[\norm{\AA'\xx}^{2}>\sigma^{2}(\tr[\boldsymbol\Sigma]+2\sqrt{\tr[\boldsymbol\Sigma^{2}]t}+2\norm{\boldsymbol\Sigma}t)]\le e^{-t}$;
with $\xx=\gg$, $\sigma=1$ and $\AA'=\NN^{1/2}$ this is the first bound.
For the second, rotating $\gg$ to the eigenbasis of $\NN$ (which preserves
its law) writes $\gg^{\top}\NN\gg=\sum_i\lambda_i\gg_i^{2}$, and Lemma~1
of~\cite{LM00:journal} states, for independent standard Gaussians $\gg_i$
and nonnegative $\lambda_i$, that
$\Pr[\sum_i\lambda_i(\gg_i^{2}-1)\le-2(\sum_i\lambda_i^{2})^{1/2}\sqrt t]\le e^{-t}$.
\end{proof}

The residual of \Cref{lem:refresh} is cubic in a random matrix whose sixth
moment is bounded by the following lemma; it is proved here because its
hypothesis is tailored to the centered retention variables of
\Cref{eq:references} below.
\begin{lemma}[sixth moment of a centered rank-one sum]
\label{lem:moment}
Let $\QQ_i\defeq\zeta_i\uu_i\uu_i^{\top}$, $i\in I$, with $\uu_i\in\R^{d}$
and independent real random variables $\zeta_i$ with $\Ex[\zeta_i]=0$ and
$\zeta_i\ge-1$ almost surely, such that
$\sum_{i\in I}\uu_i\uu_i^{\top}\preceq\II$ and, for some $a\le1$,
\[
\Ex[|\zeta_i|^{p}]\,\norm{\uu_i}^{2p-2}\ \le\ a^{p-1}
\qquad\text{for every }i\in I\text{ and }p=2,\dots,6 .
\]
Then
$\Ex\bigl[\norm{\sum_{i\in I}\QQ_i}_{S_6}^{6}\bigr]
=\Ex\bigl[\tr[(\sum_{i\in I}\QQ_i)^{6}]\bigr]\le10^{2}a^{3}d$.
\end{lemma}
\begin{proof}
The equality holds because $\sum_i\QQ_i$ is symmetric.
Since $(\uu_i\uu_i^{\top})^{p}=\norm{\uu_i}^{2p-2}\uu_i\uu_i^{\top}$,
\[
\sum_i\Ex[\QQ_i^{p}]=\sum_i\Ex[\zeta_i^{p}]\norm{\uu_i}^{2p-2}\uu_i\uu_i^{\top},
\qquad
\bigl|\Ex[\zeta_i^{p}]\norm{\uu_i}^{2p-2}\bigr|\le a^{p-1},
\]
so for $p=2,\dots,6$
\begin{align*}
-a^{p-1}\II\preceq-a^{p-1}\sum_i\uu_i\uu_i^{\top}
&\preceq\sum_i\Ex[\QQ_i^{p}]\preceq a^{p-1}\sum_i\uu_i\uu_i^{\top}\preceq a^{p-1}\II,\\
\Bigl|\tr\Bigl[\sum_i\Ex[\QQ_i^{p}]\Bigr]\Bigr|&\le a^{p-1}\tr\Bigl[\sum_i\uu_i\uu_i^{\top}\Bigr]\le a^{p-1}d .
\end{align*}
Four consequences are used repeatedly.
\begin{itemize}[leftmargin=*]
\item $\sum_{i,j}(\uu_i^{\top}\uu_j)^{2}=\tr[(\sum_i\uu_i\uu_i^{\top})^{2}]\le\tr[\sum_i\uu_i\uu_i^{\top}]\le d$,
as $0\preceq\sum_i\uu_i\uu_i^{\top}\preceq\II$.
\item $\tr[(\sum_i\Ex[\QQ_i^{2}])^{q}]\le a^{q-1}\tr[\sum_i\Ex[\QQ_i^{2}]]\le a^{q}d$
for $q=1,2,3$, as $0\preceq\sum_i\Ex[\QQ_i^{2}]\preceq a\II$.
\item For a fixed symmetric $\NN$ with $0\preceq\NN\preceq a'\II$,
\[
\sum_i\Ex[\QQ_i\NN\QQ_i]
=\sum_i\Ex[\zeta_i^{2}]\,(\uu_i^{\top}\NN\uu_i)\,\uu_i\uu_i^{\top}
\preceq a'\sum_i\Ex[\zeta_i^{2}]\norm{\uu_i}^{2}\uu_i\uu_i^{\top}
=a'\sum_i\Ex[\QQ_i^{2}] .
\]
\item For all $i,j$ (not necessarily distinct),
$\tr[\Ex[\QQ_i^{2}]\,\Ex[\QQ_j^{2}]]\ge0$ and
$\sum_{i,j}\tr[\Ex[\QQ_i^{2}]\,\Ex[\QQ_j^{2}]]=\tr[(\sum_i\Ex[\QQ_i^{2}])^{2}]\le a^{2}d$;
for $i\neq j$, $\tr[\Ex[\QQ_i^{2}]\,\Ex[\QQ_j^{2}]]=\Ex[\norm{\QQ_i\QQ_j}_F^{2}]$
by independence.
\end{itemize}

\emph{Expansion.}
$\Ex[\tr[(\sum_i\QQ_i)^{6}]]=\sum_{i_1,\dots,i_6\in I}\Ex[\tr[\QQ_{i_1}\cdots\QQ_{i_6}]]$,
and a term vanishes unless every index occurring in it occurs at least
twice, because the $\QQ_i$ are independent and centered.
Group the remaining terms by the partition of the six cyclic positions
into blocks of equal indices, with distinct indices on distinct blocks and
every block of size at least $2$.
For each partition we exhibit a nonnegative function of the block indices
that bounds the absolute value of the term when the block indices are
distinct, and sum it over \emph{all} tuples of block indices; this
dominates the sum over distinct tuples.
The partitions are, up to rotation, the following $1+15+10+15=41$.
\begin{itemize}[leftmargin=*,itemsep=1pt]
\item One block, the word $iiiiii$:
$\sum_i\Ex[\tr[\QQ_i^{6}]]=\tr[\sum_i\Ex[\QQ_i^{6}]]\le a^{5}d$.
\item Blocks of sizes $4$ and $2$ ($15$ partitions), the words $iiiijj$,
$iiijij$ and $iijiij$ and their rotations.
For $iiiijj$,
\[
\sum_{i,j}\tr[\Ex[\QQ_i^{4}]\,\Ex[\QQ_j^{2}]]
=\tr\Bigl[\Bigl(\sum_i\Ex[\QQ_i^{4}]\Bigr)\Bigl(\sum_j\Ex[\QQ_j^{2}]\Bigr)\Bigr]
\le a\,\tr\Bigl[\sum_i\Ex[\QQ_i^{4}]\Bigr]\le a^{4}d ,
\]
as both matrices are positive semidefinite and the second is at most
$a\II$.
For $iiijij$ and $iijiij$,
\[
\tr[\QQ_i^{3}\QQ_j\QQ_i\QQ_j]=\tr[\QQ_i^{2}\QQ_j\QQ_i^{2}\QQ_j]
=\zeta_i^{4}\zeta_j^{2}\norm{\uu_i}^{4}(\uu_i^{\top}\uu_j)^{4}
\le\zeta_i^{4}\norm{\uu_i}^{6}\cdot\zeta_j^{2}\norm{\uu_j}^{2}\cdot(\uu_i^{\top}\uu_j)^{2}
\]
by $(\uu_i^{\top}\uu_j)^{2}\le\norm{\uu_i}^{2}\norm{\uu_j}^{2}$; for $i\neq j$
the expectation is at most $a^{3}\cdot a\cdot(\uu_i^{\top}\uu_j)^{2}$ by
independence and the moment hypothesis, and summing over all $i,j$ gives
at most $a^{4}d$.
\item Blocks of sizes $3$ and $3$ ($10$ partitions), the words $iiijjj$,
$iijijj$ and $ijijij$ and their rotations.
For $iiijjj$, with $\Ex[\QQ_i^{3}]=\Ex[\zeta_i^{3}]\norm{\uu_i}^{4}\uu_i\uu_i^{\top}$
and $|\Ex[\zeta_i^{3}]|\norm{\uu_i}^{4}\le a^{2}$,
\[
\sum_{i,j}\bigl|\tr[\Ex[\QQ_i^{3}]\,\Ex[\QQ_j^{3}]]\bigr|
\le\sum_{i,j}a^{2}\cdot a^{2}\cdot(\uu_i^{\top}\uu_j)^{2}\le a^{4}d .
\]
For $iijijj$ and $ijijij$,
\begin{align*}
\bigl|\tr[\QQ_i^{2}\QQ_j\QQ_i\QQ_j^{2}]\bigr|
&=|\zeta_i|^{3}|\zeta_j|^{3}\norm{\uu_i}^{2}\norm{\uu_j}^{2}(\uu_i^{\top}\uu_j)^{4}
\le|\zeta_i|^{3}\norm{\uu_i}^{4}\cdot|\zeta_j|^{3}\norm{\uu_j}^{4}\cdot(\uu_i^{\top}\uu_j)^{2},\\
\bigl|\tr[(\QQ_i\QQ_j)^{3}]\bigr|
&=|\zeta_i|^{3}|\zeta_j|^{3}(\uu_i^{\top}\uu_j)^{6}
\le|\zeta_i|^{3}\norm{\uu_i}^{4}\cdot|\zeta_j|^{3}\norm{\uu_j}^{4}\cdot(\uu_i^{\top}\uu_j)^{2},
\end{align*}
with expectations at most $a^{2}\cdot a^{2}\cdot(\uu_i^{\top}\uu_j)^{2}$ for
$i\neq j$, each summing to at most $a^{4}d$.
\item Three blocks of size $2$ with two equal letters cyclically adjacent
($11$ partitions): rotate the word to $kkW$ with $W$ a word of length four
in $i,j$, each letter twice.
Taking the expectation over $\zeta_k$ first (for $k\notin\{i,j\}$) and
summing over $k$ leaves $\tr[(\sum_k\Ex[\QQ_k^{2}])\,W]$, and the six
possible $W$ give:
\begin{itemize}[leftmargin=*,itemsep=1pt]
\item $W\in\{\QQ_i^{2}\QQ_j^{2},\QQ_j^{2}\QQ_i^{2}\}$:
$\sum_{i,j}\tr[(\sum_k\Ex[\QQ_k^{2}])\,\Ex[\QQ_i^{2}]\,\Ex[\QQ_j^{2}]]
=\tr[(\sum_k\Ex[\QQ_k^{2}])^{3}]\le a^{3}d$;
\item $W\in\{\QQ_i\QQ_j^{2}\QQ_i,\QQ_j\QQ_i^{2}\QQ_j\}$: by the third
consequence with $\NN=\sum_j\Ex[\QQ_j^{2}]\preceq a\II$,
\[
\sum_i\Ex\Bigl[\QQ_i\Bigl(\sum_j\Ex[\QQ_j^{2}]\Bigr)\QQ_i\Bigr]\preceq a\sum_i\Ex[\QQ_i^{2}],
\]
so the sum is at most $a\,\tr[(\sum_k\Ex[\QQ_k^{2}])^{2}]\le a^{3}d$;
\item $W\in\{\QQ_i\QQ_j\QQ_i\QQ_j,\QQ_j\QQ_i\QQ_j\QQ_i\}$:
$|\tr[(\sum_k\Ex[\QQ_k^{2}])(\QQ_i\QQ_j)^{2}]|\le\norm{\sum_k\Ex[\QQ_k^{2}]}\cdot\norm{\QQ_i\QQ_j}_F^{2}\le a\norm{\QQ_i\QQ_j}_F^{2}$,
whose expectation summed over $i\neq j$ is at most
$a\sum_{i,j}\tr[\Ex[\QQ_i^{2}]\Ex[\QQ_j^{2}]]\le a^{3}d$ by the fourth
consequence.
\end{itemize}
\item Three blocks of size $2$ with no two equal letters cyclically adjacent
($4$ partitions): the word $ijkijk$ and the three rotations-with-relabeling
of $ijkikj$.
For $ijkijk$, with $i,j,k$ distinct,
\[
\bigl|\tr[(\QQ_i\QQ_j\QQ_k)^{2}]\bigr|
\le\norm{\QQ_i\QQ_j\QQ_k}_F^{2}
=\tr[\QQ_k\QQ_j\QQ_i^{2}\QQ_j\QQ_k],
\]
and taking the expectation over $\zeta_i$, then $\zeta_j$, then $\zeta_k$,
with the third consequence applied twice ($\sum_i\Ex[\QQ_i^{2}]\preceq a\II$,
then $\sum_j\Ex[\QQ_j(\sum_i\Ex[\QQ_i^{2}])\QQ_j]\preceq a\sum_j\Ex[\QQ_j^{2}]\preceq a^{2}\II$),
gives
\[
\sum_{i,j,k}\Ex\bigl[\tr[\QQ_k\QQ_j\QQ_i^{2}\QQ_j\QQ_k]\bigr]
\le a^{2}\,\tr\Bigl[\sum_k\Ex[\QQ_k^{2}]\Bigr]\le a^{3}d .
\]
For $ijkikj$, a rotation gives
$\tr[\QQ_i\QQ_j\QQ_k\QQ_i\QQ_k\QQ_j]=\tr[(\QQ_j\QQ_i\QQ_j)(\QQ_k\QQ_i\QQ_k)]$,
and $\QQ_j\QQ_i\QQ_j=\zeta_j^{2}(\uu_j^{\top}\QQ_i\uu_j)\uu_j\uu_j^{\top}$ with
$\uu_j^{\top}\QQ_i\uu_j=\zeta_i(\uu_i^{\top}\uu_j)^{2}$, so
\begin{align*}
\bigl|\tr[(\QQ_j\QQ_i\QQ_j)(\QQ_k\QQ_i\QQ_k)]\bigr|
&=\zeta_j^{2}\zeta_k^{2}|\uu_j^{\top}\QQ_i\uu_j|\,|\uu_k^{\top}\QQ_i\uu_k|\,(\uu_j^{\top}\uu_k)^{2}\\
&\le\zeta_j^{2}\norm{\uu_j}^{2}|\uu_j^{\top}\QQ_i\uu_j|\cdot\zeta_k^{2}\norm{\uu_k}^{2}|\uu_k^{\top}\QQ_i\uu_k| .
\end{align*}
Conditionally on $\zeta_i$ the factors indexed by $j\neq k$ are
independent, so the sum over distinct $j,k\neq i$ of the conditional
expectations is at most
\begin{align*}
\Ex\Bigl[\Bigl(\sum_{j}\zeta_j^{2}\norm{\uu_j}^{2}|\uu_j^{\top}\QQ_i\uu_j|\Bigr)^{2}\Bigm|\zeta_i\Bigr]
&\le\sum_j\Ex[\zeta_j^{4}]\norm{\uu_j}^{4}(\uu_j^{\top}\QQ_i\uu_j)^{2}
+\Bigl(\sum_j\Ex[\zeta_j^{2}]\norm{\uu_j}^{2}|\uu_j^{\top}\QQ_i\uu_j|\Bigr)^{2}.
\end{align*}
With $(\uu_j^{\top}\QQ_i\uu_j)^{2}=\zeta_i^{2}(\uu_i^{\top}\uu_j)^{4}\le\zeta_i^{2}\norm{\uu_i}^{2}\norm{\uu_j}^{2}(\uu_i^{\top}\uu_j)^{2}$
and $\Ex[\zeta_j^{4}]\norm{\uu_j}^{6}\le a^{3}$, the first sum is at most
\[
a^{3}\zeta_i^{2}\norm{\uu_i}^{2}\,\uu_i^{\top}\Bigl(\sum_j\uu_j\uu_j^{\top}\Bigr)\uu_i
\le a^{3}\zeta_i^{2}\norm{\uu_i}^{4},
\]
and with $\Ex[\zeta_j^{2}]\norm{\uu_j}^{2}\le a$ the second is at most
$(a|\zeta_i|\sum_j(\uu_i^{\top}\uu_j)^{2})^{2}\le a^{2}\zeta_i^{2}\norm{\uu_i}^{4}$.
Taking the expectation over $\zeta_i$ and summing over $i$, both are at
most $a^{2}\sum_i\Ex[\zeta_i^{2}]\norm{\uu_i}^{4}=a^{2}\tr[\sum_i\Ex[\QQ_i^{2}]]\le a^{3}d$,
so this word contributes at most $2a^{3}d$.
\end{itemize}
Each of the $41$ partitions contributes at most $2a^{3}d$, as $a\le1$, and
$82a^{3}d\le10^{2}a^{3}d$.
\end{proof}

\subsection{The residual estimator}

\paragraph{The sampled references.}
The estimator works on two sparse random matrices built from $G$ by
independent retention of its edges.
Each weight is split as $\ww_e=\ww^{\mathrm{s}}_e+\ww^{\mathrm{x}}_e$ into a
sampled part $\ww^{\mathrm{s}}_e\ge0$ and an exactly kept part
$\ww^{\mathrm{x}}_e\ge0$ (one of them may be zero), and
$E_{\mathrm{s}}\defeq\{e\in E:\ww^{\mathrm{s}}_e>0\}$.
Every $e\in E_{\mathrm{s}}$ has an \emph{intensity} $\varrho_e>0$ and is
\emph{retained} independently with probability $p_e\defeq1-e^{-\varrho_e}$.
The \emph{centered retention variable} $\zeta_e$ of $e$ and the
\emph{sparse reference} $\LLtil$ and the \emph{heavy reference} $\MMtil$ are
\begin{equation}\label{eq:references}
\begin{aligned}
\zeta_e&\defeq
\begin{cases}
1/p_e-1,&e\in E_{\mathrm{s}}\text{ retained},\\
-1,&e\in E_{\mathrm{s}}\text{ not retained},\\
0,&e\notin E_{\mathrm{s}},
\end{cases}\\
\LLtil&\defeq\sum_{e\in E}\bigl(\ww^{\mathrm{x}}_e+(1+\zeta_e)\ww^{\mathrm{s}}_e\bigr)\cchi_e\cchi_e^{\top},\\
\MMtil&\defeq\DD(\Gl)+\sum_{e\in\Gh}\bigl(\ww^{\mathrm{x}}_e+(1+\zeta_e)\ww^{\mathrm{s}}_e\bigr)\cchi_e\cchi_e^{\top},\\
\QQ_e&\defeq\zeta_e\ww^{\mathrm{s}}_e\LL(G)^{\dag/2}\cchi_e\cchi_e^{\top}\LL(G)^{\dag/2},
\qquad
\QQ^{\mathrm{heavy}}_e\defeq\zeta_e\ww^{\mathrm{s}}_e\MM^{\dag/2}\cchi_e\cchi_e^{\top}\MM^{\dag/2}.
\end{aligned}
\end{equation}
Thus $\Ex[\zeta_e]=0$, $\zeta_e\ge-1$, $(1+\zeta_e)\ww^{\mathrm{s}}_e$ is
$\ww^{\mathrm{s}}_e/p_e$ if $e$ is retained and $0$ otherwise,
$\Ex[\LLtil]=\LL(G)$, $\Ex[\MMtil]=\MM$, and
\begin{equation}\label{eq:deviations}
\LL(G)^{\dag/2}\LLtil\LL(G)^{\dag/2}-\PPi=\sum_{e\in E}\QQ_e,
\qquad
\MM^{\dag/2}\MMtil\MM^{\dag/2}-\MM^{\dag/2}\MM\MM^{\dag/2}=\sum_{e\in\Gh}\QQ^{\mathrm{heavy}}_e ,
\end{equation}
since $\sum_e\ww_e\LL(G)^{\dag/2}\cchi_e\cchi_e^{\top}\LL(G)^{\dag/2}=\PPi$ and
$\DD(\Gl)+\sum_{e\in\Gh}\ww_e\cchi_e\cchi_e^{\top}=\MM$.
We write $\LLtil=\BBtil^{\top}\WWtil\BBtil$, where $\BBtil$ has the row
$\cchi_e^{\top}$ for every edge $e$ with a positive coefficient in
\Cref{eq:references} and $\WWtil$ is the diagonal matrix of these
coefficients; $E(\LLtil)$ denotes this set of edges.
The intensities are required to dominate the effective resistances:
\begin{equation}\label{eq:intensity}
\varrho_e\ \ge\ 10^{8}\log n\cdot\ww^{\mathrm{s}}_e\Reff(G)_e
\qquad\text{for every }e\in E_{\mathrm{s}} .
\end{equation}
The references are \emph{accurate} if
\begin{equation}\label{eq:accurate}
\LLtil\approx_{10^{-3}}\LL(G)
\qquad\text{and}\qquad
\MMtil\approx_{10^{-3}}\MM .
\end{equation}

\paragraph{The attempt.}
The pseudocode of $\textsc{Attempt}$ is given in \Cref{fig:attempt}.
\begin{figure}[H]\centering
\framebox{\begin{minipage}{0.96\linewidth}\small
$(\xxrand_j,\xxres_j,\xxproxy_j)_{j\le10^{3}\log n}=\textsc{Attempt}(\LLtil,\MMtil,\ZZ(\LLtil),\ZZ(\MMtil))$\\[2pt]
\underline{Input}: the references $\LLtil=\BBtil^{\top}\WWtil\BBtil$ and
$\MMtil$; the fixed operators $\ZZ(\LLtil)$, $\ZZ(\MMtil)$;\\
reads the light degrees $\dd(\Gl)$.\\
\underline{Output}: the vectors $\xxrand_j,\xxres_j,\xxproxy_j\in\R^{V}$,
$j\le10^{3}\log n$;\\
the verdict \emph{accepted} or \emph{rejected}.
\begin{enumerate}[leftmargin=*,itemsep=1pt]
\item Draw $\gg_j\in\R^{E(\LLtil)}$, $j\le10^{3}\log n$, with independent
standard Gaussian entries.
\item For every $j\le10^{3}\log n$:
\begin{enumerate}[label=(\alph*),leftmargin=*]
\item $\xxrand_j\leftarrow\ZZ(\LLtil)\BBtil^{\top}\WWtil^{1/2}\gg_j$.
\item $\xxres_j\leftarrow(\II-\ZZ(\MMtil)\LLtil)^{3}\xxrand_j$.
\item $\xxproxy_j\leftarrow\xxrand_j-\xxres_j$.
\end{enumerate}
\item The verdict is \emph{accepted} iff
\begin{align*}
\frac1{10^{3}\log n}\sum_j\xxrand_j{}^{\top}\LLtil\xxrand_j&\le1.1n,\\
\frac1{10^{3}\log n}\sum_j\xxres_j{}^{\top}\LLtil\xxres_j&\le10^{8}(n/\dheavy+n/\log^{3}n),\\
\frac1{10^{3}\log n}\sum_j\xxproxy_j{}^{\top}\DD(\Gl)\xxproxy_j&\le10^{3}n .
\end{align*}
\item Return $(\xxrand_j,\xxres_j,\xxproxy_j)_{j\le10^{3}\log n}$ and the
verdict.
\end{enumerate}
\end{minipage}}
\caption{One attempt of a refresh: $10^{3}\log n$ random vertex vectors,
their cubic residuals and proxies, and the three acceptance tests.
Its guarantees are in \Cref{lem:refresh}.}
\label{fig:attempt}
\end{figure}

The analysis conjugates one application of $\II-\ZZ(\MMtil)\LLtil$ by
$\LL(G)^{1/2}$ and splits it into a deterministic part and a random part:
\begin{equation}\label{eq:cubic}
\begin{aligned}
\YYbar&\defeq\PPi-\LL(G)^{1/2}\MM^{\dag}\LL(G)^{1/2},\\
\YYtil&\defeq-\LL(G)^{1/2}\MM^{\dag}\LL(G)^{1/2}\Bigl(\sum_{e\in E}\QQ_e\Bigr)\\
&\quad-\LL(G)^{1/2}\bigl(\ZZ(\MMtil)-\MMtil^{\dag}\bigr)\LL(G)^{1/2}\,\LL(G)^{\dag/2}\LLtil\LL(G)^{\dag/2}\\
&\quad+\LL(G)^{1/2}\MMtil^{\dag}\MM^{1/2}\Bigl(\sum_{e\in\Gh}\QQ^{\mathrm{heavy}}_e\Bigr)\MM^{\dag/2}\LL(G)^{1/2}\,\LL(G)^{\dag/2}\LLtil\LL(G)^{\dag/2}.
\end{aligned}
\end{equation}
Here and below $\ZZ(\MMtil)$ stands for $\MM^{\dag/2}\MM\MM^{\dag/2}\ZZ(\MMtil)\MM^{\dag/2}\MM\MM^{\dag/2}$
and $\ZZ(\LLtil)$ for $\PPi\ZZ(\LLtil)\PPi$; the first paragraph of the
proof below justifies this.

\begin{lemma}[the refresh]
\label{lem:refresh}
Let $G=(V,E,\ww)$ be connected with the split of \Cref{def:heavy},
$\log n\ge20$, and let $\LLtil,\MMtil$ be the references of
\Cref{eq:references} with intensities satisfying \Cref{eq:intensity}.
Let $\ZZ(\LLtil)\leftarrow\textsc{LsolverKLPSS}(\LLtil,\frac1{40})$ and
$\ZZ(\MMtil)\leftarrow\textsc{LsolverKLPSS}(\MMtil,1/\log n)$, and call the
operators \emph{accurate} if they satisfy the conclusion of \Cref{lem:chain}:
$\ZZ(\LLtil)$ is symmetric with $\ZZ(\LLtil)\approx_{1/40}\LLtil^{\dag}$ on
the range of $\LLtil$ and $\ZZ(\MMtil)$ is symmetric with
$\ZZ(\MMtil)\approx_{1/\log n}\MMtil^{\dag}$ on the range of $\MMtil$.
An \emph{attempt} is a run of $\textsc{Attempt}(\LLtil,\MMtil,\ZZ(\LLtil),\ZZ(\MMtil))$
(\Cref{fig:attempt}) whose Gaussian vectors are independent of the
retention variables and of the randomness of the two solver calls.
\begin{enumerate}[label=(\alph*),leftmargin=*]
\item (acceptance)
The probability, over the retention variables, that the references are
accurate and
\begin{equation}\label{eq:goodrefs}
\Bigl\|\sum_{e\in E}\QQ_e\Bigr\|_{S_6}^{6}
+\Bigl\|\sum_{e\in\Gh}\QQ^{\mathrm{heavy}}_e\Bigr\|_{S_6}^{6}
\ >\ 10^{-16}\,\frac{n}{\log^{3}n}
\end{equation}
is at most $10^{-2}$.
If the references and the operators are accurate and \Cref{eq:goodrefs}
fails, an attempt is accepted with probability at least $1-3n^{-15}$ over
its Gaussian vectors.
If the references are accurate, an accepted attempt satisfies
\begin{equation}\label{eq:energies}
\begin{aligned}
\frac1{10^{3}\log n}\sum_j\xxrand_j{}^{\top}\LL(G)\xxrand_j&\le1.2n,\\
\frac1{10^{3}\log n}\sum_j\xxres_j{}^{\top}\LL(G)\xxres_j&\le2\cdot10^{8}(n/\dheavy+n/\log^{3}n),\\
\frac1{10^{3}\log n}\sum_j\xxproxy_j{}^{\top}\DD(\Gl)\xxproxy_j&\le10^{3}n .
\end{aligned}
\end{equation}
Consequently an attempt with fresh retention variables, fresh solver calls
and fresh Gaussian vectors is accepted with probability at least
$1-10^{-2}-3n^{-15}$ minus the probability that the references or the
operators fail to be accurate.
\item (all pairs)
If the references and the operators are accurate, then with probability
$1-n^{-14}$ over the Gaussian vectors, for all $u,v\in V$,
\[
\frac1{10^{3}\log n}\sum_j\bigl((\xxrand_j)_u-(\xxrand_j)_v\bigr)^{2}
\in\bigl[\tfrac12,2\bigr]\cdot(\oone_u-\oone_v)^{\top}\LL(G)^{\dag}(\oone_u-\oone_v).
\]
\item (cost)
If $\LLtil$ and $\MMtil$ have $O(n\log n)$ nonzero entries, an attempt
takes $\Odouble(n\log^{2}n)$ time beyond the two calls of
$\textsc{LsolverKLPSS}$, on a machine with a unit-cost primitive returning
a standard Gaussian: $O(n\log^{2}n)$ Gaussian draws and arithmetic
operations, $10^{3}\log n$ applications of $\ZZ(\LLtil)$ and
$3\cdot10^{3}\log n$ applications of $\ZZ(\MMtil)$.
\end{enumerate}
\end{lemma}
\begin{proof}
\emph{Conventions.}
By \Cref{lem:topk}(d), $\MM^{\dag/2}\MM\MM^{\dag/2}$ is the orthogonal
projection onto the range of $\MM$, which is $\R^{V}$ if $\Gl\neq\emptyset$
and $\oone^{\perp}$ if $\Gl=\emptyset$; when the references are accurate,
$\MMtil$ has the same range as $\MM$, and $\LLtil$ has range $\oone^{\perp}$.
In \Cref{fig:attempt}, $\ZZ(\LLtil)$ is applied to the vector
$\BBtil^{\top}\WWtil^{1/2}\gg_j\perp\oone$ only, and $\ZZ(\MMtil)$ to vectors
of the form $\LLtil\xx\perp\oone$ only; their outputs enter the attempt and
the statement only through the quadratic forms of $\LLtil$ and $\LL(G)$,
through $\DD(\Gl)$, which vanishes unless the range of $\MM$ is $\R^{V}$,
through differences of two coordinates, and through further applications
of $\II-\ZZ(\MMtil)\LLtil$, which maps $\oone$ to $\oone$.
Hence replacing $\ZZ(\LLtil)$ by $\PPi\ZZ(\LLtil)\PPi$ and $\ZZ(\MMtil)$ by
$\MM^{\dag/2}\MM\MM^{\dag/2}\ZZ(\MMtil)\MM^{\dag/2}\MM\MM^{\dag/2}$ changes
every $\xxrand_j$, $\xxres_j$, $\xxproxy_j$ by a multiple of $\oone$ and
none of the quantities of the statement; we make this replacement, as
announced after \Cref{eq:cubic}.
The accuracy of the operators then reads, as Loewner inequalities on $\R^{V}$,
\begin{equation}\label{eq:ops}
\tfrac{39}{40}\LLtil^{\dag}\preceq\ZZ(\LLtil)\preceq\tfrac{41}{40}\LLtil^{\dag},
\qquad
\Bigl(1-\frac1{\log n}\Bigr)\MMtil^{\dag}\preceq\ZZ(\MMtil)\preceq\Bigl(1+\frac1{\log n}\Bigr)\MMtil^{\dag},
\end{equation}
and $\ZZ(\LLtil)$, $\ZZ(\MMtil)$ are supported on the ranges of $\LLtil$,
$\MMtil$.

\emph{The cubic identity.}
Since $\LLtil\oone=0$, $(\II-\ZZ(\MMtil)\LLtil)\oone=\oone$, so
$\LL(G)^{1/2}(\II-\ZZ(\MMtil)\LLtil)(\II-\PPi)=0$ and, as
$\PPi=\LL(G)^{\dag/2}\LL(G)^{1/2}$,
\[
\LL(G)^{1/2}(\II-\ZZ(\MMtil)\LLtil)
=\LL(G)^{1/2}(\II-\ZZ(\MMtil)\LLtil)\PPi
=\bigl(\LL(G)^{1/2}(\II-\ZZ(\MMtil)\LLtil)\LL(G)^{\dag/2}\bigr)\LL(G)^{1/2}.
\]
Iterating three times,
\begin{equation}\label{eq:iterate}
\LL(G)^{1/2}(\II-\ZZ(\MMtil)\LLtil)^{3}
=\bigl(\LL(G)^{1/2}(\II-\ZZ(\MMtil)\LLtil)\LL(G)^{\dag/2}\bigr)^{3}\LL(G)^{1/2}.
\end{equation}
We show $\LL(G)^{1/2}(\II-\ZZ(\MMtil)\LLtil)\LL(G)^{\dag/2}=\YYbar+\YYtil$
when the references are accurate.
As $\PPi\LLtil\PPi=\LLtil$, $\LLtil=\LL(G)^{1/2}(\LL(G)^{\dag/2}\LLtil\LL(G)^{\dag/2})\LL(G)^{1/2}$,
so
\[
\LL(G)^{1/2}(\II-\ZZ(\MMtil)\LLtil)\LL(G)^{\dag/2}
=\PPi-\LL(G)^{1/2}\ZZ(\MMtil)\LL(G)^{1/2}\,\LL(G)^{\dag/2}\LLtil\LL(G)^{\dag/2}.
\]
Since $\MM$ and $\MMtil$ have the same range,
$\MMtil^{\dag}-\MM^{\dag}=\MMtil^{\dag}(\MM-\MMtil)\MM^{\dag}$, and by
\Cref{eq:deviations} and
$\MM^{1/2}(\MM^{\dag/2}\MMtil\MM^{\dag/2}-\MM^{\dag/2}\MM\MM^{\dag/2})\MM^{1/2}=\MMtil-\MM$,
\[
\ZZ(\MMtil)=\MM^{\dag}+\bigl(\ZZ(\MMtil)-\MMtil^{\dag}\bigr)
-\MMtil^{\dag}\MM^{1/2}\Bigl(\sum_{e\in\Gh}\QQ^{\mathrm{heavy}}_e\Bigr)\MM^{1/2}\MM^{\dag}.
\]
Substituting this, $\LL(G)^{\dag/2}\LLtil\LL(G)^{\dag/2}=\PPi+\sum_e\QQ_e$
for the first term, $\LL(G)^{1/2}\PPi=\LL(G)^{1/2}$ and
$\MM^{1/2}\MM^{\dag}=\MM^{\dag/2}$ gives exactly $\YYbar+\YYtil$ of
\Cref{eq:cubic}.
Consequently, by \Cref{eq:iterate}, for every $j$
\begin{equation}\label{eq:resenergy}
\xxres_j{}^{\top}\LL(G)\xxres_j=\norm{(\YYbar+\YYtil)^{3}\LL(G)^{1/2}\xxrand_j}^{2},
\qquad
\LL(G)^{1/2}(\II-\ZZ(\MMtil)\LLtil)^{3}\LL(G)^{\dag/2}=(\YYbar+\YYtil)^{3},
\end{equation}
the second because $\YYbar\PPi=\YYbar$ and $\YYtil\PPi=\YYtil$.

\emph{The deterministic part.}
For $\uu\in\R^{V}$, $\LL(G)^{1/2}\uu\perp\oone$, so \Cref{lem:topk}(e) gives
\[
\uu^{\top}\LL(G)^{1/2}\MM^{\dag}\LL(G)^{1/2}\uu
\le2\uu^{\top}\LL(G)^{1/2}\LL(G)^{\dag}\LL(G)^{1/2}\uu=2\norm{\PPi\uu}^{2};
\]
hence $0\preceq\LL(G)^{1/2}\MM^{\dag}\LL(G)^{1/2}\preceq2\PPi$,
$-\PPi\preceq\YYbar\preceq\PPi$ and $\norm{\YYbar}\le1$.
With $\LL(G)=\MM-\AA(\Gl)$ and $r\defeq\mathrm{rank}(\MM)\in\{n-1,n\}$,
\begin{align*}
\norm{\YYbar}_F^{2}
&=\tr[\PPi]-2\tr[\MM^{\dag}\LL(G)]+\tr[\MM^{\dag}\LL(G)\MM^{\dag}\LL(G)]\\
&=(n-1)-2\bigl(r-\tr[\MM^{\dag}\AA(\Gl)]\bigr)\\
&\quad+\bigl(r-2\tr[\MM^{\dag}\AA(\Gl)]+\tr[\MM^{\dag}\AA(\Gl)\MM^{\dag}\AA(\Gl)]\bigr)\\
&=(n-1-r)+\tr[\MM^{\dag}\AA(\Gl)\MM^{\dag}\AA(\Gl)]
\le\frac{32n}{\dheavy}
\end{align*}
by \Cref{lem:topk}(g), where we used
$\tr[\MM^{\dag}\MM]=r$ and $\MM^{\dag}\MM\MM^{\dag}\AA(\Gl)=\MM^{\dag}\AA(\Gl)$.

\emph{Norm bounds.}
Suppose the references and the operators are accurate.
From \Cref{eq:deviations} and \Cref{eq:accurate},
$\LL(G)^{\dag/2}\LLtil\LL(G)^{\dag/2}$ lies between $0.999\PPi$ and
$1.001\PPi$ and $\MM^{\dag/2}\MMtil\MM^{\dag/2}$ between
$0.999\MM^{\dag/2}\MM\MM^{\dag/2}$ and $1.001\MM^{\dag/2}\MM\MM^{\dag/2}$,
so
\begin{align*}
\Bigl\|\sum_{e\in E}\QQ_e\Bigr\|&\le10^{-3},\\
\Bigl\|\sum_{e\in\Gh}\QQ^{\mathrm{heavy}}_e\Bigr\|&\le10^{-3},\\
\norm{\LL(G)^{\dag/2}\LLtil\LL(G)^{\dag/2}}&\le1.001,\\
\norm{\LL(G)^{1/2}\MM^{\dag}\LL(G)^{1/2}}&\le2,\\
\norm{\MM^{\dag/2}\LL(G)^{1/2}}=\norm{\LL(G)^{1/2}\MM^{\dag}\LL(G)^{1/2}}^{1/2}&\le1.5,\\
\norm{\LL(G)^{1/2}\MMtil^{\dag/2}}^{2}=\norm{\MMtil^{\dag/2}\LL(G)\MMtil^{\dag/2}}
&\le2\norm{\MMtil^{\dag/2}\MM\MMtil^{\dag/2}}\le\frac{2}{0.999}\le2.01,\\
\norm{\MMtil^{\dag/2}\MM^{1/2}}^{2}=\norm{\MMtil^{\dag/2}\MM\MMtil^{\dag/2}}&\le\frac1{0.999}\le1.002,\\
\norm{\LL(G)^{1/2}\MMtil^{\dag}\MM^{1/2}}
&\le\norm{\LL(G)^{1/2}\MMtil^{\dag/2}}\cdot\norm{\MMtil^{\dag/2}\MM^{1/2}}\\
&\le\sqrt{2.01\cdot1.002}\le1.5,
\end{align*}
using $\LL(G)\preceq2\MM$ (\Cref{lem:topk}(c)) and $\MM\preceq\MMtil/0.999$.
As $\ZZ(\MMtil)$ and $\MMtil^{\dag}$ are supported on the range of $\MMtil$,
$\ZZ(\MMtil)-\MMtil^{\dag}=\MMtil^{\dag/2}(\MMtil^{1/2}\ZZ(\MMtil)\MMtil^{1/2}-\MMtil^{\dag/2}\MMtil\MMtil^{\dag/2})\MMtil^{\dag/2}$,
and the middle factor has norm at most $1/\log n$ by \Cref{eq:ops}; so
\[
\norm{\LL(G)^{1/2}(\ZZ(\MMtil)-\MMtil^{\dag})\LL(G)^{1/2}}
\le\norm{\LL(G)^{1/2}\MMtil^{\dag/2}}^{2}\cdot\frac1{\log n}\le\frac{2.01}{\log n}.
\]
Applying these to the three terms of \Cref{eq:cubic}, with $\log n\ge20$,
\[
\norm{\YYtil}\le2\cdot10^{-3}+\frac{2.01\cdot1.001}{\log n}+1.5\cdot10^{-3}\cdot1.5\cdot1.001
\le2\cdot10^{-3}+0.101+2.3\cdot10^{-3}\le\tfrac12 .
\]
For the Schatten norm, $\norm{\NN\NN'}_{S_6}\le\norm{\NN}\cdot\norm{\NN'}_{S_6}$,
$\norm{\NN\NN'}_{S_6}\le\norm{\NN}_{S_6}\cdot\norm{\NN'}$ and
\[
\norm{\LL(G)^{\dag/2}\LLtil\LL(G)^{\dag/2}}_{S_6}\le1.001\norm{\PPi}_{S_6}\le1.001\,n^{1/6}
\]
give
\[
\norm{\YYtil}_{S_6}
\le2\Bigl\|\sum_{e\in E}\QQ_e\Bigr\|_{S_6}
+\frac{2.02\,n^{1/6}}{\log n}
+2.3\Bigl\|\sum_{e\in\Gh}\QQ^{\mathrm{heavy}}_e\Bigr\|_{S_6},
\]
and $(x+y+z)^{6}\le3^{5}(x^{6}+y^{6}+z^{6})$ for nonnegative $x,y,z$ gives
\begin{equation}\label{eq:S6}
\norm{\YYtil}_{S_6}^{6}
\le243\Bigl(149\Bigl(\Bigl\|\sum_{e\in E}\QQ_e\Bigr\|_{S_6}^{6}
+\Bigl\|\sum_{e\in\Gh}\QQ^{\mathrm{heavy}}_e\Bigr\|_{S_6}^{6}\Bigr)
+\frac{69\,n}{\log^{6}n}\Bigr),
\end{equation}
as $2^{6}\le2.3^{6}\le149$ and $2.02^{6}\le69$.

\emph{The sixth moments.}
\begin{claim}\label{clm:zeta}
For every $e\in E_{\mathrm{s}}$ and $p=2,\dots,6$,
$\Ex[|\zeta_e|^{p}]\,(\ww^{\mathrm{s}}_e\Reff(G)_e)^{p-1}\le(10^{-7}/\log n)^{p-1}$.
\end{claim}
\begin{proof}[Proof of \Cref{clm:zeta}]
$\Ex[|\zeta_e|^{p}]=(1-p_e)+p_e(1/p_e-1)^{p}$ and
$\ww^{\mathrm{s}}_e\Reff(G)_e\le10^{-8}\varrho_e/\log n$ by \Cref{eq:intensity}.
\begin{itemize}[leftmargin=*]
\item If $\varrho_e<1$: $p_e=1-e^{-\varrho_e}\ge(1-e^{-1})\varrho_e\ge0.63\varrho_e$
by concavity, and $\Ex[|\zeta_e|^{p}]\le1+p_e^{1-p}\le2p_e^{1-p}$, so
\[
\Ex[|\zeta_e|^{p}](\ww^{\mathrm{s}}_e\Reff(G)_e)^{p-1}
\le2\Bigl(\frac{\ww^{\mathrm{s}}_e\Reff(G)_e}{p_e}\Bigr)^{p-1}
\le2\Bigl(\frac{10^{-8}}{0.63\log n}\Bigr)^{p-1}
\le\Bigl(\frac{10^{-7}}{\log n}\Bigr)^{p-1}.
\]
\item If $\varrho_e\ge1$: $p_e\ge1-e^{-1}\ge0.63$ and $1/p_e-1=e^{-\varrho_e}/p_e$, so
$\Ex[|\zeta_e|^{p}]\le e^{-\varrho_e}+p_e^{1-p}e^{-p\varrho_e}\le e^{-\varrho_e}(1+0.63^{1-p}e^{-(p-1)})\le2e^{-\varrho_e}$,
the last because $0.63^{1-p}e^{-(p-1)}=(0.63e)^{1-p}\le1$; hence
\[
\Ex[|\zeta_e|^{p}](\ww^{\mathrm{s}}_e\Reff(G)_e)^{p-1}
\le2\varrho_e^{p-1}e^{-\varrho_e}\Bigl(\frac{10^{-8}}{\log n}\Bigr)^{p-1}
\le2\Bigl(\frac{p-1}{e}\Bigr)^{p-1}\Bigl(\frac{10^{-8}}{\log n}\Bigr)^{p-1}
\le\Bigl(\frac{10^{-7}}{\log n}\Bigr)^{p-1},
\]
as $\max_{\varrho>0}\varrho^{p-1}e^{-\varrho}=((p-1)/e)^{p-1}$ and
$2((p-1)/e)^{p-1}\le10^{p-1}$ for $p\le6$.
\qedhere
\end{itemize}
\end{proof}
\Cref{lem:moment} applies to $\QQ_e$, $e\in E_{\mathrm{s}}$ (the other
$\QQ_e$ vanish), with the vectors
$\uu_e=(\ww^{\mathrm{s}}_e)^{1/2}\LL(G)^{\dag/2}\cchi_e$, for which
$\norm{\uu_e}^{2}=\ww^{\mathrm{s}}_e\Reff(G)_e$ and
$\sum_e\uu_e\uu_e^{\top}\preceq\sum_e\ww_e\LL(G)^{\dag/2}\cchi_e\cchi_e^{\top}\LL(G)^{\dag/2}=\PPi\preceq\II$,
and $a=10^{-7}/\log n$ by \Cref{clm:zeta}; and to $\QQ^{\mathrm{heavy}}_e$,
$e\in\Gh\cap E_{\mathrm{s}}$, with the vectors
$\uu_e=(\ww^{\mathrm{s}}_e)^{1/2}\MM^{\dag/2}\cchi_e$, for which
$\norm{\uu_e}^{2}=\ww^{\mathrm{s}}_e\cchi_e^{\top}\MM^{\dag}\cchi_e\le2\ww^{\mathrm{s}}_e\Reff(G)_e$
(\Cref{lem:topk}(e)),
$\sum_{e\in\Gh}\uu_e\uu_e^{\top}\preceq\MM^{\dag/2}\LL(\Gh)\MM^{\dag/2}\preceq\MM^{\dag/2}\MM\MM^{\dag/2}\preceq\II$,
and $a=2\cdot10^{-7}/\log n$ (the factor $2^{p-1}$ from the norm bound).
Therefore
\[
\Ex\Bigl[\Bigl\|\sum_{e\in E}\QQ_e\Bigr\|_{S_6}^{6}
+\Bigl\|\sum_{e\in\Gh}\QQ^{\mathrm{heavy}}_e\Bigr\|_{S_6}^{6}\Bigr]
\le10^{2}\bigl(10^{-21}+8\cdot10^{-21}\bigr)\frac{n}{\log^{3}n}
=9\cdot10^{-19}\,\frac{n}{\log^{3}n},
\]
the expectation over the retention variables alone, and Markov's inequality
bounds the probability of \Cref{eq:goodrefs} by $9\cdot10^{-3}\le10^{-2}$.
This is the first sentence of (a).

\emph{The residual on good references.}
Suppose the references and the operators are accurate and
\Cref{eq:goodrefs} fails.
By \Cref{eq:S6} and $\log^{6}n\ge\log^{3}n$,
$\norm{\YYtil}_{S_6}^{6}\le243(149\cdot10^{-16}+69)\,n/\log^{3}n\le1.7\cdot10^{4}\,n/\log^{3}n$.
The matrix $(\YYbar+\YYtil)^{3}$ is the sum of eight words in $\YYbar$ and
$\YYtil$.
The word $\YYtil^{3}$ has $\norm{\YYtil^{3}}_F\le\norm{\YYtil}_{S_6}^{3}$
(H\"older's inequality for Schatten norms, $\frac16+\frac16+\frac16=\frac12$),
and each of the other seven contains a factor $\YYbar$ and has Frobenius
norm at most $\norm{\YYbar}_F\max\{1,\norm{\YYtil}\}^{2}\le\norm{\YYbar}_F$,
by $\norm{\YYbar}\le1$ and $\norm{\YYtil}\le\frac12$.
Hence, by $(7x+y)^{2}\le98x^{2}+2y^{2}$,
\begin{equation}\label{eq:resbound}
\norm{(\YYbar+\YYtil)^{3}}_F^{2}
\le98\norm{\YYbar}_F^{2}+2\norm{\YYtil}_{S_6}^{6}
\le98\cdot\frac{32n}{\dheavy}+3.4\cdot10^{4}\,\frac{n}{\log^{3}n}
\le10^{5}\Bigl(\frac n{\dheavy}+\frac n{\log^{3}n}\Bigr).
\end{equation}

\emph{The three statistics.}
Given the references and the operators, each $\xxrand_j$ is a Gaussian
vector with covariance $\ZZ(\LLtil)\LLtil\ZZ(\LLtil)$, and the three
statistics of Step 3 of \Cref{fig:attempt} are the averages over
$j\le10^{3}\log n$ of the independent quadratic forms
$\gg_j^{\top}\NN_1\gg_j$, $\gg_j^{\top}\NN_2\gg_j$, $\gg_j^{\top}\NN_3\gg_j$
with
\begin{equation}\label{eq:forms}
\begin{aligned}
\NN_1&\defeq\WWtil^{1/2}\BBtil\ZZ(\LLtil)\,\LLtil\,\ZZ(\LLtil)\BBtil^{\top}\WWtil^{1/2},\\
\NN_2&\defeq\WWtil^{1/2}\BBtil\ZZ(\LLtil)\bigl((\II-\ZZ(\MMtil)\LLtil)^{3}\bigr)^{\top}\LLtil(\II-\ZZ(\MMtil)\LLtil)^{3}\ZZ(\LLtil)\BBtil^{\top}\WWtil^{1/2},\\
\NN_3&\defeq\WWtil^{1/2}\BBtil\ZZ(\LLtil)\bigl(\II-(\II-\ZZ(\MMtil)\LLtil)^{3}\bigr)^{\top}\DD(\Gl)\bigl(\II-(\II-\ZZ(\MMtil)\LLtil)^{3}\bigr)\ZZ(\LLtil)\BBtil^{\top}\WWtil^{1/2}.
\end{aligned}
\end{equation}
Suppose the references and the operators are accurate.
By \Cref{eq:ops}, $\LLtil^{1/2}\ZZ(\LLtil)\LLtil^{1/2}$ lies between
$\frac{39}{40}\PPi$ and $\frac{41}{40}\PPi$, so
\begin{equation}\label{eq:cov}
\ZZ(\LLtil)\LLtil\ZZ(\LLtil)
=\LLtil^{\dag/2}\bigl(\LLtil^{1/2}\ZZ(\LLtil)\LLtil^{1/2}\bigr)^{2}\LLtil^{\dag/2}
\in\Bigl[\Bigl(\frac{39}{40}\Bigr)^{2},\Bigl(\frac{41}{40}\Bigr)^{2}\Bigr]\cdot\LLtil^{\dag}
\subseteq[0.94,\,1.07]\cdot\LL(G)^{\dag},
\end{equation}
using $\LLtil^{\dag}\in[1/1.001,1/0.999]\cdot\LL(G)^{\dag}$.
The nonzero eigenvalues of $\NN_1$ are those of
\[
\LLtil^{1/2}\ZZ(\LLtil)\LLtil\ZZ(\LLtil)\LLtil^{1/2}=(\LLtil^{1/2}\ZZ(\LLtil)\LLtil^{1/2})^{2},
\]
so
\[
\tr[\NN_1]\le\Bigl(\frac{41}{40}\Bigr)^{2}(n-1)\le1.06n,
\qquad
\norm{\NN_1}\le\Bigl(\frac{41}{40}\Bigr)^{2}\le1.07 .
\]
For $\NN_2$, by \Cref{eq:cov}, $\LLtil\preceq1.001\LL(G)$ and
\Cref{eq:resenergy},
\begin{align*}
\tr[\NN_2]
&=\tr\bigl[\bigl((\II-\ZZ(\MMtil)\LLtil)^{3}\bigr)^{\top}\LLtil(\II-\ZZ(\MMtil)\LLtil)^{3}\,\ZZ(\LLtil)\LLtil\ZZ(\LLtil)\bigr]\\
&\le1.07\cdot1.001\,\tr\bigl[\bigl((\II-\ZZ(\MMtil)\LLtil)^{3}\bigr)^{\top}\LL(G)(\II-\ZZ(\MMtil)\LLtil)^{3}\LL(G)^{\dag}\bigr]\\
&=1.08\,\norm{\LL(G)^{1/2}(\II-\ZZ(\MMtil)\LLtil)^{3}\LL(G)^{\dag/2}}_F^{2}
=1.08\,\norm{(\YYbar+\YYtil)^{3}}_F^{2},
\end{align*}
and likewise, the nonzero eigenvalues of $\NN_2$ being those of
$\LLtil^{1/2}(\II-\ZZ(\MMtil)\LLtil)^{3}\,\ZZ(\LLtil)\LLtil\ZZ(\LLtil)\,((\II-\ZZ(\MMtil)\LLtil)^{3})^{\top}\LLtil^{1/2}$,
\[
\norm{\NN_2}\le1.07\cdot1.001\,\norm{\LL(G)^{1/2}(\II-\ZZ(\MMtil)\LLtil)^{3}\LL(G)^{\dag/2}}^{2}
=1.08\,\norm{(\YYbar+\YYtil)^{3}}^{2}\le1.08\cdot(1+\tfrac12)^{6}\le13 .
\]
For $\NN_3$ let $\Gl\neq\emptyset$ (otherwise $\DD(\Gl)=0$ and $\NN_3=0$),
so that $\MM$ and $\MMtil$ are positive definite.
Expanding the cube,
\[
\II-(\II-\ZZ(\MMtil)\LLtil)^{3}
=\ZZ(\MMtil)\bigl(3\II-3\LLtil\ZZ(\MMtil)+\LLtil\ZZ(\MMtil)\LLtil\ZZ(\MMtil)\bigr)\LLtil .
\]
By \Cref{eq:ops}, $\norm{\MMtil^{1/2}\ZZ(\MMtil)\MMtil^{1/2}}\le1+1/\log n\le1.05$;
by $\LLtil\preceq1.001\LL(G)\preceq2.002\MM\preceq2.01\MMtil$,
$\norm{\MMtil^{-1/2}\LLtil\MMtil^{-1/2}}\le2.01$; hence
\begin{align*}
&\norm{\MMtil^{-1/2}\bigl(3\II-3\LLtil\ZZ(\MMtil)+\LLtil\ZZ(\MMtil)\LLtil\ZZ(\MMtil)\bigr)\MMtil^{1/2}}\\
&\quad=\norm{3\II-3(\MMtil^{-1/2}\LLtil\MMtil^{-1/2})(\MMtil^{1/2}\ZZ(\MMtil)\MMtil^{1/2})
+\bigl((\MMtil^{-1/2}\LLtil\MMtil^{-1/2})(\MMtil^{1/2}\ZZ(\MMtil)\MMtil^{1/2})\bigr)^{2}}\\
&\quad\le3+3\cdot2.01\cdot1.05+(2.01\cdot1.05)^{2}\le15.2 .
\end{align*}
Since $\LL(G)^{\dag}\preceq1.001\LLtil^{\dag}$, we have
$\LLtil\LL(G)^{\dag}\LLtil\preceq1.001\LLtil\preceq2.02\MMtil$, and therefore
\begin{align*}
&\bigl(\II-(\II-\ZZ(\MMtil)\LLtil)^{3}\bigr)\LL(G)^{\dag}\bigl(\II-(\II-\ZZ(\MMtil)\LLtil)^{3}\bigr)^{\top}\\
&\quad\preceq2.02\,\ZZ(\MMtil)\bigl(3\II-3\LLtil\ZZ(\MMtil)+\LLtil\ZZ(\MMtil)\LLtil\ZZ(\MMtil)\bigr)\MMtil\bigl(\cdots\bigr)^{\top}\ZZ(\MMtil)\\
&\quad\preceq2.02\cdot15.2^{2}\,\ZZ(\MMtil)\MMtil\ZZ(\MMtil)
\preceq2.02\cdot15.2^{2}\cdot1.05^{2}\,\MMtil^{-1}
\preceq\frac{2.02\cdot15.2^{2}\cdot1.05^{2}}{0.999}\,\MM^{-1}
\preceq520\,\MM^{-1},
\end{align*}
where the second step writes
$\ZZ(\MMtil)(\cdots)\MMtil^{1/2}=\MMtil^{-1/2}(\MMtil^{1/2}\ZZ(\MMtil)\MMtil^{1/2})(\MMtil^{-1/2}(\cdots)\MMtil^{1/2})$
and bounds the last factor by its norm.
Consequently, by \Cref{eq:cov}, \Cref{lem:topk}(f) and
$\DD(\Gl)\preceq\MM$ (\Cref{lem:topk}(c)),
\begin{align*}
\tr[\NN_3]
&\le1.07\,\tr\bigl[\DD(\Gl)\bigl(\II-(\II-\ZZ(\MMtil)\LLtil)^{3}\bigr)\LL(G)^{\dag}\bigl(\II-(\II-\ZZ(\MMtil)\LLtil)^{3}\bigr)^{\top}\bigr]\\
&\le1.07\cdot520\,\tr[\DD(\Gl)\MM^{-1}]\le560n,\\
\norm{\NN_3}
&\le1.07\,\norm{\DD(\Gl)^{1/2}\bigl(\II-(\II-\ZZ(\MMtil)\LLtil)^{3}\bigr)\LL(G)^{\dag}\bigl(\II-(\II-\ZZ(\MMtil)\LLtil)^{3}\bigr)^{\top}\DD(\Gl)^{1/2}}\\
&\le1.07\cdot520\,\norm{\DD(\Gl)^{1/2}\MM^{-1}\DD(\Gl)^{1/2}}\le560 .
\end{align*}

\emph{The tests.}
Fix $\NN\in\{\NN_1,\NN_2,\NN_3\}$.
The stacked vector $(\gg_1,\dots,\gg_{10^{3}\log n})$ has independent standard
Gaussian entries, and $\sum_j\gg_j^{\top}\NN\gg_j$ is its quadratic form with
the block diagonal matrix with $10^{3}\log n$ diagonal blocks $\NN$, of
trace $10^{3}\log n\cdot\tr[\NN]$, operator norm $\norm{\NN}$ and trace of
the square $10^{3}\log n\cdot\tr[\NN^{2}]\le10^{3}\log n\cdot\norm{\NN}\tr[\NN]$.
\Cref{lem:quadtail} with $t=15\ln n$, divided by $10^{3}\log n\ge10^{3}\ln n$,
gives with probability at least $1-n^{-15}$
\begin{equation}\label{eq:tail}
\frac1{10^{3}\log n}\sum_j\gg_j^{\top}\NN\gg_j
\le\tr[\NN]+2\sqrt{0.015\,\norm{\NN}\tr[\NN]}+0.03\norm{\NN} .
\end{equation}
\begin{itemize}[leftmargin=*]
\item For $\NN_1$ the right-hand side is at most
$1.06n+2\sqrt{0.015\cdot1.07\cdot1.06n}+0.04\le1.06n+0.27\sqrt n+0.04\le1.1n$,
as $n\ge2^{20}$.
\item For $\NN_2$, if moreover \Cref{eq:goodrefs} fails,
$\tr[\NN_2]\le1.08\cdot10^{5}(n/\dheavy+n/\log^{3}n)$ by \Cref{eq:resbound},
and $2\sqrt{0.015\cdot13\,\tr[\NN_2]}\le0.04\tr[\NN_2]+4.9$ by
$2\sqrt{xy}\le x+y$, so the right-hand side is at most
$1.04\cdot1.08\cdot10^{5}(n/\dheavy+n/\log^{3}n)+5.3\le10^{8}(n/\dheavy+n/\log^{3}n)$,
as $n/\log^{3}n\ge1$.
\item For $\NN_3$ the right-hand side is at most
$560n+2\sqrt{0.015\cdot560\cdot560n}+17\le560n+137\sqrt n+17\le10^{3}n$.
\end{itemize}
So if the references and the operators are accurate and \Cref{eq:goodrefs}
fails, the three tests pass simultaneously with probability at least
$1-3n^{-15}$, which is the second sentence of (a).
If the references are accurate then $\LLtil\succeq0.999\LL(G)$, so on
acceptance the first two energies of \Cref{eq:energies} are at most
$1.1n/0.999\le1.2n$ and $10^{8}(n/\dheavy+n/\log^{3}n)/0.999\le2\cdot10^{8}(n/\dheavy+n/\log^{3}n)$,
and the third test is the third energy; this is the third sentence of (a).
The last sentence of (a) is the union bound over the failure of the
references to be accurate, the failure of the operators to be accurate,
the event that the references are accurate and \Cref{eq:goodrefs} holds
(probability at most $10^{-2}$), and the failure of the tests
(probability at most $3n^{-15}$).

\emph{All pairs.}
Suppose the references and the operators are accurate and fix $u\neq v$.
The variables $(\xxrand_j)_u-(\xxrand_j)_v=(\oone_u-\oone_v)^{\top}\ZZ(\LLtil)\BBtil^{\top}\WWtil^{1/2}\gg_j$,
$j\le10^{3}\log n$, are independent centered Gaussians with the common
variance
$(\oone_u-\oone_v)^{\top}\ZZ(\LLtil)\LLtil\ZZ(\LLtil)(\oone_u-\oone_v)$, which
lies in $[0.94,1.06]\cdot(\oone_u-\oone_v)^{\top}\LL(G)^{\dag}(\oone_u-\oone_v)$
by \Cref{eq:cov} and $(\oone_u-\oone_v)\perp\oone$.
Divided by this variance, $\sum_j((\xxrand_j)_u-(\xxrand_j)_v)^{2}$ has the
law of $\gg^{\top}\gg$ for a standard Gaussian vector
$\gg\in\R^{10^{3}\log n}$.
\Cref{lem:quadtail} with $\NN=\II$ and $t=10^{3}\log n/80$ gives
\[
\Pr\Bigl[\Bigl|\frac{\gg^{\top}\gg}{10^{3}\log n}-1\Bigr|>\frac14\Bigr]
\le2\exp\Bigl(-\frac{10^{3}\log n}{80}\Bigr)
\le2n^{-18}\le n^{-16},
\]
because $2/\sqrt{80}+2/80\le\frac14$ and $e^{12.5\log n}=n^{12.5/\ln2}\ge n^{18}$.
A union bound over the fewer than $n^{2}$ pairs gives, with probability
$1-n^{-14}$, for all $u,v$ an average in
$[\frac34\cdot0.94,\frac54\cdot1.06]\cdot(\oone_u-\oone_v)^{\top}\LL(G)^{\dag}(\oone_u-\oone_v)\subseteq[\frac12,2]\cdot(\oone_u-\oone_v)^{\top}\LL(G)^{\dag}(\oone_u-\oone_v)$,
which is (b).

\emph{Cost.}
$E(\LLtil)$ has $O(n\log n)$ edges.
Step 1 of \Cref{fig:attempt} draws $10^{3}\log n\cdot|E(\LLtil)|=O(n\log^{2}n)$
Gaussians.
For each $j$, Step 2(a) forms $\BBtil^{\top}\WWtil^{1/2}\gg_j$ in
$O(|E(\LLtil)|)$ operations and applies $\ZZ(\LLtil)$ once, in
$\Odouble(n\log n\log40)=\Odouble(n\log n)$ time by \Cref{lem:chain}; Step 2(b)
applies $\LLtil$ three times, $O(n\log n)$ each, and $\ZZ(\MMtil)$ three
times, $\Odouble(n\log n\log\log n)=\Odouble(n\log n)$ each by \Cref{lem:chain}
(the factor $\log\log n$ is hidden by $\Odouble$); Step 2(c) takes $O(n)$.
Step 3 evaluates two quadratic forms of $\LLtil$ per vector, $O(n\log n)$
each, and one of $\DD(\Gl)$, $O(n)$ given $\dd(\Gl)$.
Summed over the $10^{3}\log n$ vectors this is $\Odouble(n\log^{2}n)$, which
is (c).
\end{proof}

\section{References from intensities}\label{sec:references}

The references $\LLtil$ and $\MMtil$ of a refresh are drawn by retaining
every sampled edge independently, with a probability of the form
$1-e^{-\varrho_e}$ for an intensity $\varrho_e>0$.
The centered retention variable $\zeta_e$ of a sampled edge $e$ is that of
\Cref{eq:references}: $\zeta_e=1/(1-e^{-\varrho_e})-1$ if $e$ is retained and
$\zeta_e=-1$ otherwise, so that $\Ex[\zeta_e]=0$, $\zeta_e\ge-1$, and the
weight of $e$ in the reference is $\ww_e(1+\zeta_e)$.
The accuracy of the references is a matrix concentration bound, the master
tail bound of Tropp, and the per-edge moments of the retained weights are
what the sixth-moment estimate of the residual stream
(the proof of \Cref{lem:refresh}(a)) uses about the references; both are collected here.

\begin{lemma}[references from intensities]\label{lem:bernoulli}
Let $G=(V,E,\ww)$ be connected with the split $\Gh,\Gl$ and $\MM$ of
\Cref{def:heavy}, and let $E=E_{\mathrm{s}}\cup E_{\mathrm{x}}$ be a partition
of its edges.
Let intensities $\varrho_e>0$, $e\in E_{\mathrm{s}}$, satisfy
\[
\varrho_e\ \ge\ 10^{8}\log n\cdot\ww_e\Reff(G)_e
\quad\text{for every }e\in E_{\mathrm{s}},
\qquad
\sum_{e\in E_{\mathrm{s}}}\varrho_e\ \le\ 10^{17}n\log n .
\]
Retain every $e\in E_{\mathrm{s}}$ independently with probability
$1-e^{-\varrho_e}$, with weight $\ww_e/(1-e^{-\varrho_e})$, and keep every
$e\in E_{\mathrm{x}}$ with its weight; let $\LLtil$ be the Laplacian of the
result and $\MMtil\defeq\DD(\Gl)+\LLtil(\Gh)$, where $\LLtil(\Gh)$ is the
part of $\LLtil$ on the heavy edges.
Then:
\begin{enumerate}[label=(\alph*),leftmargin=*]
\item with probability $1-2n^{-14}$,
$\LLtil\approx_{10^{-3}}\LL(G)$ and $\MMtil\approx_{10^{-3}}\MM$;
\item with probability $1-n^{-14}$, the retained edges number at most
$2\cdot10^{17}n\log n+|E_{\mathrm{x}}|$;
\item (per-edge moments) for every $e\in E_{\mathrm{s}}$ and every integer
$p$ with $2\le p\le6$,
\begin{align*}
\Ex\bigl[|\zeta_e|^{p}\bigr]
\bigl(\ww_e\cchi_e^{\top}\LL(G)^{\dag}\cchi_e\bigr)^{p-1}
&\le\Bigl(\frac{10^{-7}}{\log n}\Bigr)^{p-1},\\
\Ex\bigl[|\zeta_e|^{p}\bigr]
\bigl(\ww_e\cchi_e^{\top}\MM^{\dag}\cchi_e\bigr)^{p-1}
&\le\Bigl(\frac{2\cdot10^{-7}}{\log n}\Bigr)^{p-1}.
\end{align*}
\end{enumerate}
\end{lemma}

\begin{proof}
Throughout, $\log n\ge\ln n$ because logarithms are to the base $2$.
The retained weight of a sampled edge is $\ww_e(1+\zeta_e)$ and the exactly
kept edges have their weights on both sides, so
\begin{equation}\label{eq:bernoulli-centered}
\LLtil-\LL(G)=\sum_{e\in E_{\mathrm{s}}}\zeta_e\ww_e\cchi_e\cchi_e^{\top},
\qquad
\MMtil-\MM=\LLtil(\Gh)-\LL(\Gh)
=\sum_{e\in E_{\mathrm{s}}\cap E(\Gh)}\zeta_e\ww_e\cchi_e\cchi_e^{\top}.
\end{equation}
The matrix concentration tool is the master tail bound of Tropp; it is
stated here because \Cref{clm:bernoulli-generic} applies it to summands
whose norms are not uniformly bounded, which rules out the matrix Chernoff
bound with a uniform norm bound.

\begin{claim}[Theorem~3.6 of~\cite{T12:journal}]\label{clm:master}
Consider a finite sequence $\{\XX_k\}$ of independent, random,
self-adjoint matrices.
For all $t\in\R$,
\[
\Pr\Bigl[\lambda_{\max}\Bigl(\sum_k\XX_k\Bigr)\ge t\Bigr]
\le\inf_{\theta>0}\ e^{-\theta t}\cdot
\tr\Bigl[\exp\Bigl(\sum_k\log\Ex\bigl[e^{\theta\XX_k}\bigr]\Bigr)\Bigr].
\]
\end{claim}

The scalar input to \Cref{clm:master} is the moment generating function of
a centered retention variable.

\begin{claim}\label{clm:mgf}
Let $\zeta$ be the centered retention variable of intensity $\varrho>0$.
Then for every real $x$ with $|x|\le10^{-3}\varrho$,
\[
\log\Ex\bigl[e^{x\zeta}\bigr]\ \le\ \frac{0.51\,x^{2}}{\varrho}.
\]
\end{claim}

\begin{proof}[Proof of \Cref{clm:mgf}]
The function $g(x)\defeq\Ex[e^{x\zeta}]
=e^{-\varrho}e^{-x}+(1-e^{-\varrho})e^{x e^{-\varrho}/(1-e^{-\varrho})}$
has $g(0)=1$, $g'(0)=\Ex[\zeta]=0$ and
\[
g''(y)=e^{-\varrho}e^{-y}
+\frac{e^{-2\varrho}}{1-e^{-\varrho}}\,e^{y e^{-\varrho}/(1-e^{-\varrho})}.
\]
We show $g''(y)\le1.002/\varrho$ for every $|y|\le10^{-3}\varrho$.
Since $e^{-\varrho}/(1-e^{-\varrho})=1/(e^{\varrho}-1)\le1/\varrho$, the
exponent of the first term is at most $|y|\le10^{-3}\varrho$ and the
exponent of the second term is at most $|y|/\varrho\le10^{-3}$.
\begin{itemize}[leftmargin=*]
\item If $\varrho\le1$, both exponents are at most $10^{-3}$, so
\[
g''(y)\le\Bigl(e^{-\varrho}+\frac{e^{-2\varrho}}{1-e^{-\varrho}}\Bigr)
e^{10^{-3}}
=\frac{e^{-\varrho}}{1-e^{-\varrho}}\,e^{10^{-3}}
\le\frac{1.0011}{\varrho}.
\]
\item If $\varrho>1$, the first term is at most
$e^{-\varrho}e^{10^{-3}\varrho}=e^{-0.999\varrho}$ and the second is at most
$1.0011\,e^{-2\varrho}/(1-e^{-\varrho})$.
The function $\varrho e^{-0.999\varrho}$ is maximized at $\varrho=1/0.999$
with value at most $0.37$, and $\varrho e^{-2\varrho}/(1-e^{-\varrho})$ is
decreasing on $\varrho>1$ with value at most $0.22$ at $\varrho=1$.
Hence $g''(y)\le(0.37+1.0011\cdot0.22)/\varrho\le1.002/\varrho$.
\end{itemize}
By Taylor's theorem, $g(x)=1+\frac12g''(y)x^{2}$ for some $y$ between $0$
and $x$, so $g(x)\le1+0.501x^{2}/\varrho$, and
$\log g(x)\le g(x)-1\le0.51x^{2}/\varrho$.
\end{proof}

Both references are handled by one computation.

\begin{claim}\label{clm:bernoulli-generic}
Let $\NN$ be a symmetric positive semidefinite $n\times n$ matrix and
$F\subseteq E_{\mathrm{s}}$ a set of sampled edges such that
\[
\sum_{e\in F}\ww_e\NN^{\dag/2}\cchi_e\cchi_e^{\top}\NN^{\dag/2}\preceq\II
\qquad\text{and}\qquad
\ww_e\cchi_e^{\top}\NN^{\dag}\cchi_e\le\frac{2\cdot10^{-8}\varrho_e}{\log n}
\quad\text{for every }e\in F .
\]
Then with probability at least $1-n^{-14}$,
\[
\Bigl\lVert\sum_{e\in F}\zeta_e\ww_e\NN^{\dag/2}\cchi_e\cchi_e^{\top}
\NN^{\dag/2}\Bigr\rVert\le10^{-3}.
\]
\end{claim}

\begin{proof}[Proof of \Cref{clm:bernoulli-generic}]
The summands $\zeta_e\ww_e\NN^{\dag/2}\cchi_e\cchi_e^{\top}\NN^{\dag/2}$,
$e\in F$, are independent symmetric random matrices with mean zero.
Put $\theta\defeq5\cdot10^{4}\log n$ in \Cref{clm:master}.
For a fixed $e\in F$, the matrix
$\theta\zeta_e\ww_e\NN^{\dag/2}\cchi_e\cchi_e^{\top}\NN^{\dag/2}$ has rank
one with the single nonzero eigenvalue
$\theta\zeta_e\ww_e\cchi_e^{\top}\NN^{\dag}\cchi_e$, so its exponential is
$\II$ plus
$\bigl(e^{\theta\zeta_e\ww_e\cchi_e^{\top}\NN^{\dag}\cchi_e}-1\bigr)$ times
the projection
$\NN^{\dag/2}\cchi_e\cchi_e^{\top}\NN^{\dag/2}
/(\cchi_e^{\top}\NN^{\dag}\cchi_e)$,
and taking the expectation and then the matrix logarithm gives
\[
\log\Ex\bigl[e^{\pm\theta\zeta_e\ww_e\NN^{\dag/2}\cchi_e\cchi_e^{\top}
\NN^{\dag/2}}\bigr]
=\log\Ex\bigl[e^{\pm\theta\zeta_e\ww_e\cchi_e^{\top}\NN^{\dag}\cchi_e}\bigr]
\cdot\frac{\NN^{\dag/2}\cchi_e\cchi_e^{\top}\NN^{\dag/2}}
{\cchi_e^{\top}\NN^{\dag}\cchi_e}.
\]
The hypothesis gives
$\theta\ww_e\cchi_e^{\top}\NN^{\dag}\cchi_e
\le5\cdot10^{4}\log n\cdot2\cdot10^{-8}\varrho_e/\log n=10^{-3}\varrho_e$,
so \Cref{clm:mgf} applies with $x=\pm\theta\ww_e\cchi_e^{\top}\NN^{\dag}\cchi_e$
and yields
\begin{align*}
\log\Ex\bigl[e^{\pm\theta\zeta_e\ww_e\NN^{\dag/2}\cchi_e\cchi_e^{\top}
\NN^{\dag/2}}\bigr]
&\preceq\frac{0.51\,\theta^{2}\ww_e\cchi_e^{\top}\NN^{\dag}\cchi_e}{\varrho_e}
\cdot\ww_e\NN^{\dag/2}\cchi_e\cchi_e^{\top}\NN^{\dag/2}\\
&\preceq0.51\cdot25\cdot10^{8}\log^{2}n\cdot\frac{2\cdot10^{-8}}{\log n}
\cdot\ww_e\NN^{\dag/2}\cchi_e\cchi_e^{\top}\NN^{\dag/2}\\
&=25.5\log n\cdot\ww_e\NN^{\dag/2}\cchi_e\cchi_e^{\top}\NN^{\dag/2}.
\end{align*}
Summing over $e\in F$ and using the first hypothesis, the sum of the
logarithms is at most $25.5\log n\cdot\II$, whose exponential has trace at
most $n\,e^{25.5\log n}$.
\Cref{clm:master} with $t=10^{-3}$, for either sign, gives
\[
\Pr\Bigl[\lambda_{\max}\Bigl(\pm\sum_{e\in F}\zeta_e\ww_e\NN^{\dag/2}
\cchi_e\cchi_e^{\top}\NN^{\dag/2}\Bigr)\ge10^{-3}\Bigr]
\le e^{-50\log n}\cdot n\,e^{25.5\log n}
=n\,e^{-24.5\log n}\le n^{-23.5},
\]
using $\log n\ge\ln n$ in the last step.
The two signs together fail with probability at most
$2n^{-23.5}\le n^{-14}$.
\end{proof}

\emph{Part (a), the Laplacian.}
Apply \Cref{clm:bernoulli-generic} with $\NN=\LL(G)$ and $F=E_{\mathrm{s}}$.
The first hypothesis holds because
$\sum_{e\in E_{\mathrm{s}}}\ww_e\LL(G)^{\dag/2}\cchi_e\cchi_e^{\top}
\LL(G)^{\dag/2}\preceq\LL(G)^{\dag/2}\LL(G)\LL(G)^{\dag/2}\preceq\II$, and the
second
because $\ww_e\cchi_e^{\top}\LL(G)^{\dag}\cchi_e=\ww_e\Reff(G)_e
\le10^{-8}\varrho_e/\log n$.
So with probability $1-n^{-14}$,
$\norm{\LL(G)^{\dag/2}(\LLtil-\LL(G))\LL(G)^{\dag/2}}\le10^{-3}$ by
\Cref{eq:bernoulli-centered}.
Every $\cchi_e$ is orthogonal to $\oone$, and $G$ is connected, so the
range of $\LLtil-\LL(G)$ lies in the range of $\LL(G)$; hence for every
$\xx\in\R^{V}$,
$\xx^{\top}(\LLtil-\LL(G))\xx
=(\LL(G)^{1/2}\xx)^{\top}\LL(G)^{\dag/2}(\LLtil-\LL(G))\LL(G)^{\dag/2}
(\LL(G)^{1/2}\xx)$, whose absolute value is at most
$10^{-3}\xx^{\top}\LL(G)\xx$.
This is $\LLtil\approx_{10^{-3}}\LL(G)$.

\emph{Part (a), the heavy reference.}
First, $\LL(G)\preceq2\MM$: writing $\LL(\Gl)=\DD(\Gl)-\AA(\Gl)$,
\begin{align*}
\DD(\Gl)+\AA(\Gl)
&=\sum_{e=uv\in E(\Gl)}\ww_e(\oone_u+\oone_v)(\oone_u+\oone_v)^{\top}
\succeq0,\\
\LL(G)=\LL(\Gh)+\LL(\Gl)
&\preceq\LL(\Gh)+2\DD(\Gl)\preceq2\MM .
\end{align*}
In particular the kernel of $\MM$ is contained in the kernel of $\LL(G)$,
which is the span of $\oone$, so every $\cchi_e$ lies in the range of
$\MM$, and for every $e$, by the variational characterization of the
pseudoinverse on the range,
\[
\cchi_e^{\top}\MM^{\dag}\cchi_e
=\max_{\xx\in\R^{V}}\bigl(2\cchi_e^{\top}\xx-\xx^{\top}\MM\xx\bigr)
\le\max_{\xx\in\R^{V}}\bigl(2\cchi_e^{\top}\xx-\tfrac12\xx^{\top}\LL(G)\xx\bigr)
=2\cchi_e^{\top}\LL(G)^{\dag}\cchi_e
=2\Reff(G)_e .
\]
Apply \Cref{clm:bernoulli-generic} with $\NN=\MM$ and
$F=E_{\mathrm{s}}\cap E(\Gh)$.
The first hypothesis holds because $\DD(\Gl)\succeq0$ gives
$\sum_{e\in E(\Gh)}\ww_e\MM^{\dag/2}\cchi_e\cchi_e^{\top}\MM^{\dag/2}
=\MM^{\dag/2}\LL(\Gh)\MM^{\dag/2}\preceq\MM^{\dag/2}\MM\MM^{\dag/2}\preceq\II$,
and the second because
$\ww_e\cchi_e^{\top}\MM^{\dag}\cchi_e\le2\ww_e\Reff(G)_e
\le2\cdot10^{-8}\varrho_e/\log n$.
So with probability $1-n^{-14}$,
$\norm{\MM^{\dag/2}(\MMtil-\MM)\MM^{\dag/2}}\le10^{-3}$ by
\Cref{eq:bernoulli-centered}, and as the range of $\MMtil-\MM$ lies in the
span of the $\cchi_e$ and hence in the range of $\MM$, the same
computation as for the Laplacian gives $\MMtil\approx_{10^{-3}}\MM$.
The two events hold together with probability at least $1-2n^{-14}$.

\emph{Part (b), the number of retained edges.}
The number of retained sampled edges is
$\sum_{e\in E_{\mathrm{s}}}(1+\zeta_e)(1-e^{-\varrho_e})$, a sum of
independent indicators with
\begin{align*}
\Ex\bigl[e^{\theta(1+\zeta_e)(1-e^{-\varrho_e})}\bigr]
&=1+(1-e^{-\varrho_e})(e^{\theta}-1)\\
&\le\exp\bigl((e^{\theta}-1)(1-e^{-\varrho_e})\bigr)
\le\exp\bigl((e^{\theta}-1)\varrho_e\bigr)
\quad\text{for every }\theta>0 .
\end{align*}
With $\theta=\ln2$ and Markov's inequality applied to the exponential of
the sum, the sum exceeds $2\cdot10^{17}n\log n$ with probability at most
\[
\exp\Bigl(-2\ln2\cdot10^{17}n\log n+\sum_{e\in E_{\mathrm{s}}}\varrho_e\Bigr)
\le\exp\bigl(-(2\ln2-1)\cdot10^{17}n\log n\bigr)
\le\exp\bigl(-3\cdot10^{16}n\log n\bigr)\le n^{-14}.
\]
Adding the $|E_{\mathrm{x}}|$ exactly kept edges gives the count.
The same computation with any bound $R\ge\sum_{e\in E_{\mathrm{s}}}\varrho_e$
in place of $10^{17}n\log n$ gives at most $2R+|E_{\mathrm{x}}|$ retained
edges with probability $1-e^{-0.38R}$.

\emph{Part (c), the per-edge moments.}
Fix $e\in E_{\mathrm{s}}$ and $2\le p\le6$.
From the distribution of $\zeta_e$,
\[
\Ex\bigl[|\zeta_e|^{p}\bigr]
=e^{-\varrho_e}\cdot1
+(1-e^{-\varrho_e})\Bigl(\frac{e^{-\varrho_e}}{1-e^{-\varrho_e}}\Bigr)^{p}
=e^{-\varrho_e}+\frac{e^{-p\varrho_e}}{(1-e^{-\varrho_e})^{p-1}} .
\]
Let $\lambda\in\{1,2\}$ and let $\NN\in\{\LL(G),\MM\}$ be the matrix with
$\ww_e\cchi_e^{\top}\NN^{\dag}\cchi_e\le\lambda\cdot10^{-8}\varrho_e/\log n$
($\lambda=1$ for $\LL(G)$ by the hypothesis on $\varrho_e$, $\lambda=2$ for
$\MM$ by the inequality $\cchi_e^{\top}\MM^{\dag}\cchi_e\le2\Reff(G)_e$
proved in part (a)).
We show
$\Ex[|\zeta_e|^{p}](\ww_e\cchi_e^{\top}\NN^{\dag}\cchi_e)^{p-1}
\le(\lambda\cdot10^{-7}/\log n)^{p-1}$.
\begin{itemize}[leftmargin=*]
\item If $\varrho_e<1$: the function $(1-e^{-\varrho})/\varrho$ is
decreasing, so $1-e^{-\varrho_e}\ge(1-e^{-1})\varrho_e\ge0.63\varrho_e$,
and
$\Ex[|\zeta_e|^{p}]\le1+(1-e^{-\varrho_e})^{1-p}\le2(1-e^{-\varrho_e})^{1-p}$.
Hence, with
$\ww_e\cchi_e^{\top}\NN^{\dag}\cchi_e
\le\lambda\cdot10^{-8}(1-e^{-\varrho_e})/(0.63\log n)
\le1.6\lambda\cdot10^{-8}(1-e^{-\varrho_e})/\log n$,
\[
\Ex\bigl[|\zeta_e|^{p}\bigr]
\bigl(\ww_e\cchi_e^{\top}\NN^{\dag}\cchi_e\bigr)^{p-1}
\le2\Bigl(\frac{1.6\lambda\cdot10^{-8}}{\log n}\Bigr)^{p-1}
=2\cdot0.16^{p-1}\Bigl(\frac{\lambda\cdot10^{-7}}{\log n}\Bigr)^{p-1}
\le\Bigl(\frac{\lambda\cdot10^{-7}}{\log n}\Bigr)^{p-1}.
\]
\item If $\varrho_e\ge1$: $1-e^{-\varrho_e}\ge1-e^{-1}$, so
$\Ex[|\zeta_e|^{p}]\le e^{-\varrho_e}\bigl(1+(e^{-1}/(1-e^{-1}))^{p-1}\bigr)
\le e^{-\varrho_e}(1+0.59^{p-1})\le2e^{-\varrho_e}$.
The function $\varrho^{p-1}e^{-\varrho}$ is maximized at $\varrho=p-1$, so
\begin{align*}
\Ex\bigl[|\zeta_e|^{p}\bigr]
\bigl(\ww_e\cchi_e^{\top}\NN^{\dag}\cchi_e\bigr)^{p-1}
&\le2\varrho_e^{p-1}e^{-\varrho_e}
\Bigl(\frac{\lambda\cdot10^{-8}}{\log n}\Bigr)^{p-1}\\
&\le2\Bigl(\frac{p-1}{10e}\Bigr)^{p-1}
\Bigl(\frac{\lambda\cdot10^{-7}}{\log n}\Bigr)^{p-1}
\le\Bigl(\frac{\lambda\cdot10^{-7}}{\log n}\Bigr)^{p-1},
\end{align*}
because $(p-1)/(10e)\le5/(10e)<\frac12$ for $p\le6$.
\end{itemize}
This is part (c) with $\lambda=1$ for $\LL(G)$ and $\lambda=2$ for $\MM$.
\end{proof}

\section{The maintained state}\label{sec:state}

The structure keeps three kinds of information: the current graph in
per-vertex trees, the retained part of the current block in a
$\textsc{QueryDrivenPartition}$, and the \emph{record} of the previous
refresh, from which the intensities of the current block are read.
\Cref{fig:state} lists them; the record is defined in \Cref{sec:refresh}.

\begin{figure}[H]\centering
\framebox{\begin{minipage}{0.96\linewidth}\small
\textbf{State} of the structure during block $k$, with $\dheavy=\log^{3}n$.
\begin{itemize}[leftmargin=*,itemsep=1pt]
\item For every vertex $v$: a balanced tree $T_v$ of the current edges at
$v$ keyed by current weight; the $\dheavy$ largest keys are the
\emph{current heavy} edges at $v$, every other edge at $v$ is
\emph{current light}.
An edge is current heavy if it is current heavy at either endpoint.
\item For every vertex $v$: a balanced tree $T^{\mathrm{light}}_v$ of the
current light edges at $v$ that are edges of $\Gret_k$, augmented with
subtree sums of the \emph{active weight}: the weight of the edge in
$\Gret_k$ if the edge has been light at $v$ since the previous refresh,
and $0$ otherwise.
\item The list $\mathrm{Changed}$ of the edges whose heavy/light status at
some endpoint changed since the previous refresh, and the list $\Gnew_k$
of the held new edges and weight increases.
\item $P$, a $\textsc{QueryDrivenPartition}$ (\Cref{lem:partition}) holding
exactly $\Gret_k$ with its weights.
\item The record of refresh $k-1$ (\Cref{def:record}): the sparse reference
$\LLtil_{k-1}$, the heavy list $E(\Gh_{k-1})$, the vectors
$\xxrand_j,\xxres_j,\xxproxy_j$, $j\le10^{3}\log n$, stored by coordinate, the
proxy masses $\xxvert_v$, and the packed distance oracle $D$ of
\Cref{lem:packed} on the rows of $[\xxrand_1\cdots\xxrand_{10^{3}\log n}]$.
\end{itemize}
\textbf{Updates} during block $k$.
\begin{itemize}[leftmargin=*,itemsep=1pt]
\item $\textsc{Insert}(e,\ww_e)$: insert $e$ into $T_u,T_v$ and into
$\Gnew_k$; if $e$ becomes current light at an endpoint, insert it into
$T^{\mathrm{light}}$ there with active weight $0$ and append it to
$\mathrm{Changed}$.
\item $\textsc{Delete}(e)$: remove $e$ from $T_u,T_v$, from
$T^{\mathrm{light}}_u,T^{\mathrm{light}}_v$, from $\Gnew_k$ if present, and
$P.\textsc{Delete}(e)$ if present.
\item $\textsc{Weight}(e,\ww'_e)$ with $\ww'_e>\ww_e$: update the key in
$T_u,T_v$ and append the increase $\ww'_e-\ww_e$ to $\Gnew_k$.
With $\ww'_e<\ww_e$: update the key in $T_u,T_v$, the active weights in
$T^{\mathrm{light}}_u,T^{\mathrm{light}}_v$ (if active), and replace the copy in
$P$ by $P.\textsc{Delete}(e)$, $P.\textsc{Insert}(e,\ww'_e)$.
\item After every update, for each endpoint $v$ touched: the edge that left
and the edge that entered the $\dheavy$ largest keys of $T_v$ (at most one
each) move between heavy and light at $v$; move it out of or into
$T^{\mathrm{light}}_v$, set its active weight to $0$ there, and append it to
$\mathrm{Changed}$.
\end{itemize}
\end{minipage}}
\caption{The maintained state and the processing of the updates of a block;
the query is \Cref{fig:block}.}
\label{fig:state}
\end{figure}

\begin{lemma}[updates]\label{lem:updates}
An update costs $O(\log n)$ time plus amortized $O(\log^{3}n)$ in $P$, and
appends at most $4$ edges to $\mathrm{Changed}$.
At every moment $P$ holds $\Gret_k$, $\Gnew_k$ holds the new edges and
increases with $\LL(G)=\LL(\Gret_k)+\LL(\Gnew_k)$, and the active weights of
$T^{\mathrm{light}}_v$ are positive exactly on the edges that are light at
$v$ both at the previous refresh and now, where they equal the weight in
$\Gret_k$.
\end{lemma}
\begin{proof}
An update changes the key of one edge at two vertices.
At a vertex $v$, the set of the $\dheavy$ largest keys of $T_v$ changes by at
most one edge leaving and one entering, found in $O(\log n)$ by rank
queries in $T_v$; so at most two status changes per endpoint, four in all,
each an $O(\log n)$ tree operation.
The invariants on $P$ and $\Gnew_k$ are immediate from the case list of
\Cref{fig:state}: a deletion removes the retained copy, a decrease replaces
it by the smaller weight, an increase or an insertion is held, and the sum
$\LL(\Gret_k)+\LL(\Gnew_k)$ is the current Laplacian in each case.
An active weight is set to $0$ whenever the status changes and is otherwise
the weight of the copy in $P$, which is the weight in $\Gret_k$.
\end{proof}

\section{The refresh of one block}\label{sec:refresh}

Throughout this section $G=(V,E,\ww)$ is the connected graph at the query of
block $k$, with the heavy/light split $\Gh,\Gl$ and $\MM=\LL(\Gh)+\DD(\Gl)$ of
\Cref{def:heavy} for $\dheavy=\log^{3}n$, read off the trees $T_v$.
The refresh takes two sparse references and produces the record of the
block.
How the references are sampled is the subject of \Cref{sec:transfer}; here
they are given.

The references are those of \Cref{eq:references}, drawn as in
\Cref{lem:bernoulli}; an attempt of the refresh is a run of
$\textsc{Attempt}$ (\Cref{fig:attempt}), and \Cref{lem:refresh} gives its
acceptance probability, the energies \Cref{eq:energies} of an accepted
attempt, its all-pairs guarantee and its cost.

\begin{definition}[the record of a refresh]\label{def:record}
For an accepted attempt, the \emph{proxy masses} are
$\xxvert_v\defeq\frac1{10^{3}\log n}\sum_j(\xxproxy_j)_v^{2}$, the packed oracle
$D$ of \Cref{lem:packed} is built on the rows of
$\XX\defeq[\xxrand_1\cdots\xxrand_{10^{3}\log n}]$, and the \emph{public resistance
estimate} is
\begin{equation}\label{eq:score}
\Reffref_{uv}\ \defeq\
\begin{cases}
D(u,v)/(10^{3}\log n), & uv\in\Gh,\\[3pt]
\dfrac1{10^{3}\log n}\displaystyle\sum_j2\bigl((\xxres_j)_u-(\xxres_j)_v\bigr)^{2}
+4\xxvert_u+4\xxvert_v, & uv\in\Gl .
\end{cases}
\end{equation}
The \emph{record} of the refresh is $\LLtil$, the list $E(\Gh)$, the vectors
$\xxrand_j,\xxres_j,\xxproxy_j$, the masses $\xxvert$ and the oracle $D$.
\end{definition}

\begin{lemma}[the public estimate]\label{lem:public}
On the events of \Cref{lem:refresh}(b) and \Cref{lem:packed}:
\begin{enumerate}[label=(\alph*),leftmargin=*]
\item $\Reffref_e\ge\frac12\Reff(G)_e$ for every edge $e$, and the light branch
of \Cref{eq:score} evaluated at any pair $u,v$ is at least
$\frac12(\oone_u-\oone_v)^{\top}\LL(G)^{\dag}(\oone_u-\oone_v)$;
\item $\sum_e\ww_e\Reffref_e\le10^{7}n$;
\item $\Reffref_e$ is evaluated in $\Odouble(1)$ time for a heavy edge and
$O(\log n)$ time for a light edge, from the record.
\end{enumerate}
\end{lemma}
\begin{proof}
(a) For a heavy edge, $D(u,v)\ge\norm{\XX_u-\XX_v}_2^{2}=\sum_j((\xxrand_j)_u-(\xxrand_j)_v)^{2}$
and \Cref{lem:refresh}(b) gives the bound.
For the light branch, $\xxrand_j=\xxres_j+\xxproxy_j$ and $(x+y)^{2}\le2x^{2}+2y^{2}$
give
\[
((\xxrand_j)_u-(\xxrand_j)_v)^{2}\le2((\xxres_j)_u-(\xxres_j)_v)^{2}+4(\xxproxy_j)_u^{2}+4(\xxproxy_j)_v^{2},
\]
so the light branch dominates the average of \Cref{lem:refresh}(b).

(b) For the heavy edges, $D(u,v)\le10^{3}\norm{\XX_u-\XX_v}_2^{2}$ and
\Cref{eq:energies} give
$\sum_{e\in\Gh}\ww_e\Reffref_e\le10^{3}\cdot\frac1{10^{3}\log n}\sum_j\xxrand_j{}^{\top}\LL(G)\xxrand_j\le1.2\cdot10^{3}n$.
For the light edges, summing the light branch against $\ww_e$ gives
$2\cdot\frac1{10^{3}\log n}\sum_j\Ecal2(\xxres_j,\Gl)+4\sum_v\dd(\Gl)_v\xxvert_v$,
which by \Cref{eq:energies} is at most
$4\cdot10^{8}(n/\dheavy+n/\log^{3}n)+4\cdot10^{3}n\le4\cdot10^{6}n+4\cdot10^{3}n$
as $\dheavy\ge200$.
The total is at most $10^{7}n$.

(c) The heavy branch is one oracle query; the light branch reads
$10^{3}\log n$ stored coordinates at each endpoint and the two masses.
\end{proof}

\section{From one block to the next}\label{sec:transfer}

Fix $k\ge2$, write $G=G_k$, $\Gret=\Gret_k$, $\Gnew=\Gnew_k$, and let the
record of refresh $k-1$ be that of $G_{k-1}$ with its split
$\Gh_{k-1},\Gl_{k-1}$.
All statements are on the events of \Cref{lem:refresh}(b), \Cref{lem:packed}
and \Cref{lem:bernoulli} for refresh $k-1$, and block $k$ is stable
(\Cref{def:stable}).

\begin{lemma}[transfer]\label{lem:transfer}
If block $k$ is stable:
\begin{enumerate}[label=(\alph*),leftmargin=*]
\item for every edge $e$ of $\Gret$, with its weight $\ww^{\mathrm{ret}}_e$ in
$\Gret$,
$\Reff(G)_e\le\Reff(\Gret)_e\le10\Reff(G_{k-1})_e\le20\Reffref_e$, where
$\Reffref$ is the estimate of refresh $k-1$;
\item for every $\xx\in\R^{V}$,
$\frac1{10}\Ecal2(\xx,G_{k-1})\le\Ecal2(\xx,\Gret)\le\Ecal2(\xx,G_{k-1})$, and
$\Ecaltil(\xx)\defeq\frac13\xx^{\top}\LLtil_{k-1}\xx$ satisfies
$\Ecaltil(\xx)\approx_{10}\Ecal2(\xx,\Gret)$.
\end{enumerate}
\end{lemma}
\begin{proof}
(a) $\LL(G)=\LL(\Gret)+\LL(\Gnew)\succeq\LL(\Gret)$ gives the first
inequality (Rayleigh monotonicity; $G$ is connected, and $\Gret$ has the
kernel of $G_{k-1}$ by \Cref{def:stable}, which is $\oone$).
$\frac1{10}\LL(G_{k-1})\preceq\LL(\Gret)\preceq\LL(G_{k-1})$ gives
$\LL(\Gret)^{\dag}\preceq10\LL(G_{k-1})^{\dag}$, the second inequality.
The third is \Cref{lem:public}(a).

(b) The two-sided bound is $\frac1{10}\LL(G_{k-1})\preceq\LL(\Gret)\preceq\LL(G_{k-1})$
evaluated at $\xx$.
As $\LLtil_{k-1}\approx_{10^{-3}}\LL(G_{k-1})$,
$\Ecaltil(\xx)\in[\frac{0.999}3,\frac{1.001}3]\Ecal2(\xx,G_{k-1})$, and
$\Ecal2(\xx,\Gret)/\Ecaltil(\xx)\in[\frac{0.3}{1.001},\frac{3}{0.999}]\subseteq[\frac1{10},10]$.
\end{proof}

\paragraph{The partition queries of the block.}
At the query of block $k$, with $P$ holding $\Gret$, the structure computes
$\Ecaltil_j\defeq\Ecaltil(\xxres_j)$ for the $10^{3}\log n$ residual vectors of
refresh $k-1$ and runs $\textsc{PartitionTwo}(\xxres_j;\Ecaltil_j)$ on $P$ for
each $j$; \Cref{lem:transfer}(b) is the hypothesis $\Ecaltil_j\approx_{10}\Ecal2(\xxres_j,\Gret)$
of \Cref{lem:partition}.
Write $\Ccal$ for the clusters after the last query (all queries are between
updates, so later queries only split earlier clusters, the centers and vertex
energies of every $\xxres_j$ are inherited, and the certificate of each query
holds on the final clusters because a split only divides the degrees), $C(e)$
for the cluster of $e$ in $\Ccal$, and $\Ecut$ for the union of the cut edges
returned, which are held outside $P$ until the end of the block.
The certificates give
\begin{equation}\label{eq:Z}
\begin{aligned}
\sum_j\sum_{C\in\Ccal}\sum_{v\in V(C)}\deg_C(v)\Ecalvert2(\xxres_j,C,v)
&\le Z\ \defeq\ 10^{8}\log^{4}n\sum_j\Ecaltil_j\\
&\le 10^{8}\log^{4}n\cdot\frac{1.001}3\cdot10^{3}\log n\cdot2\cdot10^{8}\Bigl(\frac n{\dheavy}+\frac n{\log^{3}n}\Bigr),
\end{aligned}
\end{equation}
the last by \Cref{eq:energies} for refresh $k-1$ and \Cref{lem:transfer}(b).

\paragraph{The three streams.}
For an edge $e=uv$ of $\Gret$ that is in no cut set and whose status at both
endpoints is unchanged since refresh $k-1$, write $\ww_e$ for its weight in
$\Gret$ and
\[
r_e\ \defeq\ \frac1{10^{3}\log n}\sum_j\bigl((\xxres_j)_u-(\xxres_j)_v\bigr)^{2},
\qquad
\pi_e\ \defeq\ \frac1Z\sum_j\bigl(\Ecalvert2(\xxres_j,C(e),u)+\Ecalvert2(\xxres_j,C(e),v)\bigr),
\]
so that $\sum_e\pi_e\le1$ by \Cref{eq:Z} and, by the envelope
\Cref{eq:envelope} summed over $j$,
\begin{equation}\label{eq:resenv}
\ww_e\,r_e\ \le\ \frac{4Z}{10^{3}\log n}\,\pi_e .
\end{equation}
Given a rate $c>0$, the routine $\textsc{Proposals}(c)$ of \Cref{fig:proposals}
proposes edges of $\Gret$ through three independent streams and lists the
remaining edges directly.

\begin{figure}[H]\centering
\framebox{\begin{minipage}{0.96\linewidth}\small
$(E',q)=\textsc{Proposals}(c)$\\[2pt]
\underline{Input}: the state of \Cref{fig:state} after the partition queries
of the block, with $P$ holding $\Gret$ minus $\Ecut$; a rate $c>0$.\\
\underline{Output}: a set $E'$ of edges of $\Gret$ and, for each $e\in E'$, its
proposal probability $q_e$.
\begin{enumerate}[leftmargin=*,itemsep=1pt]
\item (heavy stream) For every $e=uv$ in the stored list $E(\Gh_{k-1})$ that is
still an edge of $\Gret$: propose $e$ independently with probability
$\min\{1,\ c\,\ww_e\,D(u,v)/(10^{3}\log n)\}$.
\item (vertex stream) For every vertex $v$: run a Poisson process of rate
$8c\,\xxvert_v$ per unit of active weight along the active weights of
$T^{\mathrm{light}}_v$, and propose every edge hit.
\item (residual stream) Build a tree over the copies $(C,v)$, $C\in\Ccal$,
with weights $\deg_C(v)\sum_j\Ecalvert2(\xxres_j,C,v)$.
Draw $\mathrm{Poisson}(8cZ/(10^{3}\log n))$ times: with probability
$1-\sum_{C,v}\deg_C(v)\sum_j\Ecalvert2(\xxres_j,C,v)/Z$ return nothing, else
draw $(C,v)$ from the tree with probability proportional to its weight, then
a uniformly random live edge of $C$ at $v$ (\Cref{lem:partition}); propose
every edge returned.
\item (direct list) Propose, with $q_e=1$, every edge of $\Gret$ in $\Ecut$ or
in $\mathrm{Changed}$.
\item For every proposed $e$, compute $q_e$ as $1$ minus the product over the
streams of Steps 1--3 that cover $e$ of the probability that the stream
misses $e$: $1-\min\{1,c\ww_eD(u,v)/(10^{3}\log n)\}$ for the heavy
stream, $e^{-8c\xxvert_u\ww_e}$ and $e^{-8c\xxvert_v\ww_e}$ for the vertex
stream at an endpoint where $e$ is active, and $e^{-8cZ\pi_e/(10^{3}\log n)}$
for the residual stream.
Return $E'$ and $q$.
\end{enumerate}
\end{minipage}}
\caption{The proposals of a block from the record of the previous refresh.
Step 1 is the heavy route, Step 2 the vertex route (the proxy terms of
\Cref{eq:score}), Step 3 the residual route through the certificates of the
partition queries, Step 4 the edges those routes do not cover.}
\label{fig:proposals}
\end{figure}

\begin{lemma}[proposals]\label{lem:proposals}
On the events named at the start of the section, $\textsc{Proposals}(c)$
proposes every edge $e$ of $\Gret$ independently with a probability $q_e$
that it returns, with
\[
q_e\ \ge\ \frac12\min\{1,\ c\,\ww_e\Reffref_e\}\ \ge\ \frac12\min\Bigl\{1,\ \frac{c}{20}\,\ww_e\Reff(G)_e\Bigr\},
\]
the expected number of proposals is at most
$10^{4}cn+10^{23}c\log^{4}n\,(n/\dheavy+n/\log^{3}n)+|\Ecut|+|\mathrm{Changed}|$, and
the time is $\Odouble(n\dheavy)+O(n\log^{2}n)$ plus $O(\log n)$ per proposal.
\end{lemma}
\begin{proof}
\emph{Independence and the value of $q_e$.}
The heavy stream is an independent coin per edge.
In the vertex stream, the arrivals of a Poisson process along the active
weights at $v$ hit disjoint edges as independent Poisson variables of means
$8c\xxvert_v\ww_e$ (the active weight of an active edge is $\ww_e$); the
processes at different vertices are independent.
In the residual stream, the number of draws returning $e$ is Poisson with
mean $8cZ/(10^{3}\log n)$ times the probability that one draw returns $e$,
which is $\deg_{C(e)}(u)\sum_j\Ecalvert2(\xxres_j,C(e),u)/Z\cdot1/\deg_{C(e)}(u)$
plus the same at $v$, that is $\pi_e$, so the mean is $8cZ\pi_e/(10^{3}\log n)$;
thinning a Poisson number of independent draws gives independent counts
over the edges.
The three streams are independent, so $e$ is proposed independently of the
other edges with the probability $q_e$ computed in Step 5, and the edges of
Step 4 have $q_e=1$.

\emph{The lower bound.}
Let $e=uv$ be an edge of $\Gret$ not in Step 4; then $e$ is in no cut set
and its status at both endpoints is that of refresh $k-1$.
If $e\in\Gh_{k-1}$, the heavy stream alone gives
$q_e\ge\min\{1,c\ww_eD(u,v)/(10^{3}\log n)\}=\min\{1,c\ww_e\Reffref_e\}$ by
\Cref{eq:score}.
If $e\in\Gl_{k-1}$, it is active at both endpoints, so by \Cref{eq:resenv}
the three exponents of Step 5 sum to at least
$8c\ww_e(\xxvert_u+\xxvert_v)+2c\ww_er_e\ge c\ww_e\Reffref_e$ (the light branch of
\Cref{eq:score} is $2r_e+4\xxvert_u+4\xxvert_v$), hence
$q_e\ge1-e^{-c\ww_e\Reffref_e}\ge\frac12\min\{1,c\ww_e\Reffref_e\}$, as
$1-e^{-x}\ge x/2$ for $x\le1$ and $1-e^{-x}\ge1-e^{-1}\ge\frac12$ for $x\ge1$;
the heavy case satisfies the same bound.
The second inequality of the lemma is \Cref{lem:transfer}(a).

\emph{The number of proposals.}
The expected count is the sum of the stream means.
Heavy: $\sum_{e\in\Gh_{k-1}}c\ww_eD(u,v)/(10^{3}\log n)\le c\sum_{e\in\Gh_{k-1}}\ww^{(k-1)}_e\Reffref_e\le1.2\cdot10^{3}cn$,
by $\ww_e\le\ww^{(k-1)}_e$ and the heavy part of the proof of
\Cref{lem:public}(b).
Vertex: $\sum_v8c\xxvert_v\cdot(\text{active weight at }v)\le8c\sum_v\xxvert_v\dd(\Gl_{k-1})_v\le8\cdot10^{3}cn$,
since the active edges at $v$ are light at $v$ in $G_{k-1}$ with weights at
most those in $G_{k-1}$, and
$\sum_v\dd(\Gl_{k-1})_v\xxvert_v=\frac1{10^{3}\log n}\sum_j\xxproxy_j{}^{\top}\DD(\Gl_{k-1})\xxproxy_j\le10^{3}n$
by \Cref{eq:energies}.
Residual: $8cZ/(10^{3}\log n)\le10^{23}c\log^{4}n(n/\dheavy+n/\log^{3}n)$ by
\Cref{eq:Z}.
Direct: $|\Ecut|+|\mathrm{Changed}|$.

\emph{Time.}
Step 1 scans at most $n\dheavy$ stored edges at $\Odouble(1)$ each
(\Cref{lem:packed}).
Step 2 generates each arrival and locates it in $T^{\mathrm{light}}_v$ by the
subtree sums in $O(\log n)$.
Step 3 computes the weight of each of the at most $10^{3}n\log^{2}n$ copies
from $10^{3}\log n$ stored vertex energies, $O(n\log^{3}n)$ in all, builds the
tree in $O(n\log^{2}n)$, and a draw costs $O(\log n)$ for the copy and
$O(\log n)$ for the live edge.
Step 4 is a scan of two lists.
Step 5 reads, per proposed edge, one oracle value, two masses, and
$\pi_e$ from the stored vertex energies of the two copies, $O(\log n)$ in all
(the $10^{3}\log n$ terms of $\pi_e$).
\end{proof}

\section{The block procedure and the proof of the block lemma}\label{sec:proof}

\begin{figure}[H]\centering
\framebox{\begin{minipage}{0.96\linewidth}\small
$H=\textsc{Query}(\epsilon)$ at the end of block $k$\\[2pt]
\underline{Input}: the state of \Cref{fig:state}; $n^{-10}\le\epsilon\le\frac12$.\\
\underline{Output}: a graph $H$ with $\LL(H)\approx_{\epsilon}\LL(G_k)$; the record of
refresh $k$ replaces that of refresh $k-1$.
\begin{enumerate}[leftmargin=*,itemsep=1pt]
\item If $k\ge2$: for $j\le10^{3}\log n$ compute
$\Ecaltil_j\leftarrow\frac13\xxres_j{}^{\top}\LLtil_{k-1}\xxres_j$ and run
$\textsc{PartitionTwo}(\xxres_j;\Ecaltil_j)$ on $P$; hold the returned
$\Ecut$.
\item If $k\ge2$: $(E',q)\leftarrow\textsc{Proposals}(c)$ with
$c\defeq2\cdot10^{10}\log n+4\cdot10^{4}\epsilon^{-2}\log n$.
\item For every $e\in E'$, with $\ww_e$ its weight in $\Gret_k$ and $\Reffref_e$
from \Cref{eq:score} (heavy branch iff $e\in E(\Gh_{k-1})$):
\[
\varrho_e\leftarrow10^{10}\log n\cdot\ww_e\Reffref_e,\qquad
p_e\leftarrow\min\{1,\ 2\cdot10^{4}\epsilon^{-2}\log n\cdot\ww_e\Reffref_e\};
\]
with two independent coins, put $e$ into $R$ with probability
$(1-e^{-\varrho_e})/q_e$ and into $O$ with probability $p_e/q_e$.
\item $\LLtil_k\leftarrow\sum_{e\in R}\frac{\ww_e}{1-e^{-\varrho_e}}\cchi_e\cchi_e^{\top}+\LL(\Gnew_k)$;
$H\leftarrow\{e\in O\text{ with weight }\ww_e/p_e\}\cup\Gnew_k$, parallel
copies merged by adding weights.
Read $E(\Gh_k)$ off the trees $T_v$ and put
$\MMtil_k\leftarrow\DD(\Gl_k)+\LLtil_k(\Gh_k)$.
\item (refresh) $\ZZ(\LLtil_k)\leftarrow\textsc{LsolverKLPSS}(\LLtil_k,\frac1{40})$,
$\ZZ(\MMtil_k)\leftarrow\textsc{LsolverKLPSS}(\MMtil_k,1/\log n)$; run one attempt of
\Cref{lem:refresh} (\Cref{fig:attempt}); if it is rejected, redraw for
every proposed edge the coin of Step 3 that decides its membership in $R$
and repeat Steps 4--5, at most $10^{2}\log n$ times in all; build the
record of \Cref{def:record} from the accepted attempt.
\item Reinsert $\Ecut$ into $P$; for every edge of $\Gnew_k$:
$P.\textsc{Insert}$ (a new edge) or $P.\textsc{Delete}$ then $P.\textsc{Insert}$ with the
current weight (an increase); for every edge of $\mathrm{Changed}\cup\Gnew_k$
set its active weight in $T^{\mathrm{light}}$ to its current weight where it
is current light; empty $\mathrm{Changed}$ and $\Gnew_k$.
Return $H$.
\end{enumerate}
\end{minipage}}
\caption{The query of a block.
For $k=1$, $\Gret_1$ is empty, $R=\emptyset$ and $\LLtil_1=\LL(\Gnew_1)$.}
\label{fig:block}
\end{figure}

\begin{lemma}[the two samples]\label{lem:samples}
On the events named at the start of \Cref{sec:transfer}, with probability
$1-n^{-11}$ over the coins of Steps 2--3 of \Cref{fig:block}:
\begin{enumerate}[label=(\alph*),leftmargin=*]
\item $\LLtil_k\approx_{10^{-3}}\LL(G_k)$, $\MMtil_k\approx_{10^{-3}}\LL(\Gh_k)+\DD(\Gl_k)$,
both with at most $10^{18}n\log n$ nonzero entries, and the intensities
$\varrho_e$ satisfy the hypothesis of \Cref{lem:bernoulli} for $G_k$;
\item $\LL(H)\approx_{\epsilon}\LL(G_k)$ and $H$ has at most
$10^{12}\epsilon^{-2}n\log n$ edges.
\end{enumerate}
\end{lemma}
\begin{proof}
Let $e$ be an edge of $\Gret_k$ with weight $\ww_e$ there.
By \Cref{lem:proposals}, $q_e\ge\frac12\min\{1,c\ww_e\Reffref_e\}$ with
$c=2\cdot10^{10}\log n+4\cdot10^{4}\epsilon^{-2}\log n$, so
$q_e\ge\min\{1,\varrho_e\}\ge1-e^{-\varrho_e}$ and $q_e\ge p_e$: the two
thinning probabilities of Step 3 are at most $1$, and $e$ enters $R$ with
probability exactly $1-e^{-\varrho_e}$ and $O$ with probability exactly
$p_e$, independently over the edges (independent proposals, independent
coins).

(a) By \Cref{lem:transfer}(a), $\Reff(G_k)_e\le20\Reffref_e$, so
$\varrho_e\ge5\cdot10^{8}\log n\,\ww_e\Reff(G_k)_e$; and by
\Cref{lem:public}(b) for $G_{k-1}$ and $\ww_e\le\ww^{(k-1)}_e$,
$\sum_e\varrho_e\le10^{10}\log n\cdot10^{7}n=10^{17}n\log n$.
\Cref{lem:bernoulli} with $E_{\mathrm{s}}$ the edges of $\Gret_k$,
$E_{\mathrm{x}}=\Gnew_k$ and $G=G_k$ (a multigraph when an edge has a retained
and an increase copy; the lemma is a statement about sums of rank-one
matrices and does not use simplicity) gives the approximations with
probability $1-2n^{-14}$ and at most $2\cdot10^{17}n\log n+n$ retained edges
with probability $1-n^{-14}$, two nonzero entries each.

(b) $p_e\ge\min\{1,10^{3}\epsilon^{-2}\log n\cdot\ww_e\Reff(G_k)_e\}$ by
\Cref{lem:transfer}(a), and the edges of $\Gnew_k$ have $p_e=1$;
\Cref{lem:sparsify} on the multigraph $\Gret_k\cup\Gnew_k$, whose Laplacian
is $\LL(G_k)$, gives $\LL(H)\approx_{\epsilon}\LL(G_k)$ with probability
$1-n^{-12}$ (merging parallel copies does not change the Laplacian) and
at most $2(2\cdot10^{4}\epsilon^{-2}\log n\cdot10^{7}n+n)\le10^{12}\epsilon^{-2}n\log n$
edges with probability $1-\exp(-10^{3}\log n/3)$.
\end{proof}

\begin{proof}[Proof of \Cref{lem:stable}]
\emph{Correctness and probability.}
Call block $k$ \emph{good} if the events of \Cref{lem:bernoulli}(a) and
\Cref{lem:samples}(a) for the $\LLtil_k,\MMtil_k$ of every trial of Step 5, of
\Cref{lem:chain} for the two operators of every trial, of
\Cref{lem:refresh}(b) and \Cref{lem:packed} for the accepted attempt, and of
\Cref{lem:samples}(b) for $H$ all hold, and the refresh accepts within
$10^{2}\log n$ trials.
Each trial of Step 5 uses fresh retention coins, fresh solver randomness and
fresh Gaussian vectors, so by \Cref{lem:refresh}(a), \Cref{lem:bernoulli}(a)
and \Cref{lem:chain} it is accepted with probability at least
$1-10^{-2}-3n^{-15}-2n^{-14}-2n^{-14}\ge1-\frac1{50}$.
Condition on the history up to the query of block $k$: the graph $G_k$, the
state, and the record of refresh $k-1$ are then fixed, and the randomness
of block $k$ -- the Gaussian vectors, the coins of \Cref{fig:proposals} and
Step 3, the solver and oracle randomness -- is fresh and independent of them.
If block $k-1$ is good, the hypotheses of \Cref{lem:transfer},
\Cref{lem:proposals} and \Cref{lem:samples} hold for every fixed $G_k$
that is stable (they are properties of the record and of
$G_{k-1}$, not of the adversary's choices), so conditionally on the history
block $k$ is good with probability at least
$1-10^{2}\log n\,(n^{-11}+4n^{-14})-2n^{-14}-n^{-12}-n^{-15}-50^{-10^{2}\log n}\ge1-n^{-10}$;
block $1$ is good with the same probability since $\Gret_1=\emptyset$ makes
\Cref{lem:samples}(a) trivial.
By induction and a union bound, all $Q$ blocks are good with probability
$1-Qn^{-10}$, and on a good block $H$ is as claimed.
The adversary sees $H$, which depends on the record of refresh $k-1$, but
the properties of that record used in block $k$ are the deterministic
consequences of block $k-1$ being good, and nothing about the fresh
randomness of block $k$ is revealed before $G_k$ is chosen; this is the
whole adaptivity argument.

\emph{Time of the updates.}
By \Cref{lem:updates}, $O(\log n)$ per update plus amortized $O(\log^{3}n)$
per operation of $P$; the operations of $P$ in block $k$ are the at most $n$
updates, the reinsertion of $\Ecut$ (amortized $10^{-1}$ per operation of $P$,
\Cref{lem:partition}), and at most $2n$ operations in Step 6, so $O(n)$
operations per block and $O(U\log^{3}n+Qn\log^{3}n)$ in all.

\emph{Time of a query, with $\dheavy=\log^{3}n$.}
Step 1: $10^{3}\log n$ energies against $\LLtil_{k-1}$, $O(n\log n)$ nonzeros
each, $O(n\log^{2}n)$; $10^{3}\log n$ partition queries at amortized
$O(n\log^{2}n)$ each, $O(n\log^{3}n)$.
Step 2: by \Cref{lem:proposals} with $c=O(\epsilon^{-2}\log n)$, the time is
$\Odouble(n\log^{3}n)+O(n\log^{2}n)$ plus $O(\log n)$ per proposal, and the
expected number of proposals is
\[
O\bigl(\epsilon^{-2}\,n\log^{2}n\bigr)+|\Ecut|+|\mathrm{Changed}|
=O\bigl(\epsilon^{-2}\,n\log^{2}n\bigr),
\]
using $\log^{4}n(n/\dheavy+n/\log^{3}n)=2n\log n$ and $|\mathrm{Changed}|\le4n$;
the Poisson counts exceed twice their means with probability
$\exp(-\Omega(n))$, so the proposals cost $O(\epsilon^{-2}n\log^{3}n)$.
Step 3: $O(\log n)$ per proposed edge (\Cref{lem:public}(c)), the same bound.
Step 4: $O(n\log n)$ for $\LLtil_k$ and $H$, $O(n\dheavy)=O(n\log^{3}n)$ for
the heavy list and $\MMtil_k$.
Step 5: a trial redraws one coin per proposed edge, $O(\epsilon^{-2}n\log^{2}n)$
by the count of Step 2, rebuilds $\LLtil_k$ and $\MMtil_k$ in $O(n\log n)$,
makes two calls of \Cref{lem:chain}, $O(n\log^{2}n)$, and runs one attempt,
$\Odouble(n\log^{2}n)$ (\Cref{lem:refresh}(c)); there are at most $10^{2}\log n$
trials on a good block, with expected number at most $2$ (each is accepted
with probability at least $1-\frac1{50}$), so $\Odouble(\epsilon^{-2}n\log^{2}n)$ in
expectation and $\Odouble(\epsilon^{-2}n\log^{3}n)$ always; the oracle of \Cref{lem:packed} with $c=12$ in
$O(n\log^{2}n)$ (on the event of \Cref{lem:refresh}(b) the nonzero distances
between rows of $\XX$ lie in $[n^{-6},n^{6}]$, as every resistance of a
connected graph with weights in $[n^{-10},n^{10}]$ lies in $[n^{-11},n^{11}]$);
the proxy masses in $O(n\log n)$.
Step 6: $O(n)$ operations of $P$, counted above, and $O(n\log n)$ tree
operations.
In all a query costs $\Odouble(\epsilon^{-2}n\log^{3}n)$, which summed over
the blocks is the second term of the theorem.
\end{proof}

\section{Without the stability hypothesis}\label{sec:general}

Stability (\Cref{def:stable}) enters \Cref{lem:transfer} only: it is what
makes the previous estimate an upper bound on the new resistances and the previous
energies estimates of the new ones.
It can be enforced by the structure itself, at the price of $O(\log n)$
internal refreshes per block.

\paragraph{The halving schedule.}
The updates of a block are now recorded and applied at its query.
Let $G_{k-1}$ be the graph at the previous query and $G_k$ the graph after
the block's updates, and write $\delta_e\defeq\ww^{(k-1)}_e-\ww^{(k)}_e>0$
for the edges $e$ of $G_{k-1}$ whose weight decreased or that were deleted
($\ww^{(k)}_e\defeq0$), $F$ for the set of these edges, and $\Gnew_k$ for the
new part as in \Cref{sec:main}.
Put $t\defeq30\log n$ and, for $0\le i\le t$,
\begin{align*}
\LL(G_{k,i})&\defeq\LL(G_k)+2^{-i}\sum_{e\in F}\delta_e\cchi_e\cchi_e^{\top},
\qquad\text{so that}\\
G_{k,0}&=G_{k-1}+\Gnew_k,\qquad
\LL(G_{k,t})=\LL(G_k)+n^{-30}\sum_{e\in F}\delta_e\cchi_e\cchi_e^{\top}.
\end{align*}
The structure runs $t+1$ internal blocks: block $i$ ($1\le i\le t$) applies
$\textsc{Weight}(e,\ww^{(k)}_e+2^{-i}\delta_e)$ for every $e\in F$ (block $1$
also inserts $\Gnew_k$) and calls $\textsc{Query}(\frac12)$, discarding its
output; block $t+1$ applies $\textsc{Weight}(e,\ww^{(k)}_e)$ for $e\in F$ with
$\ww^{(k)}_e>0$ and $\textsc{Delete}(e)$ for the others, and calls
$\textsc{Query}(\epsilon)$, whose output is returned.
The intermediate weights lie in $[n^{-40},n^{10}]$; the imported statements
are used with this polynomial range, which changes the number of weight
classes of \Cref{lem:partition} from $20\log n+1$ to $50\log n+1$ and nothing
else in their proofs, as every bound there is stated for a polynomial range.

\begin{lemma}[the internal blocks are stable]\label{lem:halving}
With $G_{k,0}$ in the role of $G_{k-1}$, every internal block satisfies
\Cref{def:stable}: for $1\le i\le t$, $\LL(G_{k,i})\succeq\frac12\LL(G_{k,i-1})$,
and $\LL(G_k)\succeq\frac12\LL(G_{k,t})$.
\end{lemma}
\begin{proof}
For $1\le i\le t$, $\LL(G_{k,i-1})-\LL(G_{k,i})=2^{-i}\sum_{e\in F}\delta_e\cchi_e\cchi_e^{\top}\preceq\LL(G_{k,i})$,
because $G_{k,i}$ contains every $e\in F$ with weight at least
$2^{-i}\delta_e$; so $\LL(G_{k,i-1})\preceq2\LL(G_{k,i})$.
For the last block, $G_k$ is connected, so for every edge $e$,
$\cchi_e\cchi_e^{\top}\preceq\Reff(G_k)_e\LL(G_k)$ (for $\xx\perp\oone$,
$(\cchi_e^{\top}\xx)^{2}\le(\cchi_e^{\top}\LL(G_k)^{\dag}\cchi_e)(\xx^{\top}\LL(G_k)\xx)$
by Cauchy--Schwarz in the $\LL(G_k)$ inner product), and
$\Reff(G_k)_e\le(n-1)n^{10}$ (a path of at most $n-1$ edges of weight at
least $n^{-10}$), $\delta_e\le n^{10}$, $|F|\le n$; hence
\[
\LL(G_{k,t})-\LL(G_k)=n^{-30}\sum_{e\in F}\delta_e\cchi_e\cchi_e^{\top}
\preceq n^{-30}\cdot n\cdot n^{10}\cdot n^{11}\,\LL(G_k)\preceq\LL(G_k).
\qedhere
\]
\end{proof}

\begin{proof}[Proof of \Cref{thm:main}]
The block's updates are fixed when its query arrives, so the internal
blocks are a sequence of $t+1\le31\log n$ blocks of the model of
\Cref{sec:main}, each with at most $|F|+|\Gnew_k|\le2n$ updates and one
query, which satisfy \Cref{def:stable} by \Cref{lem:halving}; the graphs
$G_{k,i}$ are connected since $G_k$ is.
\Cref{lem:stable} applies to the whole sequence of internal blocks: each is
good with probability $1-n^{-10}$ conditionally on the past, so all
$(t+1)Q$ of them are good with probability $1-31Q\log n\cdot n^{-10}\ge1-Qn^{-9}$,
and on a good last block $H$ is as claimed.
The time is that of \Cref{lem:stable} for $(t+1)Q$ blocks with
$(t+1)U$ updates, the internal queries at $\epsilon=\frac12$ and the last at
$\epsilon_k$: $\Odouble((t+1)U\log^{3}n+(t+1)Qn\log^{3}n+n\log^{3}n\sum_k\epsilon_k^{-2})$.
\end{proof}

\paragraph{Toward one logarithm less.}
The following identity shows what a block of arbitrary deletions does to
the resistances, and how much of it the same mechanisms could absorb in a
single refresh.

\begin{lemma}[deleted edges]\label{lem:woodbury}
Let $G=(V,E,\ww)$ be connected, $F\subseteq E$ with $G-F$ connected, and let
$\SSS\defeq\WW_F^{1/2}\BB_F\LL(G)^{\dag}\BB_F^{\top}\WW_F^{1/2}\in\R^{F\times F}$ be
the leverage matrix of $F$.
Then $0\preceq\SSS\prec\II$,
\begin{align*}
\LL(G-F)^{\dag}-\LL(G)^{\dag}
&=\LL(G)^{\dag}\BB_F^{\top}\WW_F^{1/2}(\II-\SSS)^{-1}\WW_F^{1/2}\BB_F\LL(G)^{\dag}
\ \succeq\ 0,\\
\sum_{e\in E\setminus F}\ww_e\bigl(\Reff(G-F)_e-\Reff(G)_e\bigr)
&=\sum_{e\in F}\ww_e\Reff(G)_e
=\tr[\SSS] .
\end{align*}
Consequently, if $\yy_1,\dots,\yy_{s}$ are independent Gaussian vectors with
covariance $\LL(G-F)^{\dag}-\LL(G)^{\dag}$, $s=10^{3}\log n$, then with
probability $1-n^{-14}$ for every edge $e$ of $G-F$
\begin{align*}
\frac1{s}\sum_{i\le s}(\cchi_e^{\top}\yy_i)^{2}
&\in\bigl[\tfrac12,2\bigr]\cdot\bigl(\Reff(G-F)_e-\Reff(G)_e\bigr),\\
\sum_{e\in E\setminus F}\ww_e\,\frac1s\sum_{i\le s}(\cchi_e^{\top}\yy_i)^{2}
&\le2\,|F| .
\end{align*}
\end{lemma}
\begin{proof}
$\SSS\succeq0$ is clear, and $\SSS\preceq\II$ because
$\sum_{e\in F}\ww_e\LL(G)^{\dag/2}\cchi_e\cchi_e^{\top}\LL(G)^{\dag/2}\preceq\PPi$,
whose nonzero eigenvalues are those of $\SSS$; a unit eigenvector $\zz$ of
$\SSS$ with eigenvalue $1$ would give a vector
$\LL(G)^{\dag/2}\BB_F^{\top}\WW_F^{1/2}\zz\perp\oone$ in the kernel of
$\LL(G)-\LL(F)=\LL(G-F)$, contradicting connectivity, so $\SSS\prec\II$.
The identity is the Woodbury formula for
$\LL(G-F)=\LL(G)-\BB_F^{\top}\WW_F\BB_F$ on $\oone^{\perp}$.
For the trace identity,
$\sum_{e\notin F}\ww_e\cchi_e^{\top}\LL(G-F)^{\dag}\cchi_e=\tr[\LL(G-F)^{\dag}\LL(G-F)]=n-1$
and
$\sum_{e\notin F}\ww_e\cchi_e^{\top}\LL(G)^{\dag}\cchi_e=\tr[\LL(G)^{\dag}\LL(G)]-\sum_{e\in F}\ww_e\Reff(G)_e$;
subtracting gives $\sum_{e\in F}\ww_e\Reff(G)_e=\tr[\SSS]\le|F|$, as every
leverage is at most $1$.
The all-pairs statement is the $\chi^{2}$ tail with $s=10^{3}\log n$ degrees
of freedom and a union bound over $n^{2}$ pairs, exactly as in
\Cref{lem:refresh}(b), and the energy bound follows from it and the trace
identity.
\end{proof}
Thus after a block that deletes $F$ ($|F|\le n$), the retained edges whose
leverage rose are exactly the edges with large energy under the vectors
$\yy_i$, the total energy is at most $2|F|$, and the certificate of
\Cref{lem:partition} with the $\yy_i$ as query vectors would find them and
sample them with $O(\log^{4}n\cdot|F|\cdot\log n)$ proposals per block, the
same accounting as the residual stream.
What is missing is the vectors themselves: $\yy_i=\LL(G-F)^{\dag}\BB_F^{\top}\WW_F^{1/2}(\II-\SSS)^{1/2}\gg_i$
(by the identity, $\LL(G-F)^{\dag}\BB_F^{\top}\WW_F^{1/2}=\LL(G)^{\dag}\BB_F^{\top}\WW_F^{1/2}(\II-\SSS)^{-1}$, so the covariance is the difference above)
requires solving in $\LL(G-F)$, whose only sparse approximation at hand,
$\LLtil_{k-1}$, is a preconditioner of condition number $1/(1-\lambda_{\max}(\SSS))$,
unbounded when a block deletes most of a cut; and $\SSS$ itself is a dense
$|F|\times|F|$ matrix.
Drawing the $\yy_i$ in $\Odouble(n\log^{3}n)$ would remove the factor
$\log n$ of \Cref{thm:main}; this is the open item of the verification
log.

\paragraph{The old reference does not determine the new resistances.}
The following example shows that the vectors $\yy_i$, and in fact any
constant-factor approximation of $\LL(G-F)$, cannot be computed from the
reference of the previous refresh and the deleted edges alone, however much
time is spent: the retained graph $G-F$ itself must be consulted, through
the structures of \Cref{sec:state}.
\begin{lemma}[two graphs with the same reference]\label{lem:twographs}
Let $n\ge2^{20}$ with $n\equiv2\pmod4$.
There are connected graphs $G_1,G_{11}$ on the same $n$ vertices with weights
in $\{1,2,3\}$, the same weighted degrees, and a common set $F$ of $n-1$ edges
of weight $1$ with $G_1-F$ and $G_{11}-F$ connected, such that every $\LLtil$
with $\LLtil\approx_{4\cdot10^{-4}}\LL(G_1)$ satisfies
$\LLtil\approx_{10^{-3}}\LL(G_{11})$, while no symmetric matrix $\XX$ satisfies
both $\XX\approx_{3/4}\LL(G_1-F)$ and $\XX\approx_{3/4}\LL(G_{11}-F)$.
Hence a procedure whose input is $(\LLtil,F,\ww_F,\dd(G))$ outputs a
$\frac34$-approximation of $\LL(G-F)$ with probability at most $\frac12$ on
one of the two graphs.
\end{lemma}
\begin{proof}
Let $V=A\cup B$, $A=\{a_1,\dots,a_{n/2}\}$, $B=\{b_1,\dots,b_{n/2}\}$,
$F\defeq\{a_ib_i:i\le n/2\}\cup\{a_ib_{i+1}:i<n/2\}$, and for $j\in\{1,11\}$
let $H_j\defeq\{a_{2i}b_{2i+3}:i\le j\}$, $j$ edges outside $F$ with $2j$
distinct endpoints.
Fix perfect matchings $M^A_j$ of $A\setminus\{a_{2i}:i\le j\}$ and $M^B_j$ of
$B\setminus\{b_{2i+3}:i\le j\}$ (sets of even size $n/2-j$).
$G_j$ has every pair inside $A$ and inside $B$ with weight $2$, raised to $3$
on $M^A_j\cup M^B_j$, the edges of $F$ with weight $1$ and the edges of $H_j$
with weight $1$.
Every vertex of $A$ has weighted degree $2(n/2-1)+1+\dd(F)_v$, the $1$ from
its matching edge or from its edge of $H_j$; likewise in $B$; so the degrees
do not depend on $j$, and $G_j-F$ is connected through $H_j$.

Write $Q_A\defeq\sum_{\{a,a'\}\subseteq A}(\xx_a-\xx_{a'})^{2}$, likewise $Q_B$;
then $2Q_A+2Q_B\le\xx^{\top}\LL(G_1)\xx$ and
$Q_A=\frac n2\sum_{a\in A}(\xx_a-\bar\xx_A)^{2}$ for the mean $\bar\xx_A$.
$\LL(G_{11})-\LL(G_1)$ is the sum of $\cchi_h\cchi_h^{\top}$ over the ten
edges $h\in H_{11}\setminus H_1$ plus
$\LL(M^A_{11})-\LL(M^A_1)+\LL(M^B_{11})-\LL(M^B_1)$ with unit weights.
For $h=a_pb_q$ and every $i\le n/2$,
$(\xx_{a_p}-\xx_{b_q})^{2}\le3(\xx_{a_p}-\xx_{a_i})^{2}+3(\xx_{a_i}-\xx_{b_i})^{2}+3(\xx_{b_i}-\xx_{b_q})^{2}$;
averaging over $i$ and using that the edges $a_ib_i$ have weight $1$ in $G_1$,
\[
(\xx_{a_p}-\xx_{b_q})^{2}\le\frac6n\Bigl(Q_A+\sum_i(\xx_{a_i}-\xx_{b_i})^{2}+Q_B\Bigr)
\le\frac{12}n\,\xx^{\top}\LL(G_1)\xx .
\]
For a unit perfect matching $M$ of a subset of $A$,
$\xx^{\top}\LL(M)\xx\le2\sum_{a\in A}(\xx_a-\bar\xx_A)^{2}=\frac4nQ_A\le\frac2n\xx^{\top}\LL(G_1)\xx$.
Summing, $|\xx^{\top}(\LL(G_{11})-\LL(G_1))\xx|\le\frac{128}n\xx^{\top}\LL(G_1)\xx\le2\cdot10^{-4}\xx^{\top}\LL(G_1)\xx$
for $n\ge2^{20}$, so $\LL(G_{11})\approx_{2\cdot10^{-4}}\LL(G_1)$ and
$\LLtil\approx_{4\cdot10^{-4}}\LL(G_1)$ gives $\LLtil\approx_{10^{-3}}\LL(G_{11})$.

For $\xx=\oone_A-\oone_B$ the edges inside $A$, inside $B$ and in $F$ carry no
energy in $G_j-F$, and each edge of $H_j$ carries $4$:
$\xx^{\top}\LL(G_j-F)\xx=4j$.
$\XX\approx_{3/4}\LL(G_1-F)$ forces $\xx^{\top}\XX\xx\le7$ and
$\XX\approx_{3/4}\LL(G_{11}-F)$ forces $\xx^{\top}\XX\xx\ge11$.
The input $(\LLtil,F,\ww_F,\dd(G))$ is the same for both graphs, so the
output has the same law, and the two success events are disjoint.
\end{proof}
The structure of \Cref{sec:state} holds $G-F$ exactly, in $P$ and in the
trees $T_v$, so \Cref{lem:twographs} does not close the question; it says
that the queries to $P$ must carry the information about $G-F$ that the
reference lacks -- in the example, the energy of $\oone_A-\oone_B$ in $G-F$ --
and the halving schedule obtains it by $O(\log n)$ refreshes.
Which $O(1)$ query vectors per block suffice is the open item.

\bibliographystyle{alpha}
\bibliography{../DenseGraphs}

\clearpage
\appendix
\section{The packed distance oracle}\label{sec:packed}

This section proves \Cref{lem:packed}: given the $n$ rows of the matrix
$\XX=[\xxrand_1\cdots\xxrand_{10^{3}\log n}]$ of a refresh, a data structure
built in $O(n\log^{2}n)$ time returns the squared distance of any two rows
within the factor $10^{3}$ in $\Odouble(1)$ time, so that the heavy branch of
\Cref{eq:score} and Step~1 of \Cref{fig:proposals} cost $\Odouble(1)$ per heavy
edge instead of the $O(\log n)$ of a direct evaluation.
The construction is the random-shift hashing of a Gaussian projection of
Indyk and Motwani~\cite{IM98:conf} and Datar, Immorlica, Indyk and
Mirrokni~\cite{DIIM04:conf}, read at every dyadic scale and reduced to
one parity bit per coordinate and scale; the parity bits of a vertex at a
scale are packed into $O(\log\log n)$ integers, and the number of
disagreeing bits of two such integers is read off one integer product,
the word-parallel use of multiplication of Fredman and
Willard~\cite{FW93:journal}.
No dimension-reduction theorem is used: the Gaussian projection serves
only to make every coordinate difference a Gaussian of variance the
squared distance, and the concentration is Hoeffding's
inequality~\cite{H63:journal}.

\paragraph{Model of computation.}
The oracle is a collection of records linked by pointers; every record
holds $O(1)$ pointers and integers.
An integer held by a record has at most $(4c+7)\log n$ bits, and the
operations addition, subtraction, comparison, multiplication, integer
division (rounding down) and remainder on such integers cost $O(1)$ each,
as does following a pointer.
No bitwise operation and no integer-addressed memory is used, so the
oracle runs on a pointer machine whose registers support these six
operations, and on any word RAM with word length at least $(4c+7)\log n$
and unit-cost multiplication, division and remainder.
Real numbers (the entries of $\XX$, the Gaussian draws) are held in fixed
point, as integer multiples of $2^{-w}$ for the integer $w$ fixed below, and
the Gaussian draws are treated as exact, as in \Cref{lem:refresh}; the
only rounding the analysis tracks is the rounding of the projected rows
in Step~3 of \Cref{fig:packed-pre}.
\Cref{lem:packed-model} at the end of the section states how the same
bounds are obtained with bitwise exclusive or and population count, or
with exclusive or and a lookup table, in place of multiplication.

\begin{lemma}\label{lem:packed}
Let $\XX\in\R^{n\times10^{3}\log n}$ have rows $\XX_v$ with all nonzero pairwise
distances $\norm{\XX_u-\XX_v}_2$ in $[n^{-c},n^{c}]$ for a constant $c\ge1$.
After $O(n\log^{2}n)$ randomized preprocessing creating $\Odouble(n\log n)$
records, a query $(u,v)$ returns $D(u,v)$ in $\Odouble(1)$ worst-case time with
$\norm{\XX_u-\XX_v}_2^{2}\le D(u,v)\le10^{3}\norm{\XX_u-\XX_v}_2^{2}$ for all
pairs simultaneously with probability $1-n^{-14}$; the answers are a fixed
function of the preprocessing, so adaptive queries are allowed.
\end{lemma}

\paragraph{The construction.}
The preprocessing draws $\GG\in\R^{10^{3}\log n\times10^{3}\log n}$ with
independent standard Gaussian entries and computes
\[
\YY\ \defeq\ \XX\GG\ \in\ \R^{n\times10^{3}\log n},
\]
with rows $\YY_v$, in $n\cdot(10^{3}\log n)^{2}=O(n\log^{2}n)$ operations.
Let $w$ be the least integer with $2^{w}\ge2^{6}n^{c+1}$, let $v_0$ be a fixed
vertex, and let the scales be $2^{i}$ for the integers
$i\in\{-w+6,\dots,w\}$, which are $2w-5\le2(c+1)\log n+9$ scales (the scales
are powers of two, so the offsets $2^{6}$ and $6$ are forced to be powers of
two).
For every coordinate $j$ the preprocessing draws a shift $\gamma_j$
uniformly from the multiples of $2^{-w}$ in $[0,2^{w+1})$ and stores, for
every vertex $v$, the shifted and rounded coordinate
\begin{equation}\label{eq:packed-shift}
\YYbar_{v,j}\ \defeq\ 2^{-w}\bigl\lfloor2^{w}(\YY_{v,j}-\YY_{v_0,j})\bigr\rfloor
+\gamma_j+2^{w+1},
\end{equation}
a multiple of $2^{-w}$, and for every scale $i$ the bit vector
$\bb_{v,i}\in\{0,1\}^{10^{3}\log n}$ with
\begin{equation}\label{eq:packed-bits}
(\bb_{v,i})_j\ \defeq\ \bigl\lfloor\YYbar_{v,j}/2^{i}\bigr\rfloor\bmod2 .
\end{equation}
A query $(u,v)$ finds by bisection over the scales an index $h$ with
$\norm{\bb_{u,h}-\bb_{v,h}}_1\le250\log n$ and
$\norm{\bb_{u,h-1}-\bb_{v,h-1}}_1>250\log n$ and returns $D\defeq2^{2h+1}$; the
analysis shows that, with high probability, the bits disagree in more than
$250\log n$ coordinates at every scale $2^{i}\le\norm{\XX_u-\XX_v}_2$ and in
at most $250\log n$ coordinates at every scale $2^{i}\ge10\norm{\XX_u-\XX_v}_2$,
which places $2^{h}$ in $(\norm{\XX_u-\XX_v}_2,\ 20\norm{\XX_u-\XX_v}_2)$.
The threshold $250\log n$ is not a power of ten: it has to separate the
disagreement rates $0.47$ and $0.1$ per coordinate established in the proof
with a margin that survives a union bound over $n^{2}$ pairs, and $10^{2}\log n$
is too close to the lower rate.

We first record the four facts about one coordinate that the analysis
uses: the law of a projected difference, the probability that two parity
bits disagree, the mean of that probability under a Gaussian input, and
Hoeffding's inequality.

\begin{lemma}\label{lem:packed-law}
For every pair $(u,v)$ the differences $\YY_{u,j}-\YY_{v,j}$,
$j=1,\dots,10^{3}\log n$, are independent $N(0,\norm{\XX_u-\XX_v}_2^{2})$
variables, and the pair $(\YY_{u,j}-\YY_{v_0,j},\YY_{v,j}-\YY_{v_0,j})$
depends on
$\GG$ only through its $j$-th column.
\end{lemma}
\begin{proof}
$\YY_{u,j}-\YY_{v,j}=(\XX_u-\XX_v)\gg_j$ for the $j$-th column $\gg_j$ of $\GG$,
a linear functional of a standard Gaussian vector with squared norm
$\norm{\XX_u-\XX_v}_2^{2}$, hence $N(0,\norm{\XX_u-\XX_v}_2^{2})$; the columns
$\gg_j$ are independent, and both entries of the pair are functions of
$\gg_j$.
\end{proof}

\begin{lemma}\label{lem:packed-parity}
Let $T\ge1$ and $M\ge1$ be integers, let $a,b\in\Z$, and let $g$ be uniform
on the integers of $[0,2TM)$.
Then
\[
\Pr\left[\left\lfloor\frac{a+g}{T}\right\rfloor\not\equiv
\left\lfloor\frac{b+g}{T}\right\rfloor\pmod 2\right]
\ =\ \frac{\dist(b-a,\,2T\Z)}{T}.
\]
\end{lemma}
\begin{proof}
The parity of $\lfloor m/T\rfloor$ is $2T$-periodic in $m\in\Z$, equal to
$0$ for $m\bmod2T\in[0,T)$ and to $1$ for $m\bmod2T\in[T,2T)$.
Since $g$ is uniform on $2TM$ consecutive integers, $(a+g)\bmod2T$ is
uniform on the residues modulo $2T$, so the probability equals the
fraction of residues $m\in\{0,\dots,2T-1\}$ with
$\lfloor m/T\rfloor\not\equiv\lfloor(m+d)/T\rfloor\pmod2$ for
$d\defeq(b-a)\bmod2T$.
\begin{itemize}[leftmargin=*]
\item If $d\le T$: the parities differ exactly for $m\in[T-d,T)$ (where
$m+d\in[T,T+d)$) and for $m\in[2T-d,2T)$ (where $m+d\in[2T,2T+d)$), that is
for $2d$ residues, a fraction $d/T$.
\item If $d>T$: the parity of $\lfloor(m+d)/T\rfloor$ is that of
$\lfloor(m-(2T-d))/T\rfloor$, and the substitution $m\mapsto m+(2T-d)$ shows
that the count equals the count for the difference $2T-d<T$, that is
$2(2T-d)$ residues, a fraction $(2T-d)/T$.
\end{itemize}
Both fractions equal $\dist(b-a,2T\Z)/T$, as
$\dist(b-a,2T\Z)=\min\{d,2T-d\}$.
\end{proof}

\begin{lemma}\label{lem:packed-wave}
Let $t>0$ and let $x\sim N(0,\sigma^{2})$.
The function $x\mapsto\dist(x,2t\Z)/t$ is $1/t$-Lipschitz with values in
$[0,1]$, and
\begin{align*}
\Ex\bigl[\dist(x,2t\Z)/t\bigr]&\ \ge\ 0.49 &&\text{if }t\le\sigma,\\
\Ex\bigl[\dist(x,2t\Z)/t\bigr]&\ \le\ \frac{\sigma}{t}\sqrt{\frac2\pi}
&&\text{for every }\sigma .
\end{align*}
\end{lemma}
\begin{proof}
The function $x\mapsto\dist(x,2t\Z)/t$ is continuous, even, $2t$-periodic,
equal to $|x|/t$ on $[-t,t]$, with slope $\pm1/t$ and values in $[0,1]$.
Its Fourier coefficients are, for $k\ge1$, by integration by parts,
\begin{align*}
\frac1t\int_{-t}^{t}\frac{|x|}{t}\cos\frac{k\pi x}{t}\,dx
&=\frac2t\int_{0}^{t}\frac{x}{t}\cos\frac{k\pi x}{t}\,dx
=-\frac{2}{k\pi t}\int_{0}^{t}\sin\frac{k\pi x}{t}\,dx
=-\frac{2\bigl(1-(-1)^{k}\bigr)}{k^{2}\pi^{2}},
\end{align*}
the sine coefficients vanish by evenness, and the mean is
$\frac1t\int_{-t}^{t}|x|/t\,dx\cdot\frac12=\frac12$.
The coefficients are absolutely summable, so the Fourier series converges
uniformly to the function (Corollary~2.3 of Chapter~2
of~\cite{SS03:book}):
\[
\frac{\dist(x,2t\Z)}{t}\ =\ \frac12-\frac{4}{\pi^{2}}\sum_{k\ge1\text{ odd}}
\frac{1}{k^{2}}\cos\frac{k\pi x}{t}\qquad\text{for all }x\in\R .
\]
Taking expectations termwise (the series is uniformly bounded), with the
characteristic function of $N(0,\sigma^{2})$ at $k\pi/t$,
\[
\Ex\Bigl[\cos\frac{k\pi x}{t}\Bigr]
=\exp\Bigl(-\frac{k^{2}\pi^{2}\sigma^{2}}{2t^{2}}\Bigr).
\]
If $t\le\sigma$ this is at most $e^{-\pi^{2}/2}\le0.0072$ for every $k\ge1$,
and $\sum_{k\text{ odd}}k^{-2}=\pi^{2}/8$, so
$\Ex[\dist(x,2t\Z)/t]\ge\frac12-\frac12\cdot0.0072\ge0.49$.
The upper bound is $\dist(x,2t\Z)/t\le|x|/t$ and
$\Ex[|x|]=\sigma\sqrt{2/\pi}$.
\end{proof}

\begin{lemma}[Theorem~2 of~\cite{H63:journal}; Theorem~2.8 of~\cite{BLM13:book}]
\label{lem:packed-hoeffding}
If $\xi_1,\dots,\xi_{N}$ are independent random variables with values in
$[0,1]$, then for every $\epsilon>0$,
\[
\Pr\Bigl[\sum_j\xi_j\ge\Ex\Bigl[\sum_j\xi_j\Bigr]+\epsilon N\Bigr]
\le e^{-2\epsilon^{2}N}
\quad\text{and}\quad
\Pr\Bigl[\sum_j\xi_j\le\Ex\Bigl[\sum_j\xi_j\Bigr]-\epsilon N\Bigr]
\le e^{-2\epsilon^{2}N}.
\]
\end{lemma}

The next two lemmas are the word-parallel Hamming distance: the inner
product of two bit strings is one digit of the product of the integer
holding the first string in fields of $f$ bits with the integer holding
the second string in the same fields in reverse order, because that
product is the convolution of the two strings and a field width above
$\log k$ keeps every convolution coefficient inside its own field.

\begin{lemma}\label{lem:packed-product}
Let $k\ge1$ and $f\ge1$ be integers with $2^{f}>k$, and let
$\uu,\vv\in\{0,1\}^{k}$.
Then
$0\le\sum_{j=1}^{k}\uu_j2^{(j-1)f}<2^{kf}$,
$0\le\sum_{j=1}^{k}\vv_j2^{(k-j)f}<2^{kf}$, and
\[
\sum_{j=1}^{k}\uu_j\vv_j
\ =\ \left\lfloor\frac{\bigl(\sum_{j=1}^{k}\uu_j2^{(j-1)f}\bigr)
\bigl(\sum_{j=1}^{k}\vv_j2^{(k-j)f}\bigr)}{2^{(k-1)f}}\right\rfloor\bmod2^{f}.
\]
\end{lemma}
\begin{proof}
The two bounds hold since every bit is at most $1$ and
$\sum_{j\le k}2^{(j-1)f}<2^{kf}$.
Expanding the product and grouping the terms by $d\defeq j-l$,
\[
\Bigl(\sum_{j=1}^{k}\uu_j2^{(j-1)f}\Bigr)
\Bigl(\sum_{l=1}^{k}\vv_l2^{(k-l)f}\Bigr)
=\sum_{j=1}^{k}\sum_{l=1}^{k}\uu_j\vv_l2^{(k-1+j-l)f}
=\sum_{d=-(k-1)}^{k-1}2^{(k-1+d)f}\sum_{j-l=d}\uu_j\vv_l .
\]
Each inner sum has $k-|d|\le k$ terms, each a product of bits, so it lies
in $\{0,\dots,k\}\subseteq\{0,\dots,2^{f}-1\}$.
A nonnegative integer has exactly one expansion in base $2^{f}$ with
digits in $\{0,\dots,2^{f}-1\}$, so the inner sums are the base-$2^{f}$
digits of the product, the sum for $d$ at position $k-1+d$.
Dividing by $2^{(k-1)f}$ and rounding down removes the digits at
positions below $k-1$, and the remainder modulo $2^{f}$ keeps the digit at
position $k-1$, which is $\sum_j\uu_j\vv_j$.
\end{proof}

\begin{lemma}\label{lem:packed-hamming}
Let $k,f,r\ge1$ be integers with $2^{f}>k$, and let
$\uu,\vv\in\{0,1\}^{rk}$ be given through the $2r$ integers
$\sum_{j=1}^{k}\uu_{(\ell-1)k+j}2^{(j-1)f}$ and
$\sum_{j=1}^{k}\vv_{(\ell-1)k+j}2^{(k-j)f}$,
$\ell=1,\dots,r$, and the two integers $\norm{\uu}_1$ and $\norm{\vv}_1$.
Then
\[
\norm{\uu-\vv}_1
=\norm{\uu}_1+\norm{\vv}_1
-2\sum_{\ell=1}^{r}\left(\left\lfloor
\frac{\bigl(\sum_{j=1}^{k}\uu_{(\ell-1)k+j}2^{(j-1)f}\bigr)
\bigl(\sum_{j=1}^{k}\vv_{(\ell-1)k+j}2^{(k-j)f}\bigr)}{2^{(k-1)f}}
\right\rfloor\bmod2^{f}\right),
\]
and the right-hand side is evaluated with at most $4r+3$ operations of
the model on integers below $2^{2kf}$, given the constants $2^{f}$ and
$2^{(k-1)f}$.
\end{lemma}
\begin{proof}
For bits, $|\uu_j-\vv_j|=\uu_j+\vv_j-2\uu_j\vv_j$; summing over $j$ gives
$\norm{\uu-\vv}_1=\norm{\uu}_1+\norm{\vv}_1-2\sum_j\uu_j\vv_j$, and the
inner product splits over the $r$ blocks of $k$ consecutive bits, each
given by \Cref{lem:packed-product}.
Every block costs one multiplication, one integer division, one
remainder and one addition into a running sum, and the remaining
operations are one multiplication by $2$ and two subtractions.
The largest integer handled is a product of two of the given integers,
each below $2^{kf}$ by \Cref{lem:packed-product}.
\end{proof}

\paragraph{The data structure.}
The bit vectors $\bb_{v,i}$ are stored in the fields of
\Cref{lem:packed-hamming} with the parameters
\begin{align*}
f&\defeq\text{the least integer with }2^{f}>10^{3}\log n,\\
k&\defeq\text{the largest integer with }2^{kf}\le n,\\
r&\defeq\text{the least integer with }rk\ge10^{3}\log n ,
\end{align*}
the bits of index above $10^{3}\log n$ being zero, so that
$2^{f}>10^{3}\log n\ge k$ as \Cref{lem:packed-hamming} requires, every
field integer is below $2^{kf}\le n$, and
\begin{equation}\label{eq:packed-params}
f\le\log\log n+11,\qquad
k\ge\frac{\log n}{f}-1,\qquad
r\le\frac{10^{3}\log n}{k}+1\le2\cdot10^{3}(\log\log n+11)+1
\end{equation}
for $\log n\ge2f$, which holds for $n\ge2^{40}$; for $2^{20}\le n<2^{40}$,
$k\ge1$ gives $r\le4\cdot10^{4}+1$.
In all cases $f=O(\log\log n)$ and $r=O(\log\log n)$.

The $2w-5$ scales of a vertex are arranged in a binary tree of pointers
of the same shape for every vertex, so that the bisection over the scales
of a query is a descent of two trees in lockstep without any
integer-addressed memory; the shape is defined on intervals of scale
indices.
\begin{definition}\label{def:packed-tree}
For integers $\ell<h$, $\mathrm{Tree}(\ell,h)$ is the empty tree if
$h=\ell+1$, and otherwise the binary tree whose root carries the index
$m$, the integer $(\ell+h)/2$ rounded down, with left subtree
$\mathrm{Tree}(\ell,m)$ and right subtree $\mathrm{Tree}(m,h)$.
\end{definition}
$\mathrm{Tree}(\ell,h)$ carries the indices $\ell+1,\dots,h-1$ once each,
has $h-\ell-1$ nodes, the subtree rooted at its node of index $m$ is
$\mathrm{Tree}(\ell',h')$ for an interval with $\ell'<m<h'$, and its depth
is at most $\log(h-\ell)+1$, since every level halves the interval length.

\Cref{fig:packed-state} lists the state of the oracle,
\Cref{fig:packed-pre} its preprocessing and \Cref{fig:packed-query} its
query; the proof of \Cref{lem:packed} follows the figures.

\begin{figure}[H]
\centering
\framebox{\begin{minipage}{0.95\linewidth}\small
\underline{Global variables}:
\begin{itemize}[leftmargin=*,itemsep=1pt]
\item the integers $c$, $w$ (the least integer with $2^{w}\ge2^{6}n^{c+1}$),
$f$, $k$, $r$ as above, and the constants $2^{f}$, $2^{(k-1)f}$.
\item one \emph{scale record} per scale index $i\in\{-w+6,\dots,w\}$,
holding $i$, the integer $2^{i+w}$ and the number $2^{2i+1}$, linked in
increasing order of $i$.
\item the shifts $\gamma_1,\dots,\gamma_{10^{3}\log n}$, multiples of $2^{-w}$
in $[0,2^{w+1})$, as a list.
\item for every vertex $v$ the numbers
$\YYbar_{v,1},\dots,\YYbar_{v,10^{3}\log n}$
as a list.
\item for every vertex $v$ and every scale index $i$ a
\emph{fingerprint record} holding the integer $\norm{\bb_{v,i}}_1$ and a
list of $r$ \emph{chunk records}, the $\ell$-th holding the two integers
$F_{v,i,\ell}$ and $F^{\mathrm{rev}}_{v,i,\ell}$.
\item for every vertex $v$ a tree $T_v$ of the shape
$\mathrm{Tree}(-w+6,w)$ whose node of index $i$ points to the scale
record of $i$ and to the fingerprint record of $(v,i)$; the vertex record
points to the root of $T_v$ and to the fingerprint record of
$(v,-w+6)$.
\end{itemize}
\underline{Invariants}:
\begin{itemize}[leftmargin=*,itemsep=1pt]
\item
$\YYbar_{v,j}=2^{-w}\lfloor2^{w}(\YY_{v,j}-\YY_{v_0,j})\rfloor
+\gamma_j+2^{w+1}$.
\item $(\bb_{v,i})_j=\lfloor\YYbar_{v,j}/2^{i}\rfloor\bmod2$ for
$j\le10^{3}\log n$ and $(\bb_{v,i})_j=0$ for $10^{3}\log n<j\le rk$.
\item $F_{v,i,\ell}=\sum_{j=1}^{k}(\bb_{v,i})_{(\ell-1)k+j}2^{(j-1)f}$ and
$F^{\mathrm{rev}}_{v,i,\ell}=\sum_{j=1}^{k}(\bb_{v,i})_{(\ell-1)k+j}2^{(k-j)f}$.
\end{itemize}
\end{minipage}}
\caption{State of the packed distance oracle.}
\label{fig:packed-state}
\end{figure}

\begin{figure}[H]
\centering
\framebox{\begin{minipage}{0.95\linewidth}\small
$\textsc{Preprocess}(\XX,c)$\\[2pt]
\underline{Input}: the rows $\XX_v$, lists of $10^{3}\log n$ numbers at the
vertex records, and the constant $c$;\\
writes all global variables of \Cref{fig:packed-state}.\\
\underline{Output}: none.
\begin{enumerate}[leftmargin=*,itemsep=1pt]
\item Set $w$ to the least integer with $2^{w}\ge2^{6}n^{c+1}$, $f$ to the
least integer with $2^{f}>10^{3}\log n$, $k$ to the largest integer with
$2^{kf}\le n$, $r$ to the least integer with $rk\ge10^{3}\log n$; compute
$2^{f}$, $2^{(k-1)f}$ and the scale records by repeated doubling.
\item Draw $\GG\in\R^{10^{3}\log n\times10^{3}\log n}$ with independent standard
Gaussian entries and compute $\YY_v\leftarrow\XX_v\GG$ for every $v$.
\item For $j=1,\dots,10^{3}\log n$ draw $\gamma_j$ uniformly from the
multiples of $2^{-w}$ in $[0,2^{w+1})$; for every $v$ and $j$ set
$\YYbar_{v,j}\leftarrow
2^{-w}\lfloor2^{w}(\YY_{v,j}-\YY_{v_0,j})\rfloor+\gamma_j+2^{w+1}$.
\item For every $v$ and every scale record $i$: create the fingerprint
record of $(v,i)$ with $\norm{\bb_{v,i}}_1\leftarrow0$ and $r$ chunk records
with $F\leftarrow0$, $F^{\mathrm{rev}}\leftarrow0$; then for $\ell=1,\dots,r$
and $j=(\ell-1)k+1,\dots,\ell k$ with $j\le10^{3}\log n$:
\begin{enumerate}[itemsep=1pt]
\item set $\beta$ to the remainder modulo $2$ of the quotient of the
integer $2^{w}\YYbar_{v,j}$ by $2^{i+w}$;
\item add $\beta$ to $\norm{\bb_{v,i}}_1$, add $\beta\,2^{(j-(\ell-1)k-1)f}$ to
$F_{v,i,\ell}$ and $\beta\,2^{(\ell k-j)f}$ to $F^{\mathrm{rev}}_{v,i,\ell}$,
the two powers of $2^{f}$ obtained from those of the previous $j$ by one
multiplication and one division by $2^{f}$.
\end{enumerate}
\item For every $v$ build $T_v$ of the shape $\mathrm{Tree}(-w+6,w)$ by the
recursion of \Cref{def:packed-tree}, pointing every node to its scale
record and fingerprint record; store at $v$ the root and the fingerprint
record of $(v,-w+6)$.
\end{enumerate}
\end{minipage}}
\caption{Preprocessing of the packed distance oracle.
Its guarantees are in \Cref{lem:packed}.}
\label{fig:packed-pre}
\end{figure}

\begin{figure}[H]
\centering
\framebox{\begin{minipage}{0.95\linewidth}\small
$D=\textsc{Query}(u,v)$\\[2pt]
\underline{Input}: pointers to the vertex records $u$ and $v$;\\
reads $f$, $k$, $r$, $2^{f}$, $2^{(k-1)f}$, the trees $T_u$, $T_v$, the scale
records and the fingerprint records.\\
\underline{Output}: the estimate $D$ of $\norm{\XX_u-\XX_v}_2^{2}$.
\begin{enumerate}[leftmargin=*,itemsep=1pt]
\item If $\norm{\bb_{u,-w+6}-\bb_{v,-w+6}}_1\le250\log n$, return
$D\leftarrow0$.
\item Set $h\leftarrow w$ and let the current nodes be the roots of $T_u$ and
$T_v$.
\item While the current nodes are not empty, with $m$ their common index:
\begin{enumerate}[itemsep=1pt]
\item if $\norm{\bb_{u,m}-\bb_{v,m}}_1\le250\log n$, set $h\leftarrow m$ and
move
both current nodes to their left children;
\item otherwise move both current nodes to their right children.
\end{enumerate}
\item Return $D\leftarrow2^{2h+1}$, read from the scale record of $h$.
\end{enumerate}
\end{minipage}}
\caption{Query of the packed distance oracle.
Every $\norm{\bb_{u,i}-\bb_{v,i}}_1$ is evaluated by the formula of
\Cref{lem:packed-hamming} from the fingerprint records of $(u,i)$ and
$(v,i)$, their chunk lists traversed in lockstep.
Its guarantees are in \Cref{lem:packed}.}
\label{fig:packed-query}
\end{figure}

\begin{proof}[Proof of \Cref{lem:packed}]
\emph{The good event.}
Let (i) be the event that
$|\YY_{v,j}-\YY_{v_0,j}|\le n^{c+1}$ for all $v$ and $j$, and let (ii) be
the event that for every pair $(u,v)$ with $\XX_u\neq\XX_v$ and every scale
index $i\in\{-w+6,\dots,w\}$,
\begin{align*}
\norm{\bb_{u,i}-\bb_{v,i}}_1&>250\log n
&&\text{if }2^{i}\le\norm{\XX_u-\XX_v}_2,\\
\norm{\bb_{u,i}-\bb_{v,i}}_1&\le250\log n
&&\text{if }2^{i}\ge10\norm{\XX_u-\XX_v}_2 .
\end{align*}
By \Cref{lem:packed-law}, $\YY_{v,j}-\YY_{v_0,j}$ is
$N(0,\norm{\XX_v-\XX_{v_0}}_2^{2})$
with $\norm{\XX_v-\XX_{v_0}}_2\le n^{c}$, and a Gaussian $g$ of variance at
most $n^{2c}$ has, with $\lambda=n^{1-c}$,
\[
\Pr[g\ge n^{c+1}]\ \le\ e^{-\lambda n^{c+1}}\Ex[e^{\lambda g}]
\ \le\ e^{-n^{2}+n^{2}/2}\ =\ e^{-n^{2}/2};
\]
a union bound over the $2n\cdot10^{3}\log n$ one-sided events gives
$\Pr[\text{not (i)}]\le2\cdot10^{3}n\log n\,e^{-n^{2}/2}\le n^{-40}$.

\emph{Event (ii).}
Fix a pair $(u,v)$ with $\XX_u\neq\XX_v$ and a scale index $i$.
In units of $2^{-w}$, by \Cref{eq:packed-shift,eq:packed-bits},
\[
(\bb_{u,i})_j=\left\lfloor\frac{\lfloor2^{w}(\YY_{u,j}-\YY_{v_0,j})\rfloor
+2^{2w+1}+2^{w}\gamma_j}{2^{i+w}}\right\rfloor\bmod2,
\]
and likewise for $v$, where $2^{w}\gamma_j$ is uniform on the integers of
$[0,2^{2w+1})=[0,2\cdot2^{i+w}\cdot2^{w-i})$ with $w-i\ge0$, and is independent
of $\GG$.
Conditionally on $\GG$, \Cref{lem:packed-parity} with $T=2^{i+w}$ and
$M=2^{w-i}$ gives that the $j$-th bits of $\bb_{u,i}$ and $\bb_{v,i}$
disagree with probability
\begin{equation}\label{eq:packed-disagree}
\frac{\dist\bigl(2^{-w}\lfloor2^{w}(\YY_{v,j}-\YY_{v_0,j})\rfloor
-2^{-w}\lfloor2^{w}(\YY_{u,j}-\YY_{v_0,j})\rfloor,\ 2^{i+1}\Z\bigr)}{2^{i}} ,
\end{equation}
and the $10^{3}\log n$ disagreement indicators are independent, since the
$j$-th depends on $(\GG,\gamma)$ only through the $j$-th column of $\GG$
and $\gamma_j$ (\Cref{lem:packed-law}).
The argument of $\dist$ in \Cref{eq:packed-disagree} differs from
$\YY_{v,j}-\YY_{u,j}$ by less than $2^{-w}\le2^{i}/2^{6}$, as $i\ge-w+6$, so by
the Lipschitz bound of \Cref{lem:packed-wave} the expectation of
\Cref{eq:packed-disagree} over $\GG$ differs from
$\Ex[\dist(x,2^{i+1}\Z)/2^{i}]$ for $x\sim N(0,\norm{\XX_u-\XX_v}_2^{2})$
(\Cref{lem:packed-law}) by at most $2^{-6}\le0.016$.
Hence the probability that the $j$-th bits disagree is at least
$0.49-0.016\ge0.47$ when $2^{i}\le\norm{\XX_u-\XX_v}_2$, and at most
$\frac1{10}\sqrt{2/\pi}+0.016\le0.1$ when $2^{i}\ge10\norm{\XX_u-\XX_v}_2$, by
\Cref{lem:packed-wave}.
\Cref{lem:packed-hoeffding} with $N=10^{3}\log n$ gives, in the first case,
\[
\Pr\bigl[\norm{\bb_{u,i}-\bb_{v,i}}_1\le250\log n\bigr]
\le\Pr\Bigl[\sum_j\xi_j\le\Ex\Bigl[\sum_j\xi_j\Bigr]-0.22\cdot10^{3}\log n\Bigr]
\le e^{-2\cdot0.22^{2}\cdot10^{3}\log n}\le n^{-90},
\]
and in the second case
\[
\Pr\bigl[\norm{\bb_{u,i}-\bb_{v,i}}_1>250\log n\bigr]
\le\Pr\Bigl[\sum_j\xi_j\ge\Ex\Bigl[\sum_j\xi_j\Bigr]+0.15\cdot10^{3}\log n\Bigr]
\le e^{-2\cdot0.15^{2}\cdot10^{3}\log n}\le n^{-45},
\]
where $\xi_j$ is the indicator that the $j$-th bits disagree and
$e^{-\log n}\le n^{-1}$.
A union bound over the at most $n^{2}$ pairs and the
$2w-5\le2(c+1)\log n+9\le n^{3}$
scales (as $c$ is a constant and $n\ge2^{20}$) gives
$\Pr[\text{not (ii)}]\le n^{5}\cdot n^{-45}=n^{-40}$.
Both events hold with probability at least $1-2n^{-40}\ge1-n^{-14}$; the
rest of the proof is on these events.

\emph{Correctness.}
The fingerprint records hold the integers of \Cref{lem:packed-hamming}
for $\uu=\bb_{u,i}$ and $\vv=\bb_{v,i}$ (invariants of
\Cref{fig:packed-state}), so every evaluation of
$\norm{\bb_{u,i}-\bb_{v,i}}_1$ in \Cref{fig:packed-query} is exact.
If $\XX_u=\XX_v$ then $\YY_u=\YY_v$, so $\bb_{u,i}=\bb_{v,i}$ for every $i$,
Step~1 returns $D=0=\norm{\XX_u-\XX_v}_2^{2}$.
Otherwise $n^{-c}\le\norm{\XX_u-\XX_v}_2\le n^{c}$, so
$2^{-w+6}\le2^{6}/(2^{6}n^{c+1})<n^{-c}\le\norm{\XX_u-\XX_v}_2$ and
$2^{w}\ge2^{6}n^{c+1}\ge10\norm{\XX_u-\XX_v}_2$; by (ii),
$\norm{\bb_{u,-w+6}-\bb_{v,-w+6}}_1>250\log n$, so Step~1 does not return, and
$\norm{\bb_{u,w}-\bb_{v,w}}_1\le250\log n$.
The descent of Step~3 maintains the invariant that the current nodes of
$T_u$ and $T_v$ are the roots of the subtrees $\mathrm{Tree}(\ell,h)$ of both
trees for the same interval $(\ell,h)$, with
$\norm{\bb_{u,\ell}-\bb_{v,\ell}}_1>250\log n$ and
$\norm{\bb_{u,h}-\bb_{v,h}}_1\le250\log n$: it holds at the start with
$(\ell,h)=(-w+6,w)$ by the two facts just established, the two trees have
the same shape so the current nodes carry the same index $m$ and their
left (right) children are the roots of $\mathrm{Tree}(\ell,m)$
($\mathrm{Tree}(m,h)$) in both trees, and Step~3(a) or~3(b) moves to the
child whose interval keeps the invariant.
When the current nodes are empty, $h=\ell+1$ with
$\norm{\bb_{u,\ell}-\bb_{v,\ell}}_1>250\log n$ and
$\norm{\bb_{u,h}-\bb_{v,h}}_1\le250\log n$; by (ii) the first forces
$2^{\ell}<10\norm{\XX_u-\XX_v}_2$ and the second forces
$2^{h}>\norm{\XX_u-\XX_v}_2$, so
$\norm{\XX_u-\XX_v}_2<2^{h}=2\cdot2^{\ell}<20\norm{\XX_u-\XX_v}_2$ and
\[
\norm{\XX_u-\XX_v}_2^{2}\ <\ 2\norm{\XX_u-\XX_v}_2^{2}\ <\ D=2^{2h+1}
\ <\ 2\cdot20^{2}\norm{\XX_u-\XX_v}_2^{2}\ \le\ 10^{3}\norm{\XX_u-\XX_v}_2^{2}.
\]
The query reads only stored records and draws no randomness, so $D(u,v)$
is a function of the preprocessing, and the bounds hold for all pairs
simultaneously on the events (i) and (ii).

\emph{Sizes of the integers.}
On (i),
$|2^{-w}\lfloor2^{w}(\YY_{v,j}-\YY_{v_0,j})\rfloor|\le n^{c+1}+2^{-w}<2^{w}$,
so $0<\YYbar_{v,j}<2^{w+3}$ and the integer $2^{w}\YYbar_{v,j}$ is below
$2^{2w+3}$; the scale records hold $2^{i+w}\le2^{2w}$ and
$2^{2i+1}=2^{-2w}\cdot2^{2(i+w)+1}$, whose numerator is below $2^{4w+1}$; and
the integers of \Cref{lem:packed-hamming} are below $2^{2kf}\le n^{2}$.
Since $2^{w}<2^{7}n^{c+1}$ and $n\ge2^{20}$, all integers are below
$2^{29}n^{4c+4}\le n^{4c+7}$, within the $(4c+7)\log n$ bits of the model.

\emph{Time and space.}
Step~1 of \Cref{fig:packed-pre} takes $O(w+kf)=O(\log n)$ operations and
creates $2w-5$ scale records; Step~2 takes $(10^{3}\log n)^{2}$ draws and
$n(10^{3}\log n)^{2}$ operations; Step~3 takes $10^{3}\log n$ draws and
$O(n\log n)$ operations.
Step~4 takes, per vertex and scale, $r+1$ record creations and $O(1)$
operations per $j\le10^{3}\log n$, that is
$O(n(2w-5)(10^{3}\log n+r))=O(n\log^{2}n)$ operations and $n(2w-5)(r+1)$
records in all, which is $O(n\log n\cdot\log\log n)=\Odouble(n\log n)$ by
\Cref{eq:packed-params}.
Step~5 creates $n(2w-7)$ nodes in $O(n\log n)$ operations.
The total is $O(n\log^{2}n)$ operations and $\Odouble(n\log n)$ records.
A query evaluates $\norm{\bb_{u,i}-\bb_{v,i}}_1$ once in Step~1 and once per
level of the descent, at most $\log(2w-5)+1=O(\log\log n)$ times, each
evaluation costing at most $4r+3=O(\log\log n)$ operations by
\Cref{lem:packed-hamming} and \Cref{eq:packed-params}, plus the comparison
with $250\log n$; the moves between nodes cost $O(1)$ each.
The total is $O((\log\log n)^{2})=\Odouble(1)$ operations, in the worst case.
\end{proof}

The multiplication of \Cref{lem:packed-hamming} is the only place where
the instruction set matters; the following lemma records the two
replacements available on a word RAM, so that \Cref{lem:packed} holds
with the same bounds under either of the common instruction sets.

\begin{lemma}\label{lem:packed-model}
Let the state of \Cref{fig:packed-state} be built by \Cref{fig:packed-pre}.
\begin{enumerate}[label=(\alph*),leftmargin=*]
\item If the bitwise exclusive or $\oplus$ and the population count
$\textsc{Popcount}$ (the number of ones in the binary expansion) of
integers below $n$ cost $O(1)$, then
\[
\norm{\bb_{u,i}-\bb_{v,i}}_1
=\sum_{\ell=1}^{r}\textsc{Popcount}(F_{u,i,\ell}\oplus F_{v,i,\ell}),
\]
evaluated in $3r$ operations; the integers $F^{\mathrm{rev}}$ and
$\norm{\bb_{v,i}}_1$ are then not needed.
\item If the memory is integer-addressed and $\oplus$ costs $O(1)$ but
$\textsc{Popcount}$ is not available, a table of the population counts of
the integers below $2^{s}$, $s$ the least integer with $2s\ge kf$, has
at most $2\sqrt n$ entries, is built in $O(\sqrt n)$ operations, and gives
\[
\textsc{Popcount}(x)=\textsc{Popcount}(\lfloor x/2^{s}\rfloor)
+\textsc{Popcount}(x\bmod2^{s})
\]
for $x<2^{kf}$ by two lookups, one division, one remainder and one
addition.
\end{enumerate}
In both cases the preprocessing remains $O(n\log^{2}n)$ operations and the
query $O((\log\log n)^{2})=\Odouble(1)$ operations, and \Cref{lem:packed} holds
with the same bounds.
\end{lemma}
\begin{proof}
(a) The integer $F_{v,i,\ell}$ has the bit $(\bb_{v,i})_{(\ell-1)k+j}$ at
binary position $(j-1)f$ and zeros elsewhere (invariants of
\Cref{fig:packed-state}); the exclusive or acts position by position, so
$F_{u,i,\ell}\oplus F_{v,i,\ell}$ has at position $(j-1)f$ the bit
$|(\bb_{u,i})_{(\ell-1)k+j}-(\bb_{v,i})_{(\ell-1)k+j}|$ and zeros elsewhere,
and its population count is the number of disagreements in the $\ell$-th
block; summing over the $r$ blocks gives $\norm{\bb_{u,i}-\bb_{v,i}}_1$ with
one exclusive or, one population count and one addition per block.
The integers are below $2^{kf}\le n$.

(b) The table is filled in increasing order by
$\textsc{Popcount}(0)=0$ and
$\textsc{Popcount}(x)=\textsc{Popcount}(\lfloor x/2\rfloor)+(x\bmod2)$, three
operations per entry; $2^{s}<2\cdot2^{kf/2}\le2\sqrt n$ as $2(s-1)<kf$ and
$2^{kf}\le n$.
The identity splits the binary expansion of $x<2^{kf}\le2^{2s}$ into its
$s$ low bits and its at most $s$ high bits, both below $2^{s}$.

In both cases every evaluation of $\norm{\bb_{u,i}-\bb_{v,i}}_1$ costs
$O(r)=O(\log\log n)$ operations, as in \Cref{lem:packed-hamming}, and the
correctness argument of \Cref{lem:packed} is unchanged, as it uses only
that the evaluations are exact; the table adds $O(\sqrt n)$ to the
preprocessing.
\end{proof}

\clearpage
\section{The fixed solver}\label{sec:solver}

The random vertex vectors of a refresh are
$\ZZ(\LLtil)\BBtil^{\top}\WWtil^{1/2}\gg$ for a Gaussian $\gg$; for them to be
Gaussian with covariance
$\ZZ(\LLtil)\LLtil\ZZ(\LLtil)$, the solve must be one fixed linear operator,
applied consistently to every right-hand side, and for the trace
identities of the refresh this operator must be symmetric.
The solver of Kyng, Lee, Peng, Sachdeva and Spielman~\cite{KLPSS16:conf}
builds such an operator; we state what this paper uses and derive it from
the statements of the arXiv version (arXiv:1512.01892), quoted verbatim
below with their numbers.

\begin{lemma}[from Theorem~1.2, Claim~2.4, Lemma~3.3 and Lemmas~I.3
and~I.5 of~\cite{KLPSS16:conf}]\label{lem:chain}
There is a routine $\textsc{LsolverKLPSS}$ that, on an SDDM matrix $\MM$
with $n$ rows, $m'$ nonzero entries and polynomially bounded entries, and
an error $\epsilon\in(0,\frac12]$, returns
$\ZZ(\MM)\leftarrow\textsc{LsolverKLPSS}(\MM,\epsilon)$, implicit access to
a symmetric matrix $\ZZ(\MM)$, such that with probability $1-n^{-c}$, for
the absolute constant $c>0$ of the phrase ``with high probability'' in
Lemma~I.5 of~\cite{KLPSS16:conf}:
\begin{enumerate}[label=(\alph*),leftmargin=*]
\item $\ZZ(\MM)\approx_{\epsilon}\MM^{\dag}$ on the range of $\MM$, that
is, $\xx^{\top}\ZZ(\MM)\xx\approx_{\epsilon}\xx^{\top}\MM^{\dag}\xx$ for
every $\xx$ in the range of $\MM$;
\item the construction takes $O(m'\log n)+O(n\log^{2}n)$ time;
\item one application of $\ZZ(\MM)$ to a vector takes
$O\bigl((m'+n\log n)\log(1/\epsilon)\bigr)$ time.
\end{enumerate}
The constants in (b) and (c) are $2^{O(k\log^{2}k)}$ for the phase length
$k$ of the recursive construction of~\cite{KLPSS16:conf}, which we fix at
$k=28$; see the accounting after the proof.
\end{lemma}

\paragraph{The source statements.}
In~\cite{KLPSS16:conf} a bDD matrix is a Hermitian block matrix
$\MM$ with $r\times r$ blocks satisfying
$\MM_{[i,i]}\succeq\sum_{j\ne i}\lVert\MM_{[i,j]}\rVert\cdot\II_r$ for every
block row $i$ (its Definition~3.4 with $\alpha=0$, ``for $r=1$, Laplacian
matrices are $0$-bDD''), and the approximation relation is multiplicative
in $e^{\pm\epsilon}$:
\begin{quote}
``We say that $\AA$ is an $\epsilon$-approximation of $\BB$, written
$\AA\approx_{\epsilon}\BB$, if
$e^{\epsilon}\BB\succeq\AA\succeq e^{-\epsilon}\BB$.
This relation is symmetric.''
(Section~2 of~\cite{KLPSS16:conf}.)
\end{quote}
The statements used are the following.
\begin{quote}
``\textbf{Theorem 1.2} (Sparsified Multigrid).
There is an algorithm that, when given a bDD matrix $\MM$ with $n$ block
rows and $m$ nonzero blocks, produces a solver for $\MM$ in
$O(m\log n+n\log^{2+o(1)}n)$ work and $O(n^{o(1)})$ depth so that the
solver finds $\epsilon$-approximate solutions to systems of equations in
$\MM$ in $O((m+n\log^{1+o(1)}n)\log(1/\epsilon))$ work and
$O(\log^{2}n\log\log n\log(1/\epsilon))$ depth.''

``\textbf{Claim 2.4.}
Suppose that all nonzero eigenvalues of $\MM$ are at least $\mu$ and
$\ZZ\approx_{\epsilon}(\MM+\epsilon\mu\II)^{-1}$ for some $0<\epsilon<1/2$.
Given any $\bb$ such that $\MM\xx=\bb$ for some $\xx$, we have
$\lVert\xx-\ZZ\bb\rVert_{\MM}^{2}\le6\epsilon\lVert\xx\rVert_{\MM}^{2}$.''

``\textbf{Lemma 3.3.}
Consider an $\epsilon$-SCC for $\MM^{(0)}$ where $\MM^{(i)}$ and $\ZZ^{(i)}$
can be applied to a vector in work $W_{\MM^{(i)}}$, $W_{\ZZ^{(i)}}$ and depth
$D_{\MM^{(i)}}$, $D_{\ZZ^{(i)}}$ respectively.
The algorithm
$\textsc{ApplyChain}\bigl((\MM^{(1)},\ZZ^{(1)}),\dots,(\MM^{(d)},\ZZ^{(d)});
F_1,\dots,F_d;\bb\bigr)$ corresponds to a linear operator $\WW$ acting on
$\bb$ such that
1. $\WW^{-1}\approx_{\sum_{i=1}^{d}2\epsilon_i}\MM^{(0)}$, and
2. for any vector $\bb$, $\textsc{ApplyChain}(\dots;\bb)$ runs in
$O(\sum_{i=1}^{d}(D_{\MM^{(i)}}+D_{\ZZ^{(i)}}))$ depth and
$O(\sum_{i=1}^{d}(W_{\MM^{(i)}}+W_{\ZZ^{(i)}}))$ work.''

``\textbf{Lemma 3.5.}
Given a bDD matrix $\MM$ with $n$ block-rows, and an $\alpha\ge0$,
\textsc{bDDSubset} computes a subset $F$ of size least $n/(8(1+\alpha))$
such that $\MM_{[F,F]}$ is $\alpha$-bDD.
It runs in runs in $O(m)$ expected work and $O(\log n)$ expected depth,
where $m$ is the number of nonzero blocks in $\MM$.''

``\textbf{Lemma I.3.}
For any bDD matrix $\MM^{(0)}=\xi\II+\BB^{(0)}(\BB^{(0)})^{*}$ for some
$\xi>0$ that has $n$ blocks, the algorithm
$\textsc{RecursiveConstruct}(\MM^{(0)})$ returns a Schur Complement Chain
$(\MM^{(1)},\MM^{(2)},\dots;F_1,F_2,\dots)$ whose corresponding the linear
operator $\WW$ satisfies $\WW\approx_{O(1)}(\MM^{(0)})^{-1}$.
Also, we can evaluate $\WW\bb$ in $O(\log^{2}n\log\log n)$ depth and
$2^{O(k\log^{2}k)}n\log n$ work for any vector $\bb$.''

``\textbf{Lemma I.5.}
With high probability, $\textsc{RecursiveConstruct}(\MM^{(0)})$ takes
$2^{O(\log n/k)}$ depth and $m\log n+2^{O(k\log^{2}k)}n\log^{2}n$ work.''
\end{quote}
Two displayed formulas of the proofs are also used.
The proof of Theorem~1.3 in Appendix~C of~\cite{KLPSS16:conf} writes the
operator of \textsc{ApplyChain} as
\begin{quote}
``$(\WW^{(1)})^{-1}=\UU^{\top}\operatorname{diag}\bigl((\ZZ^{(1)})^{-1},\dots,
(\ZZ^{(d-1)})^{-1},\II_{|C_{d-1}|}\bigr)\UU$''
\end{quote}
for an upper triangular $\UU$ (the displayed factorization on page~26 of
arXiv:1512.01892, with $\UU$ assembled from the identity and the blocks
$\ZZ^{(i)}\MM_{[F_i,C_i]}$), and equation~(7) of Appendix~E defines the
Jacobi operators as
\begin{quote}
``$\ZZ^{(k)}\defeq\sum_{i=0}^{k}\XX_{[F,F]}^{-1}
\bigl(-\LL_{[F,F]}\XX_{[F,F]}^{-1}\bigr)^{i}$''
\end{quote}
for the splitting $\MM_{[F,F]}=\XX_{[F,F]}+\LL_{[F,F]}$ of Lemma~3.6 there,
with $\XX_{[F,F]}$ block diagonal.
In the source the proof of Theorem~1.2 chooses $k=\log\log\log n$, so that
its $2^{O(k\log^{2}k)}$ is $\log^{o(1)}n$ and its depth is $n^{o(1)}$; this
is the only place where $k$ is fixed, and only the depth bound needs $k$ to
grow.

\begin{proof}
Let $\MM$ be an SDDM matrix: symmetric, $\MM_{u,v}\le0$ for $u\ne v$, and
$\MM_{u,u}\ge\sum_{v\ne u}|\MM_{u,v}|$ for every $u$.
Then
\begin{equation}\label{eq:sddm-gram}
\MM=\sum_{u<v}(-\MM_{u,v})(\oone_u-\oone_v)(\oone_u-\oone_v)^{\top}
+\sum_{u}\Bigl(\MM_{u,u}-\sum_{v\ne u}|\MM_{u,v}|\Bigr)\oone_u\oone_u^{\top}
\succeq0,
\end{equation}
and $\MM$ is a bDD matrix with $r=1$.

\emph{The smallest nonzero eigenvalue.}
Let $\xx$ be a unit vector orthogonal to the kernel of $\MM$, and let $H$
be the graph on $V$ with an edge $uv$ for every nonzero $\MM_{u,v}$,
$u\ne v$.
The quadratic form \Cref{eq:sddm-gram} splits over the connected
components of $H$.
\begin{itemize}[leftmargin=*]
\item If a component $C$ contains a vertex $v$ with
$\MM_{v,v}>\sum_{u\ne v}|\MM_{v,u}|$, then for every $u\in C$ a path from
$u$ to $v$ in $H$ has at most $n-1$ edges, so by the Cauchy--Schwarz
inequality
$\xx_u^{2}\le2(\xx_u-\xx_v)^{2}+2\xx_v^{2}
\le2n\sum_{ab\text{ on the path}}(\xx_a-\xx_b)^{2}+2\xx_v^{2}$, and
summing over $u\in C$,
$\sum_{u\in C}\xx_u^{2}\le2n^{2}\cdot
\xx_C^{\top}\MM\xx_C/\min\{|\MM_{a,b}|,\MM_{v,v}-\sum_{u\ne v}|\MM_{v,u}|\}$,
the minimum over the nonzero off-diagonal entries and nonzero excesses
$\MM_{v,v}-\sum_{u\ne v}|\MM_{v,u}|$ of $\MM$.
\item Otherwise $\MM$ restricted to $C$ is the Laplacian of a connected
graph, $\oone_C$ is in the kernel of $\MM$, so $\sum_{u\in C}\xx_u=0$,
$\max_{u\in C}\xx_u-\min_{u\in C}\xx_u\ge\max_{u\in C}|\xx_u|
\ge(\sum_{u\in C}\xx_u^{2})^{1/2}/\sqrt{n}$, and along a path of at most
$n-1$ edges between the maximizer and the minimizer the Cauchy--Schwarz
inequality gives
$\xx_C^{\top}\MM\xx_C\ge\min_{a\ne b,\ \MM_{a,b}\ne0}|\MM_{a,b}|
\cdot\sum_{u\in C}\xx_u^{2}/n^{2}$.
\end{itemize}
Summing over the components, every nonzero eigenvalue of $\MM$ is at least
$n^{-3}$ times the least nonzero off-diagonal entry or excess of $\MM$,
and for polynomially bounded entries this is $n^{-O(1)}$; write $\mu$ for
this lower bound, a quantity computed from $\MM$ in $O(m'+n)$ time.

\emph{The shift.}
Put $\xi\defeq\epsilon\mu/10$.
By \Cref{eq:sddm-gram}, $\MM+\xi\II=\xi\II+\BB^{(0)}(\BB^{(0)})^{\top}$ with
$\BB^{(0)}$ the matrix whose columns are
$(-\MM_{u,v})^{1/2}(\oone_u-\oone_v)$ over the edges of $H$ and
$(\MM_{u,u}-\sum_{v\ne u}|\MM_{u,v}|)^{1/2}\oone_u$ over the vertices with
positive excess, which is the input form of Lemmas~I.3 and~I.5.
Run $\textsc{RecursiveConstruct}(\MM+\xi\II)$ with phase length $k=28$; by
Lemma~I.5 this takes $O(m'\log n)+O(n\log^{2}n)$ time, and by Lemma~I.3 it
returns a chain whose operator $\WW$ satisfies, for an absolute constant
$a>0$,
\begin{equation}\label{eq:chain-W}
e^{-a}(\MM+\xi\II)^{-1}\preceq\WW\preceq e^{a}(\MM+\xi\II)^{-1},
\end{equation}
and one application of $\WW$ takes $O(n\log n)$ time, where the constants
hidden in the two $O(\cdot)$ are $2^{O(k\log^{2}k)}$ with $k=28$.

\emph{$\WW$ is a fixed symmetric matrix.}
By Lemma~3.3, \textsc{ApplyChain} is a linear operator $\WW$ in $\bb$; the
chain is fixed once built, so every application uses the same $\WW$.
In the displayed factorization
$(\WW^{(1)})^{-1}=\UU^{\top}\operatorname{diag}((\ZZ^{(1)})^{-1},\dots,
(\ZZ^{(d-1)})^{-1},\II)\UU$ of Appendix~C, each $\ZZ^{(i)}$ is the Jacobi
operator of equation~(7),
$\sum_{i=0}^{k}\XX_{[F,F]}^{-1}(-\LL_{[F,F]}\XX_{[F,F]}^{-1})^{i}
=\XX_{[F,F]}^{-1/2}\sum_{i=0}^{k}(-\XX_{[F,F]}^{-1/2}\LL_{[F,F]}
\XX_{[F,F]}^{-1/2})^{i}\XX_{[F,F]}^{-1/2}$, a symmetric matrix since
$\XX_{[F,F]}$ is a positive diagonal matrix and $\LL_{[F,F]}$ is symmetric
(for $r=1$ all matrices are real).
Hence $(\WW^{(1)})^{-1}$ and $\WW=\WW^{(1)}$ are symmetric.

\emph{From constant to $\epsilon$ accuracy.}
Let $T$ be the least integer with $(1-e^{-2a})^{T}\le\epsilon/10$, so
$T\le e^{2a}\ln(10/\epsilon)+1=O(\log(1/\epsilon))$, and define
\[
\ZZ(\MM)\defeq\sum_{i=0}^{T-1}e^{-a}
\bigl(\II-e^{-a}\WW(\MM+\xi\II)\bigr)^{i}\WW .
\]
Since $\WW$ is symmetric positive definite by \Cref{eq:chain-W},
$(\II-e^{-a}\WW(\MM+\xi\II))^{i}\WW
=\WW^{1/2}(\II-e^{-a}\WW^{1/2}(\MM+\xi\II)\WW^{1/2})^{i}\WW^{1/2}$, so
\[
\ZZ(\MM)=\WW^{1/2}\,f\bigl(\WW^{1/2}(\MM+\xi\II)\WW^{1/2}\bigr)\,\WW^{1/2},
\qquad
f(\lambda)\defeq\sum_{i=0}^{T-1}e^{-a}(1-e^{-a}\lambda)^{i}
=\frac{1-(1-e^{-a}\lambda)^{T}}{\lambda},
\]
a symmetric matrix.
By \Cref{eq:chain-W}, the eigenvalues $\lambda$ of
$\WW^{1/2}(\MM+\xi\II)\WW^{1/2}$ lie in $[e^{-a},e^{a}]$, so
$0\le1-e^{-a}\lambda\le1-e^{-2a}$ and
$(1-\epsilon/10)/\lambda\le f(\lambda)\le1/\lambda$; therefore
\[
(1-\epsilon/10)\,\WW^{1/2}\bigl(\WW^{1/2}(\MM+\xi\II)\WW^{1/2}\bigr)^{-1}
\WW^{1/2}\preceq\ZZ(\MM)\preceq
\WW^{1/2}\bigl(\WW^{1/2}(\MM+\xi\II)\WW^{1/2}\bigr)^{-1}\WW^{1/2},
\]
and both outer matrices equal $(\MM+\xi\II)^{-1}$; that is,
$\ZZ(\MM)\approx_{\epsilon/10}(\MM+\xi\II)^{-1}$.
For $\xx$ in the range of $\MM$, expanding $\xx$ in the eigenvectors of
$\MM$ with eigenvalues $\lambda_i\ge\mu\ge10\xi/\epsilon$,
\[
\xx^{\top}(\MM+\xi\II)^{-1}\xx=\sum_i\frac{(\xx^{\top}\vv_i)^{2}}{\lambda_i+\xi}
\in\Bigl[\frac{1}{1+\epsilon/10},1\Bigr]\cdot\xx^{\top}\MM^{\dag}\xx ,
\]
so
$\xx^{\top}\ZZ(\MM)\xx\in[(1-\epsilon/10)^{2},1]\cdot\xx^{\top}\MM^{\dag}\xx$,
which is part (a).
Each of the $T$ terms is accumulated with one application of $\WW$
($O(n\log n)$ time) and one multiplication by $\MM+\xi\II$ ($O(m'+n)$
time), which is part (c); part (b) is the construction time of Lemma~I.5
plus the $O(m'+n)$ time to compute $\mu$ and $\xi$.
The probability is that of Lemma~I.5, whose phrase ``with high
probability'' the source does not quantify.
\end{proof}

\paragraph{Exponent accounting for \Cref{lem:chain}.}
Every logarithmic factor of \Cref{lem:chain} is traced in
\Cref{tab:chain-accounting}; the source's $\log^{o(1)}n$ factors come from
its choice $k=\log\log\log n$ and are replaced here by the constant
$2^{O(k\log^{2}k)}$, $k=28$, as explained in the table.
\begin{table}[H]\centering\footnotesize
\begin{tabular}{@{}p{0.17\linewidth}p{0.35\linewidth}p{0.12\linewidth}%
p{0.05\linewidth}p{0.2\linewidth}@{}}
\toprule
factor & lemma or step producing it & logarithm of & exp. &
explicit or derived\\
\midrule
build, $m'\log n$ & Lemma~I.5, the term $m\log n$ (the recurrence
$W_{\mathrm{dense}}(m)\le W_{\mathrm{dense}}(m/2)+m\log n+\dots$ in its
proof), inherited from the $m\log n$ work of \textsc{Sparsify}
(Lemma~H.1) & $n$ & $1$ & explicit in Lemmas~I.5, H.1\\
build, $n\log^{2}n$ & Lemma~I.5, the term $2^{O(k\log^{2}k)}n\log^{2}n$;
in Lemma~I.4 one $\log n$ is the $O(Kn\log n/\epsilon^{2})$ density of
the output of \textsc{Sparsify} (Lemma~H.1), the other the
$O(\log n)$ linear solves per sparsification
(``$w_{\mathrm{eval}}\log n$'' in Lemma~H.1) & $n$, $n$ & $2$ & explicit
in Lemmas~I.4, I.5\\
apply, $n\log n$ & Lemma~I.3, the term $2^{O(k\log^{2}k)}n\log n$; the
$\log n$ is the density $n^{(i)}\log n^{(i)}$ of the chain matrices
(Lemma~I.2, item~2), summed geometrically over the levels & $n$ & $1$ &
explicit in Lemmas~I.2, I.3\\
apply, $\log(1/\epsilon)$ & the $T=O(\log(1/\epsilon))$ steps of the
iteration in the proof above, from $\WW\approx_{O(1)}(\MM+\xi\II)^{-1}$
(Lemma~I.3) to $\epsilon$; the source's Theorem~1.2 states the same
factor & $1/\epsilon$ & $1$ & derived (our iteration); explicit in
Theorem~1.2\\
apply, $m'\log(1/\epsilon)$ & one multiplication by $\MM+\xi\II$ per step
of the iteration; Theorem~1.2 states ``$(m+n\log^{1+o(1)}n)\log(1/\epsilon)$''
& $1/\epsilon$ & $1$ & explicit in Theorem~1.2\\
constant $2^{O(k\log^{2}k)}$ & Lemma~I.2 (densities $2^{3ck\log^{2}(jk)}$),
Lemmas~I.3, I.5; the source fixes $k=\log\log\log n$ in the proof of
Theorem~1.2 for the depth $2^{O(\log n/k)}=n^{o(1)}$, which turns the
constant into $\log^{o(1)}n$; the work sums of Lemmas~I.3 and~I.5
converge for every $k$ with $2\beta^{k}<1$, where $\beta=1-1/40$ is the
vertex reduction of Lemma~3.5 with $\alpha=4$, so $k=28$ makes the
constant absolute & none & $0$ & derived from the proofs of Lemmas~I.3,
I.5; not stated in the source\\
no $\log n$ from $\epsilon$ or $\xi$ & the construction does not depend on
$\epsilon$ or on the shift $\xi$ (Lemma~I.5 has no such term); the
$\log(1/\xi)=O(\log n)$ bits of the shift only enter the word size & --- &
$0$ & explicit\\
probability $n^{-c}$ & Lemma~I.5 ``with high probability'', inherited from
Lemma~H.1 and Theorem~H.2 (Theorem~3 of~\cite{CLMMPS15:conf}); the
exponent is set by the $O(\log n)$ rows of the random projection in the
proof of Lemma~H.1 & $n$ & $c$ & exponent not stated in the source\\
\bottomrule
\end{tabular}
\caption{Exponent accounting for \Cref{lem:chain}; ``explicit'' means the
factor is written in the quoted statement, ``derived'' that it was traced
by us through the proof named.}\label{tab:chain-accounting}
\end{table}

\clearpage
\section{Table of symbols}\label{app:symbols}

The symbols of the paper in order of first appearance.
\begin{center}\small
\begin{tabular}{@{}p{0.2\linewidth}p{0.75\linewidth}@{}}
\toprule
symbol & meaning (where defined)\\
\midrule
$\Odouble(\cdot)$ & $O(\cdot)$ up to $(\log\log n)^{O(1)}$ factors (\Cref{sec:main})\\
$G=(V,E,\ww)$, $n$ & the current graph on the fixed vertex set $V$, $|V|=n\ge2^{20}$, weights $\ww_e\in[n^{-10},n^{10}]$ (\Cref{sec:main})\\
$\LL(G)$, $\LL(G)^{\dag}$ & Laplacian and its pseudoinverse (\Cref{sec:prelim})\\
$\epsilon$, $\epsilon_k$ & accuracy of a query, $n^{-10}\le\epsilon\le\frac12$; of the $k$-th query (\Cref{sec:main})\\
$\AA\approx_{\epsilon}\BB$ & $(1-\epsilon)\BB\preceq\AA\preceq(1+\epsilon)\BB$ (\Cref{sec:prelim})\\
$H$ & the sparsifier returned by a query (\Cref{sec:main})\\
$G_k$, $\Gret_k$, $\Gnew_k$ & the graph at the $k$-th query; its retained part (edges of $G_{k-1}$ at the smaller of the two weights) and new part (new edges and increases) (\Cref{sec:main})\\
$\Reff(G)_e$ & effective resistance $\cchi_e^{\top}\LL(G)^{\dag}\cchi_e$ (\Cref{sec:main})\\
$U$, $Q$ & total numbers of updates and queries (\Cref{lem:stable})\\
$F$, $\delta_e$ & the edges deleted or decreased in a block and the amounts (\Cref{sec:main})\\
$\Reffref$ & the public resistance estimate of a refresh (\Cref{def:record})\\
$\dheavy$ & heavy threshold, $\dheavy=\log^{3}n$ (\Cref{def:heavy})\\
$\cchi_e$, $\oone$, $\oone_v$ & $\oone_u-\oone_v$ for $e=uv$, the all-ones vector, the $v$-th unit vector (\Cref{sec:prelim})\\
$\dd(G)$, $\DD(G)$, $\AA(G)$ & weighted degrees, their diagonal matrix, the adjacency matrix (\Cref{sec:prelim})\\
$\PPi$ & $\II-\frac1n\oone\oone^{\top}$ (\Cref{sec:prelim})\\
$\MM$ & $\LL(\Gh)+\DD(\Gl)$ (\Cref{def:heavy})\\
$\ZZ(\MM)$ & the fixed solver operator for $\MM$ (\Cref{lem:chain})\\
$\Gh$, $\Gl$ & heavy and light graphs (\Cref{def:heavy})\\
$\tr$, $\norm{\cdot}$, $\norm{\cdot}_F$, $\norm{\cdot}_{S_6}$ & trace, operator, Frobenius and Schatten-$6$ norms (\Cref{sec:resistance})\\
$\XX$ & the matrix $[\xxrand_1\cdots\xxrand_{10^{3}\log n}]$ (\Cref{def:record})\\
$\vol$, $\cut$ & volume of a vertex set and the edges across a cut (\Cref{def:cluster})\\
$\Ccal$, $C$, $(C,v)$ & the partition, a cluster, a copy (\Cref{def:cluster})\\
$\mathrm{center}_{\xx}(C)$ & the center of a cluster for the query vector $\xx$ (\Cref{def:cluster})\\
$\Ecal{2}(\xx,G)$, $\Ecalvert{2}(\xx,C,v)$ & energy $\xx^{\top}\LL(G)\xx$ and its vertex share in a cluster (\Cref{def:cluster})\\
$\Ecaltil$ & the energy estimate given to $\textsc{PartitionTwo}$ (\Cref{lem:partition})\\
$\Ecut$ & edges cut by a query and returned to the caller (\Cref{lem:partition})\\
$\Ex$ & expectation (\Cref{sec:resistance})\\
$\varrho_e$, $p_e$ & intensity of an edge and its retention probability $1-e^{-\varrho_e}$ (\Cref{eq:references})\\
$\zeta_e$ & centered retention variable (\Cref{eq:references})\\
$\LLtil$, $\MMtil$ & sparse and heavy references (\Cref{eq:references})\\
$\BBtil$, $\WWtil$ & incidence and weight matrices of $\LLtil=\BBtil^{\top}\WWtil\BBtil$ (\Cref{sec:resistance})\\
$\QQ_e$, $\QQ^{\mathrm{heavy}}_e$ & the deviation of one edge in the two references (\Cref{eq:references})\\
$\xxrand_j$, $\xxres_j$, $\xxproxy_j$ & random, residual and proxy vectors of an attempt (\Cref{fig:attempt})\\
$\YYbar$, $\YYtil$ & deterministic and random parts of one residual step (\Cref{eq:cubic})\\
$T_v$, $T^{\mathrm{light}}_v$ & per-vertex trees of the current edges and of the active light edges (\Cref{fig:state})\\
$\mathrm{Changed}$ & edges whose heavy/light status changed in the block (\Cref{fig:state})\\
$P$ & the $\textsc{QueryDrivenPartition}$ holding $\Gret_k$ (\Cref{fig:state})\\
$\xxvert_v$ & proxy mass at $v$ (\Cref{def:record})\\
$D(u,v)$ & the packed distance oracle (\Cref{lem:packed})\\
$q_e$ & proposal probability of an edge (\Cref{lem:proposals})\\
$R$, $O$ & the edges sampled for the reference and for the output (\Cref{fig:block})\\
$\SSS$ & $\WW_F^{1/2}\BB_F\LL(G)^{\dag}\BB_F^{\top}\WW_F^{1/2}$ (\Cref{lem:woodbury})\\
$\Z$, $\dist$ & the integers; the squared-distance table of the packed oracle (\Cref{sec:packed})\\
$\bb$, $\vv$ & bit vectors and word vectors of the packed oracle (\Cref{sec:packed})\\
\bottomrule
\end{tabular}
\end{center}

\end{document}
